%%%%%%%%%%%%%%%%%%%%%%%%%%%%%%%%%%%%%%%%%
%DIF LATEXDIFF DIFFERENCE FILE
%DIF DEL diff_work/old/thesis_proposal.tex   Thu Jun 11 13:08:50 2026
%DIF ADD thesis_proposal.tex                 Thu Jun 11 13:23:33 2026
% Uppsala University Thesis Proposal
% LaTeX Document
% Author: Andrei Olaru
%%%%%%%%%%%%%%%%%%%%%%%%%%%%%%%%%%%%%%%%%

\documentclass[12pt]{article}

% Language and encoding
\usepackage[english]{babel}
\usepackage[utf8]{inputenc}
\usepackage{csquotes}

% Mathematics
\usepackage{amsmath}

% Graphics and figures
\usepackage{graphicx}
\usepackage{float}

% Page layout
\usepackage[margin=1in,marginparwidth=2cm]{geometry}

% Bibliography
\usepackage[style=apa,backend=biber]{biblatex}
\addbibresource{references.bib}

% Diagrams
\usepackage{tikz}

% Hyperlinks
\usepackage{hyperref}
\hypersetup{
    colorlinks=true,
    linkcolor=blue,
    filecolor=magenta,
    urlcolor=cyan,
    citecolor=blue,
}

% Section formatting
\usepackage{titlesec}
\titleformat{\section}{\Large\bfseries}{\thesection}{1em}{}
\titleformat{\subsection}{\large\bfseries}{\thesubsection}{1em}{}

% Todo notes
\usepackage[colorinlistoftodos]{todonotes}

%%%%%%%%%%%%%%%%%%%%%%%%%%%%%%%%%%%%%%%%%
% TITLE PAGE
%%%%%%%%%%%%%%%%%%%%%%%%%%%%%%%%%%%%%%%%%

 %DIF > 
%DIF PREAMBLE EXTENSION ADDED BY LATEXDIFF
%DIF UNDERLINE PREAMBLE %DIF PREAMBLE
\RequirePackage[normalem]{ulem} %DIF PREAMBLE
\RequirePackage{color}\definecolor{RED}{rgb}{1,0,0}\definecolor{BLUE}{rgb}{0,0,1} %DIF PREAMBLE
\providecommand{\DIFaddtex}[1]{{\protect\color[rgb]{0,0.5,0}\uwave{#1}}} %DIF PREAMBLE
\providecommand{\DIFdeltex}[1]{{\protect\color{red}\sout{#1}}} %DIF PREAMBLE
%DIF SAFE PREAMBLE %DIF PREAMBLE
\providecommand{\DIFaddbegin}{} %DIF PREAMBLE
\providecommand{\DIFaddend}{} %DIF PREAMBLE
\providecommand{\DIFdelbegin}{} %DIF PREAMBLE
\providecommand{\DIFdelend}{} %DIF PREAMBLE
\providecommand{\DIFmodbegin}{} %DIF PREAMBLE
\providecommand{\DIFmodend}{} %DIF PREAMBLE
%DIF FLOATSAFE PREAMBLE %DIF PREAMBLE
\providecommand{\DIFaddFL}[1]{\DIFadd{#1}} %DIF PREAMBLE
\providecommand{\DIFdelFL}[1]{\DIFdel{#1}} %DIF PREAMBLE
\providecommand{\DIFaddbeginFL}{} %DIF PREAMBLE
\providecommand{\DIFaddendFL}{} %DIF PREAMBLE
\providecommand{\DIFdelbeginFL}{} %DIF PREAMBLE
\providecommand{\DIFdelendFL}{} %DIF PREAMBLE
%DIF HYPERREF PREAMBLE %DIF PREAMBLE
\providecommand{\DIFadd}[1]{\texorpdfstring{\DIFaddtex{#1}}{#1}} %DIF PREAMBLE
\providecommand{\DIFdel}[1]{\texorpdfstring{\DIFdeltex{#1}}{}} %DIF PREAMBLE
\providecommand{\DIFscaledelfig}{0.5}
%DIF HIGHLIGHTGRAPHICS PREAMBLE %DIF PREAMBLE
\RequirePackage{settobox} %DIF PREAMBLE
\RequirePackage{letltxmacro} %DIF PREAMBLE
\newsavebox{\DIFdelgraphicsbox} %DIF PREAMBLE
\newlength{\DIFdelgraphicswidth} %DIF PREAMBLE
\newlength{\DIFdelgraphicsheight} %DIF PREAMBLE
% store original definition of \includegraphics %DIF PREAMBLE
\LetLtxMacro{\DIFOincludegraphics}{\includegraphics} %DIF PREAMBLE
\providecommand{\DIFaddincludegraphics}[2][]{{\color[rgb]{0,0.5,0}\fbox{\DIFOincludegraphics[#1]{#2}}}} %DIF PREAMBLE
\providecommand{\DIFdelincludegraphics}[2][]{% %DIF PREAMBLE
\sbox{\DIFdelgraphicsbox}{\DIFOincludegraphics[#1]{#2}}% %DIF PREAMBLE
\settoboxwidth{\DIFdelgraphicswidth}{\DIFdelgraphicsbox} %DIF PREAMBLE
\settoboxtotalheight{\DIFdelgraphicsheight}{\DIFdelgraphicsbox} %DIF PREAMBLE
\scalebox{\DIFscaledelfig}{% %DIF PREAMBLE
\parbox[b]{\DIFdelgraphicswidth}{\usebox{\DIFdelgraphicsbox}\\[-\baselineskip] \rule{\DIFdelgraphicswidth}{0em}}\llap{\resizebox{\DIFdelgraphicswidth}{\DIFdelgraphicsheight}{% %DIF PREAMBLE
\setlength{\unitlength}{\DIFdelgraphicswidth}% %DIF PREAMBLE
\begin{picture}(1,1)% %DIF PREAMBLE
\thicklines\linethickness{2pt} %DIF PREAMBLE
{\color[rgb]{1,0,0}\put(0,0){\framebox(1,1){}}}% %DIF PREAMBLE
{\color[rgb]{1,0,0}\put(0,0){\line( 1,1){1}}}% %DIF PREAMBLE
{\color[rgb]{1,0,0}\put(0,1){\line(1,-1){1}}}% %DIF PREAMBLE
\end{picture}% %DIF PREAMBLE
}\hspace*{3pt}}} %DIF PREAMBLE
} %DIF PREAMBLE
\LetLtxMacro{\DIFOaddbegin}{\DIFaddbegin} %DIF PREAMBLE
\LetLtxMacro{\DIFOaddend}{\DIFaddend} %DIF PREAMBLE
\LetLtxMacro{\DIFOdelbegin}{\DIFdelbegin} %DIF PREAMBLE
\LetLtxMacro{\DIFOdelend}{\DIFdelend} %DIF PREAMBLE
\DeclareRobustCommand{\DIFaddbegin}{\DIFOaddbegin \let\includegraphics\DIFaddincludegraphics} %DIF PREAMBLE
\DeclareRobustCommand{\DIFaddend}{\DIFOaddend \let\includegraphics\DIFOincludegraphics} %DIF PREAMBLE
\DeclareRobustCommand{\DIFdelbegin}{\DIFOdelbegin \let\includegraphics\DIFdelincludegraphics} %DIF PREAMBLE
\DeclareRobustCommand{\DIFdelend}{\DIFOaddend \let\includegraphics\DIFOincludegraphics} %DIF PREAMBLE
\LetLtxMacro{\DIFOaddbeginFL}{\DIFaddbeginFL} %DIF PREAMBLE
\LetLtxMacro{\DIFOaddendFL}{\DIFaddendFL} %DIF PREAMBLE
\LetLtxMacro{\DIFOdelbeginFL}{\DIFdelbeginFL} %DIF PREAMBLE
\LetLtxMacro{\DIFOdelendFL}{\DIFdelendFL} %DIF PREAMBLE
\DeclareRobustCommand{\DIFaddbeginFL}{\DIFOaddbeginFL \let\includegraphics\DIFaddincludegraphics} %DIF PREAMBLE
\DeclareRobustCommand{\DIFaddendFL}{\DIFOaddendFL \let\includegraphics\DIFOincludegraphics} %DIF PREAMBLE
\DeclareRobustCommand{\DIFdelbeginFL}{\DIFOdelbeginFL \let\includegraphics\DIFdelincludegraphics} %DIF PREAMBLE
\DeclareRobustCommand{\DIFdelendFL}{\DIFOaddendFL \let\includegraphics\DIFOincludegraphics} %DIF PREAMBLE
%DIF AMSMATHULEM PREAMBLE %DIF PREAMBLE
\makeatletter %DIF PREAMBLE
\let\sout@orig\sout %DIF PREAMBLE
\renewcommand{\sout}[1]{\ifmmode\text{\sout@orig{\ensuremath{#1}}}\else\sout@orig{#1}\fi} %DIF PREAMBLE
\makeatother %DIF PREAMBLE
%DIF COLORLISTINGS PREAMBLE %DIF PREAMBLE
\RequirePackage{listings} %DIF PREAMBLE
\RequirePackage{color} %DIF PREAMBLE
\lstdefinelanguage{DIFcode}{ %DIF PREAMBLE
%DIF DIFCODE_UNDERLINE %DIF PREAMBLE
  moredelim=[il][\color{red}\sout]{\%DIF\ <\ }, %DIF PREAMBLE
  moredelim=[il][\color[rgb]{0,0.5,0}\uwave]{\%DIF\ >\ } %DIF PREAMBLE
} %DIF PREAMBLE
\lstdefinestyle{DIFverbatimstyle}{ %DIF PREAMBLE
	language=DIFcode, %DIF PREAMBLE
	basicstyle=\ttfamily, %DIF PREAMBLE
	columns=fullflexible, %DIF PREAMBLE
	keepspaces=true %DIF PREAMBLE
} %DIF PREAMBLE
\lstnewenvironment{DIFverbatim}[1][]{\lstset{style=DIFverbatimstyle}}{} %DIF PREAMBLE
\lstnewenvironment{DIFverbatim*}[1][]{\lstset{style=DIFverbatimstyle,showspaces=true}}{} %DIF PREAMBLE
\lstset{extendedchars=\true,inputencoding=utf8}

%DIF END PREAMBLE EXTENSION ADDED BY LATEXDIFF

\begin{document}

\begin{titlepage}

\newcommand{\HRule}{\rule{\linewidth}{0.5mm}}

\center

\textsc{\LARGE Uppsala University}\\[1.5cm]

\includegraphics[scale=.1]{otherdocuments/uppsala_University_template/Uppsala_University_seal_svg.png}\\[1cm]
\textsc{\Large Master's Thesis}\\[0.5cm]

\HRule \\[0.4cm]
{\huge \bfseries Applying Transformative Game Design Theory to Competitive Debate}\\[0.4cm]
\HRule \\[1.5cm]

\begin{minipage}{0.4\textwidth}
\begin{flushleft} \large
\emph{Author:}\\
Andrei Dan \textsc{Olaru}
\end{flushleft}
\end{minipage}
\DIFaddbegin \DIFadd{~
}\begin{minipage}{0.4\textwidth}
\begin{flushright} \large
\emph{\DIFadd{Supervisor:}}\\
\DIFadd{Josephine }\textsc{\DIFadd{Baird}}\\[0.4cm]
\emph{\DIFadd{Examiner:}}\\
\DIFadd{Sabine }\textsc{\DIFadd{Harrer}}
\end{flushright}
\end{minipage}\DIFaddend \\[2cm]

{\large \today}\\[2cm]

\vfill

\end{titlepage}
\setcounter{page}{2}

%%%%%%%%%%%%%%%%%%%%%%%%%%%%%%%%%%%%%%%%%
%DIF >  ABSTRACT
%DIF > %%%%%%%%%%%%%%%%%%%%%%%%%%%%%%%%%%%%%%%%

\DIFaddbegin 

\begin{abstract}
\DIFadd{This thesis applies transformative game design theory to competitive debate, treating major debate formats not as fixed institutions but as designed games open to deliberate redesign. Although forensics pedagogy has long documented debate's transformative effects, no prior work analyses debate formats from a game-design perspective or brings transformative role-playing-game theory to bear on them. Working at this gap, the thesis asks }\emph{\DIFadd{how}} \DIFadd{such theory can be operationalised into a concrete format, and }\emph{\DIFadd{how}} \DIFadd{participants respond to the format that results.

}

\DIFadd{Using Research through Design, and specifically an evocative design experiment, I develop and playtest }\emph{\DIFadd{Inside Out}}\DIFadd{, a new role-playing debate format built to address three tendencies observed in World Schools and British Parliamentary debate: value convergence, information density, and competitive hierarchy. The format replaces binary opposition and win/loss adjudication with several value-defending teams, a deciding Subject character, and a commendation-based system of evaluation. Two playtest cycles with experienced debaters were analysed through a pre-specified taxonomy of behavioural and self-reported clues.

}

\DIFadd{The contribution includes a documented design and an account of how its mechanics shape the experience of those who play it, including emergent effects the design did not set out to produce. Beyond the artifact, the thesis debates intermediate-level design knowledge about how one might design for transformation in deliberative games by using and the scaffold of debate.
}\end{abstract}

%DIF > %%%%%%%%%%%%%%%%%%%%%%%%%%%%%%%%%%%%%%%%
%DIF >  ACKNOWLEDGEMENTS
%DIF > %%%%%%%%%%%%%%%%%%%%%%%%%%%%%%%%%%%%%%%%

\section*{\DIFadd{Acknowledgements}}

\DIFadd{This thesis and the game on whose design it was built upon would have been different and most likely worse were there not for a few people who graced me with their time and thought. I would like to spare these few lines to thank them and acknowledge clearly their help: Nae Șovăială, for the great ideas he dared to brainstorm with me and for unshackling me from ideas either unactionable or undesireable; and for his great coffee of course; Șerban Pitic for theorycrafting debate world's issues together with me over phone-calls and late-night walks; David Moscovici for challenging me towards bold design; Radu Țigănaș, Adrian Mălăieș and Alexandru Fioșoiu for helping me organize the first playtest; the debaters at the Jean Monnet high school for participating in said playtest; my dear friends Radu Tăunean, Alexandra Badea, George Ionescu, Mihai Tonsciuc and Mihai Oprea for debating during the second playtest and giving me their honest thoughts on the format; and most of all to my partner Maria Trifon for her immense patience and support, and for all the scholarly insight she aided with and without which this would have been nothing more than an opinionated essay.

}

%DIF > %%%%%%%%%%%%%%%%%%%%%%%%%%%%%%%%%%%%%%%%
\DIFaddend % CONTENT
%%%%%%%%%%%%%%%%%%%%%%%%%%%%%%%%%%%%%%%%%

\section{INTRODUCTION}

\subsection{Brief Overview}

Through this thesis, I \DIFdelbegin \DIFdel{apply transformative game design theory and best practices to }\DIFdelend create a new competitive debating format \DIFaddbegin \DIFadd{with the use of the theory and best practices of transformative game design: the intentional design of games, and of the rituals, roles, and community structures around them, to facilitate lasting personal and social change in those who play \mbox{%DIFAUXCMD
\parencite{bowman2025transformative}}\hskip0pt%DIFAUXCMD
}\DIFaddend . This research employs Research through Design (RtD) methodology to develop and playtest an alternative debate format that addresses what I believe to be systematic problems in widely-used formats, particularly World Schools and British Parliamentary debate.

\subsection{Motivation}

As a debater, judge, coach, and organizer with 15 years of experience, my insider perspective motivates this reformist work. \DIFdelbegin \DIFdel{Debate }\DIFdelend \DIFaddbegin \DIFadd{In this thesis I argue that debate }\DIFaddend is a class of role-playing games with \DIFdelbegin \DIFdel{immense }\DIFdelend \DIFaddbegin \DIFadd{significant }\DIFaddend transformative potential, yet the design of major formats retains decades-old flaws, \DIFdelbegin \DIFdel{with game design methodology }\DIFdelend \DIFaddbegin \DIFadd{and game-design methodology is almost }\DIFaddend entirely absent from attempts to improve them. \DIFaddbegin \DIFadd{Forensics pedagogy has long documented what debate can do for those who practise it; what no one has yet done is treat its formats as games open to deliberate design. Bringing that perspective to bear is the work this thesis undertakes.
}\DIFaddend \subsection{Problem Statement}

\DIFdelbegin \DIFdel{Although a very diverse range of debate formats are played }\DIFdelend \DIFaddbegin \DIFadd{The problem this thesis addresses is, first of all, a gap in how debate is studied and improved. Debate formats are designed artifacts, played competitively }\DIFaddend all over the world, \DIFaddbegin \DIFadd{yet they have escaped }\DIFaddend the iterative game-design improvement \DIFdelbegin \DIFdel{seen in other games is absent}\DIFdelend \DIFaddbegin \DIFadd{that comparable games receive}\DIFaddend : most formats are \DIFdelbegin \DIFdel{deeply }\DIFdelend calcified, their competitive communities evolving new ``metas'', modes of play, and interpretations of rules without game-design supervision. \DIFdelbegin \DIFdel{As an example, I hypothesize that the }\DIFdelend \DIFaddbegin \DIFadd{No body of work analyses major debate formats as games to be designed, and the transformative role-playing-game theory that could inform such design has never been brought to bear on them. Addressing this gap is the central contribution of the thesis.

}

\DIFadd{The gap is worked at a concrete site: three problematic tendencies I observe in the }\DIFaddend two most popular \DIFdelbegin \DIFdel{debate }\DIFdelend formats, World Schools and British Parliamentary, whose \DIFdelbegin \DIFdel{communities in Europe are deeply intertwined, exhibit the following problematic tendencies:  }\DIFdelend \DIFaddbegin \DIFadd{European communities are closely intertwined. }\DIFaddend Firstly, they seem to push debaters towards value convergence, an exclusionary and self-right\DIFaddbegin \DIFadd{e}\DIFaddend ous holding of the same beliefs, potentially causing them to become less tolerant of people with different values and diminishing their ability to convince people from different backgrounds. Secondly, they appear to teach debaters to speak \DIFdelbegin \DIFdel{with extreme speed and }\DIFdelend \DIFaddbegin \DIFadd{rapidly and with high }\DIFaddend information density, which could diminish their ability to meaningfully support their positions to people not trained as debate judges. Thirdly, \DIFdelbegin \DIFdel{in my opinion they }\DIFdelend \DIFaddbegin \DIFadd{they appear to }\DIFaddend sustain competitive pressure that erodes psychological safety and ossifies into status hierarchies among debaters.

These \DIFdelbegin \DIFdel{observations }\DIFdelend \DIFaddbegin \DIFadd{three tendencies }\DIFaddend are not empirically validated, and this thesis does not intend to prove them, though a later section presents their echoes in literature. \DIFaddbegin \DIFadd{They function here as the design problem the new format responds to, the concrete material on which the absent game-design perspective is put to work.

}\DIFaddend 

\subsection{Research Question}

How can the application of transformative game design theory inform the creation of a competitive debate format that manages to elevate its transformative potential (\DIFaddbegin \DIFadd{such as }\DIFaddend improved critical thinking, improved discourse, \DIFdelbegin \DIFdel{self empowerment etc}\DIFdelend \DIFaddbegin \DIFadd{and self-empowerment}\DIFaddend ) by addressing the issues of value convergence, information density and competitive hierarchy?

This question decomposes \DIFdelbegin \DIFdel{in two: can a method of applying said theory to this problembe described and demonstrated? and does that method yield expected results? }\DIFdelend \DIFaddbegin \DIFadd{into two. First, }\emph{\DIFadd{how}} \DIFadd{can transformative game-design theory be operationalised into a concrete debate format: what does the design process look like when the theory is brought to bear on the problem? Second, }\emph{\DIFadd{how}} \DIFadd{do participants respond to an example resulting format: what experiences does it evoke in those who play it? The first is a question about design; the second, about the play that design produces.

}\DIFaddend 

\subsection{Scope and Limitations}

This thesis aims to demonstrate proof of concept: that alternative debate formats can be systematically designed to address \DIFaddbegin \DIFadd{identified }\DIFaddend problems, and that such designs produce observable, articulable differences in participant experience. \DIFdelbegin \DIFdel{The demonstration draws on participant observation and participant-reported experiences across }\DIFdelend \DIFaddbegin \DIFadd{These are assessed along }\DIFaddend five evaluation fields\DIFaddbegin \DIFadd{: (1) the extent that it sustains }\emph{\DIFadd{long-term engagement}} \DIFadd{(players return to it voluntarily), }\DIFaddend (\DIFdelbegin \DIFdel{engagement and appeal, productive pressure, growth direction, community dynamics, and epistemic humility) , assessing whether a format grounded in transformative game design theory can produce clues of meaningful change without sacrificing the intellectual challenge or communal value of the activity}\DIFdelend \DIFaddbegin \DIFadd{2) that it exerts }\emph{\DIFadd{productive pressure}} \DIFadd{(play stretches participants rather than merely pleasing them), (3) that it points growth in a useful }\emph{\DIFadd{direction}} \DIFadd{(the skills it builds transfer beyond the activity), (4) that it maintains }\emph{\DIFadd{container safety}} \DIFadd{(the social environment supports rather than threatens its players), and (5) that it cultivates }\emph{\DIFadd{epistemic humility}} \DIFadd{(players come to hold their own values more pluralistically). These five fields double as the categories under which playtest observations were gathered and interpreted (Section~\ref{sec:methods})}\DIFaddend .

The thesis does not aim to produce robust claims about long-term educational efficacy or entertainment value\DIFdelbegin \DIFdel{, as }\DIFdelend \DIFaddbegin \DIFadd{; }\DIFaddend such claims would require longitudinal studies with larger samples than a master's thesis permits\DIFdelbegin \DIFdel{. Its main }\DIFdelend \DIFaddbegin \DIFadd{, and, as Section~\ref{sec:methods} argues, they belong to a different mode of inquiry than the one practised here. Its }\DIFaddend contribution to knowledge is \DIFdelbegin \DIFdel{to exemplify that debateformats can be designed better if looked at as roleplaying games}\DIFdelend \DIFaddbegin \DIFadd{of the kind Research through Design characteristically produces: transferable, intermediate-level design knowledge, an account of }\emph{\DIFadd{how}} \DIFadd{one might design for transformation in an interpersonal, argumentative setting by using and subverting the structure of debate. The designed format is itself part of that contribution \mbox{%DIFAUXCMD
\parencite{zimmerman2007research}}\hskip0pt%DIFAUXCMD
}\DIFaddend .

\section{BACKGROUND AND LITERATURE REVIEW}
\label{sec:background}

\subsection{Debate as Transformative Game}

\DIFdelbegin \DIFdel{Academic debate}\DIFdelend \DIFaddbegin \DIFadd{Many activities call themselves ``debate.'' \mbox{%DIFAUXCMD
\textcite{snider2008code} }\hskip0pt%DIFAUXCMD
defines it broadly, as ``the process that determines how change should occur,'' reducing to a ``discussion }[\DIFadd{that}] \DIFadd{resolves whether a specific change should occur.'' That general definition is somewhat limitative, and Snider himself refines it, narrowing his focus to ``a formal competitive type of debate'' governed by explicit rules. It is debate in this latter sense that this thesis concerns: the leisure and educational activities performed on specific rulesets, or }\emph{\DIFadd{formats}}\DIFadd{, of which World Schools and British Parliamentary are the most widely addressed in this paper due to their wide adoption and because of my familiarity with them.

}

\DIFadd{In this thesis I argue that academic debate }\DIFaddend formats are games. Following \textcite{salen2004rules}, a game is a system in which players engage in artificial conflict, defined by rules, resulting in a quantifiable outcome. Academic debate satisfies all criteria: rules, participants, winners, losers, tournaments, and scoring systems. Alfred Snider, \DIFaddbegin \DIFadd{in }\DIFaddend whose school of debate \DIFdelbegin \DIFdel{of }\DIFdelend training I was \DIFdelbegin \DIFdel{priviledge in being apprenticedin}\DIFdelend \DIFaddbegin \DIFadd{privileged to be apprenticed}\DIFaddend , through his 1982 doctoral dissertation at the University of Kansas, along with his subsequent publication \parencite{snider1984gaming}, established the scholarly foundation for viewing debate as a game while documenting its educational benefits.

But what makes a game \textit{transformative}? \textcite{bowman2025transformative} conceptualize transformative role-playing games as those intentionally designed to encourage personal and social change arising from play. For transformation to occur, this change must expand beyond the bounds of the original experience and integrate into the participant's daily frames of reality and identity. By this standard, a game is truly transformative not merely when participants report enjoyment or in-game learning, but when the experience produces lasting changes in cognition, behavior, or self-concept.

Forensics pedagogy documents debate's transformative potential across several domains that satisfy this criterion. The most extensively studied is critical thinking: debaters develop the ability to recognize stated and unstated assumptions, evaluate causal inferences, and assess the validity of arguments \parencite{colbert1995enhancing, hill1993value}. Beyond cognitive gains, debate serves as a vehicle for empowerment, particularly for previously ``silenced student populations'' who develop their own voice through the activity \parencite{snider2003gamemaster}. The activity cultivates what \textcite{hicks1998public} calls a ``democratic ethos'': the capacity to justify convictions in public, respect minority views, and evaluate opposing arguments critically. Debate also provides a socialization mechanism, integrating students into a sub-culture with specialized language that enhances their sense of belonging to an intellectual community \parencite{cirlin1997sociological, kargacin1998state}. Crucially, these effects transfer beyond the activity into academic performance, professional communication, and civic engagement, satisfying the criterion of transformation beyond the game.

% Empirical research substantiates these claims, with quantitative evidence accumulating across critical thinking, academic achievement, and college readiness over several decades. \textcite{allen1999meta} conducted a meta-analysis of studies examining forensics and communication education, finding a significant positive relationship between such training and gains in critical thinking (longitudinal designs $r = .18$, cross-sectional designs $r = .20$), with forensics participation demonstrating the largest effect among the communication activities studied. At the individual study level, \textcite{colbert1987effects} conducted the most comprehensive investigation, tracking 275 CEDA and NDT debaters and non-debaters across nine universities over a full academic year. Using analysis of covariance controlling for pretest scores, he found that debaters' Watson-Glaser Critical Thinking Appraisal scores rose from 61.18 to 64.53 while non-debaters' scores actually declined from 52.67 to 49.14, with NDT debaters significantly outgaining non-debaters ($F = 31.77$, $p < .001$). While \textcite{hill1993value} cautioned that the broader body of debate-critical thinking research remained inconclusive due to methodological limitations and inconsistent results across sub-samples, \textcite{colbert1995enhancing} responded that the preponderance of evidence establishes ``presumptive proof'' favoring a positive debate-critical thinking relationship, particularly when properly designed studies using robust statistical controls are considered.
%
% Beyond critical thinking, debate participation correlates with substantial academic outcomes. \textcite{mezuk2009urban} analyzed data from the Chicago Debate League spanning 1997 to 2006 and found that participants were 70\% more likely to graduate and three times less likely to drop out than comparable non-participants, even after controlling for eighth-grade test scores and grade point average. Debaters were also more likely to meet ACT benchmarks for college readiness in English and reading, though not in science or mathematics. A subsequent expanded study of the same program by \textcite{mezuk2011impact}, analyzing 9,145 students of whom 2,449 participated in at least one CDL tournament, found that debaters were 2.5 times less likely to drop out, graduated at a rate of 90\% compared to 75\% for non-debaters, and achieved a mean twelfth-grade GPA of 3.23 compared to 2.83 for their peers. Most recently, \textcite{schueler2025debate} examined policy debate in Boston Public Schools serving large concentrations of economically disadvantaged students and found positive impacts on English Language Arts scores of 0.13 standard deviations, equivalent to 68\% of a full year of average ninth-grade learning. Critically, these gains were concentrated on analytical rather than rote subskills and were largest among the lowest-performing students at baseline, with additional evidence of positive effects on high school graduation and postsecondary enrollment.

Empirical research substantiates these claims. \textcite{allen1999meta} conducted a meta-analysis finding a significant positive relationship between debate participation and critical thinking gains (longitudinal $r = .18$), with debate demonstrating the largest effect among communication activities studied. At the individual study level, \textcite{colbert1987effects} found significant critical thinking gains in debaters relative to non-debaters over a full academic year.

\DIFdelbegin \DIFdel{.
}\DIFdelend 

Beyond critical thinking, debate participation correlates with substantial academic outcomes. \textcite{mezuk2009urban} found Chicago Debate League participants were 70\% more likely to graduate and three times less likely to drop out than comparable non-participants.

Forensics scholars have proposed theoretical frameworks for these effects, each grounded in established learning theory. \textcite{colbert1995enhancing} argues that critical thinking develops through what he calls intellectual ``disequilibrium,'' a Piagetian concept where confrontation with opposing positions forces learners to restructure their mental models. Debate's adversarial structure ensures this confrontation is unavoidable: a debater cannot ignore the opposing case but must engage with it directly. \textcite{cirlin1997sociological} grounds debate's communicative and socialising effects in Richard Weaver's rhetorical theory and the concept of ``sermonic'' language, while \textcite{hicks1998public}'s ``democratic ethos'' draws on Rawls's ideal of public reason. Drawing these threads together, \textcite{bellon2000research} synthesized evidence from competitive forensics, communication pedagogy, and educational psychology to argue that debate is uniquely effective precisely because it engages cognitive, communicative, social, and civic levels simultaneously. However, these forensics-specific frameworks have not been connected to broader theories of experiential and transformative learning.

% In the Theoretical Framework section of this paper, I apply \textcite{vygotsky1978mind}'s zone of proximal development to the scaffolded learning that occurs through debate coaching and graduated competition, map \textcite{kolb1984experiential}'s experiential learning cycle onto the iterative process of debating, receiving feedback, and refining strategy, and draw on research into structured debriefing in simulation and role-playing \parencite{crookall2014engaging, daniau2016transformative} to illuminate how debate's built-in oral critique enables the reflective processing essential to lasting transformation.

\DIFdelbegin \DIFdel{Although forensics pedagogy has documented debate's transformative effects, no prior work analyzes major debate formats from a game design perspective or applies transformative role-playing game theory to competitive debate. Addressing this gap is the central contribution of this thesis.
}%DIFDELCMD < 

%DIFDELCMD < %%%
\DIFdelend \subsection{Debate as Role-Playing Game}

In the World Schools format, two teams receive a motion (e.g., ``This House Would legalize the possession of drugs'') and are randomly assigned to Proposition/Government or Opposition. They deliver alternating speeches as if individuals supporting or opposing the motion \DIFaddbegin \DIFadd{\mbox{%DIFAUXCMD
\parencite{quinn2009debating}}\hskip0pt%DIFAUXCMD
}\DIFaddend . Judges listen as someone with no opinion on the motion, expert in debate rules, and averagely informed, deciding who most convinced the character they role-play. The debate ends with the judge ruling who won, explaining the decision, and offering feedback for improvement.

Understanding debate's relationship to role-playing games requires first establishing what role-playing games are and where debate fits among related practices. \textcite{zagal2024definitions} note that defining role-playing games is contentious, but certain characteristics recur: rule-structured creation and enactment of characters in a fictional world, with broad choice over what actions to attempt. \textcite{bowman2025transformative} offer a functional definition: analog role-playing games are ``co-creative experiences in which participants immerse into fictional characters and realities for a bounded period of time and improvise through spontaneous, emergent playfulness'' (p. 16). Debate fits both definitions: it places fictional defenders of positions in a fictional world where a crucial decision will be taken based on their discourse, co-created and time-bound, within which everything is playfully improvised.

Transformative game design literature explicitly positions debate among ``cousin forms'' of analog role-playing games \parencite{bowman2014educational,bowman2025transformative}, alongside simulations, educational role-playing, improv theatre, and psychodrama. The distinction rests on community practices, norms, and innovations, not on the act of role-playing itself. While debate is indeed used in educational contexts, it differs fundamentally from close cousins such as work simulations or classroom exercises in ways that bring it closer to analog role-playing games proper. Three theoretical dimensions illuminate this proximity: the depth of character immersion, the strength of alibi, and the leisure-community dimension.

First, regarding immersion: As opposed to work simulations or \DIFdelbegin \DIFdel{boardgames}\DIFdelend \DIFaddbegin \DIFadd{board games}\DIFaddend , debate requires players to genuinely inhabit an epistemic stance, constructing and defending arguments as if they believed them. This process of active perspective-taking requires what \textcite{bowman2010functions} describes as adopting a theory of mind, thinking ``as if'' one were someone else in a unique set of circumstances (p.~57). The debater must anticipate objections, identify weaknesses, and respond to challenges as someone who holds the assigned position, a cognitive demand that approaches the immersion into character central to role-playing games. As \textcite{seo2022good} observes, ``the moment one received a motion, one adopted the mindset of a person who was completely convinced of that point of view. One clung to this feeling of absolute conviction to make arguments, sink objections, and display passion'' (p.~222). I and other debaters have observed that the cognitive load of debate performance is incompatible with anything short of fully assuming the assigned opinion for the round's duration. Seo corroborates this from the inside: ``Midspeech, a speaker rarely experiences the physical strain of thinking so damn hard. The adrenaline induced a mental fogginess that numbed the hard edges of experience'' (p.~159). This passionate delivery, especially under such cognitive load, is difficult to sustain through meta-cognitive simulation alone. Affective role-assumption becomes the efficient path, organically eliciting rather than faking the emotions expressed.

\textcite{bowman2010functions} identifies nine types of relationships between players and characters, including the ``Fragmented Self,'' where a facet of the player's personality or interests becomes magnified into a central character trait. Debate characters fit this framework: they represent fragmented selves whose defining feature is epistemic commitment to a particular position. The debater role-plays as someone who believes the motion should pass or fail, adopting that argumentative perspective as the character's core identity. They find in themselves beliefs they can hold that support the position and banish from thought all other, focusing on a fragment of their identity.

Second, regarding alibi: this concept describes how game participants can act with lessened social consequences, encouraging risk-taking and willingness to fail \parencite{deterding2017alibis}. Debate provides robust alibi through its competitive frame. Because debate is ``just a game,'' debaters can advance positions they might not advocate publicly and explore uncomfortable arguments. \textcite{seo2022good} captures this vividly: the random assignment of sides provided a paradoxical ``sense of freedom,'' allowing him to ``flirt with ideas, unencumbered by expectations of consistency or deep conviction'' (p.~20). Work simulations do not offer that, as the situational enactment they demand does not require holding specific beliefs, which are a strong coordinate of identity. This enables the identity experimentation that \textcite{bowman2025transformative} identify as central to transformational containers. Debate, as a competitive game, does not offer alibi to fail intellectually; but many TTRPGs and Nordic larps similarly demand performative or emotional skill while still offering identity alibi. In mainstream North American boffer LARPs, for instance, players who ``break game'' risk derision or ostracism \parencite{brown2016culture}, while in Nordic larp communities, immersion into character is often treated as an indicator of success, with its absence experienced as personal failure \parencite{harder2018larp,pohjola2015steering}. Alibi to fail is not essential to a role-playing game; alibi to assume the mantle of a character is. On that metric, debate offers players a \DIFdelbegin \DIFdel{very }\DIFdelend strong alibi to embody a character with \DIFdelbegin \DIFdel{deeply }\DIFdelend \DIFaddbegin \DIFadd{markedly }\DIFaddend different opinions from their own: what Seo describes as tracing ``the steps of how a sensible person (us!) could arrive at conclusions that might otherwise have seemed alien'' (p.~224).

Third, regarding the leisure-community dimension: competitive debate possesses passionate discourse communities, distinguishing role-playing games from cousin forms. I am part of these communities\DIFdelbegin \DIFdel{. The }\DIFdelend \DIFaddbegin \DIFadd{, and have served on the board of ARDOR, the Romanian national debate association. I have seen hhe }\DIFaddend World Schools and British Parliamentary circuits encompass hundreds of competitions \DIFaddbegin \DIFadd{each year, anchored by flagship championships that recur on a fixed annual calendar and draw participants from dozens of countries: the World Universities Debating Championship (WUDC), the European University Debating Championship (EUDC), and, at school level, the World Schools Debating Championship (WSDC)}\DIFaddend . Debaters participate voluntarily, often for decades; they develop specialized terminology, shared references, and community norms; they attend tournaments analogous to gaming conventions. The activity is often pursued for its intrinsic satisfactions, not merely as preparation for \DIFaddbegin \DIFadd{an }\DIFaddend external goal. \DIFaddbegin \DIFadd{From this insider vantage I can attest both to the depth of voluntary commitment that sustains these communities and to the status dynamics that accompany it, taken up among the format critiques below. }\DIFaddend This leisure dimension, combined with the vibrant discourse community, positions competitive debate closer to role-playing games than to the educational simulations with which it shares surface similarities.

Furthermore, evidence for debate's role-playing dimension appears throughout its practice. Experienced coaches regularly employ explicit role-playing techniques during preparation, asking debaters to ``become'' stakeholders affected by policies. Furthermore, debate ethos as I witnessed it is summoned through traditions that explicitly stage parliamentary fiction: bench-pounding, heckling, Points of Information delivered as if holding onto wigs. The format performs a fiction of parliamentary deliberation, with debaters and judges alike participating in theatrical conventions inherited from eighteenth-century British debate culture.

Genuine differences between debate and games that are more traditionally seen as analog\DIFdelbegin \DIFdel{ue}\DIFdelend  role-playing games, such as Dungeons and Dragons, are undeniable: debate usually lacks the narrative worldbuilding, the character backstories, and the extended campaign structures of tabletop or live-action role-playing. Debate's role-playing operates through the same psychological mechanisms, justifying application of the substantial body of transformative game design scholarship to how design choices might enhance or undermine its educational effects.

\subsection{Literature Review of Format Critiques}

Game design theory addresses problems that forensics pedagogy has not solved. The largest body of work that analyzes debate as an educational tool is that of forensics pedagogy. This centers around the more predominant formats played in the US (especially American Policy). I will thus focus on reviewing what literature was written that echoes the problems I set out to solve and that forms arguments that are also applicable to the formats that I seek to improve (especially the World Schools format).

\subsubsection{The Technical-Public Sphere Tension}

A fundamental critique pervading the literature concerns the evolution of debate formats away from public sphere accessibility toward technical sphere insularity. \textcite{combs2002applying} documents how NDT and CEDA have evolved into technical sphere activities characterized by expert-level topic knowledge, speaking rates exceeding what listeners can process, sophisticated analytical models, and specially trained judges: Colbert found average CEDA debaters speaking at approximately 237 words per minute and NDT finalists at 284, rates that ``greatly exceed what is considered optimal for normal speaking'' \parencite[p.~36]{combs2002applying}.

While World Schools and British Parliamentary formats explicitly discourage fast delivery through style-based evaluation, similar pressures toward information density persist. These formats evaluate on two co-equal axes, Style and Content \DIFaddbegin \DIFadd{\mbox{%DIFAUXCMD
\parencite{harveysmith2011practical}}\hskip0pt%DIFAUXCMD
}\DIFaddend , yet increasing information density improves argumentative quality while reducing intelligibility. Because argumentative quality is perceived as more objective than style in an academic setting, the competitive meta\DIFaddbegin \DIFadd{, in my experience of these circuits, }\DIFaddend gradually shifts toward privileging Content over Style despite the rubric formally treating them as equal. The result is that skills developed may not transfer to contexts without expert judges, a difficulty captured in the audience-standard literature, where \textcite{hicks1998public} and \textcite{weiss2002audience} disagree on whether debaters should be persuasive to a regular audience or to the perfect rational listener, with educational value lost either way.

\subsubsection{Value Convergence and Ideological Homogenization}

The ``closed system'' dynamic described in \textcite{combs2002applying}, creates conditions conducive to value convergence. In academic debate, tournament practices are insulated from external feedback and judges who deviate from community norms are ``effectively excluded from the system''; those who object are rejected ``much like the human system rejects certain bacteria'' \parencite[p.~40]{combs2002applying}. This dynamic may cultivate shared assumptions that, while facilitating intra-community communication, potentially diminish debaters' capacity to engage persuasively with those holding different values. In my experience this propagates throughout the community through argument borrowing. Debaters observe which arguments succeed in out-rounds, then adopt those lines of reasoning for their strategic value rather than out of genuine engagement with the underlying ideas, producing superficial understanding and reduced capacity for authentic cross-ideological dialogue.


\subsubsection{Hierarchical Dynamics and Psychological Safety}

Forensics pedagogy is thinnest on this issue, yet several authors recognize the risk. \textcite{snider1999code} addresses safety and ethics in his Code of the Debate, and in his later practical guide warns organizers to ``avoid a counterproductive rigid and hierarchical organization,'' noting ``I've seen groups die when they let ego and status get in the way'' \parencite[p.~36]{snider2014sparking}. The International Debate Education Association's ethics guide for coaches and youth workers similarly acknowledges the ``unbalanced power dynamic'' inherent in coach-student relationships, compounded by the fact that ``older debaters often mentor younger teams'' and thus occupy dual roles of peer and authority figure \parencite{idea2024ethics}.

The mechanism behind these dynamics is illuminated by recent research. \textcite{prussikora2025cracovian} note that ``people placed in competitive intellectual situations tend to develop an adversarial mindset'' characterized by myside bias and otherside bias, distorting cognitive appraisal of the issues being debated. They identify competitive debate's ``most crucial weakness'': its rules ``do not allow for saying to the opponent: `You are right, I was wrong'\,'' instead training debaters to listen for omissions and inaccuracies rather than for understanding \parencite[p.~1474]{prussikora2025cracovian}. When win-loss records become the primary measure of status within these hermetic communities, the adversarial structure produces natural hierarchies in which winning debaters gain social authority while losing debaters are marginalized, reinforcing competitive pressure at the expense of collaborative learning.

\subsection{Gap Identification}

This literature review reveals a significant gap\DIFdelbegin \DIFdel{. Forensics }\DIFdelend \DIFaddbegin \DIFadd{: forensics }\DIFaddend pedagogy has established debate's educational benefits but has not applied \DIFdelbegin \DIFdel{game design frameworks to analyze }\DIFdelend \DIFaddbegin \DIFadd{game-design frameworks to analyse }\DIFaddend how format mechanics produce \DIFdelbegin \DIFdel{specific effects. Addressing this gap is a central contribution of this thesis}\DIFdelend \DIFaddbegin \DIFadd{those effects. It is precisely this analysis, of bringing transformative game-design theory to bear on the mechanics of competitive debate, that the next chapter undertakes}\DIFaddend .

\section{THEORETICAL FRAMEWORK}
\label{sec:theory}

\subsection{Core Theoretical Constructs}

\subsubsection{The Transformational Container}

The concept of the transformational container was adapted to role-playing game theory by \textcite{bowman2021magic}. It describes a bounded space, building upon \textcite{huizinga1958homo}'s ``magic circle'' but additionally requiring consent, safety structures, and community support oriented toward lasting personal change \parencite{bowman2021magic, baird2022transformative, bowman2025transformative}. When these conditions hold, the container provides participants enough psychological safety to risk vulnerability; the resulting discomfort becomes the material for growth. I argue that debate, in its existing forms, already constitutes a partially-realized transformational container. The format's explicit rules and conventions (speaking times, turn order, permissible modes of engagement) establish boundaries that demarcate the debate encounter from ordinary social interaction. Within these boundaries, the competitive frame provides what \textcite{deterding2017alibis} terms \textit{alibi}: because side allocation is typically random and the activity is understood as a game, debaters can explore unfamiliar arguments and adopt perspectives they might not otherwise endorse with reduced social risk. The judge functions as an authority analogous to the facilitator in transformative game design \parencite{bowman2025transformative}: a figure who maintains the rules, calibrates the encounter through feedback, and structures the post-round evaluation.

% Crucially, the container is a \textit{group} space, and group dynamics shape its transformative capacity. \textcite{yalom2020theory} identifies therapeutic factors operating in group settings---universality (discovering others share one's struggles), group cohesion, altruism, and interpersonal learning---all of which are present in debate communities. The shared experience of intellectual vulnerability, of being publicly challenged and occasionally defeated, creates bonds that support continued risk-taking. When these group dynamics function well, the container is strengthened; when status hierarchies or competitive antagonism dominate, the group forces that sustain the container erode.

% \subsubsection{Identity and the Debater-Character Relationship}

% If the container provides the bounded space within which transformation occurs, identity theory specifies \textit{what} is being transformed. \textcite{bowman2010functions} proposes a model of identity in role-playing as fundamentally fragmented: rather than possessing a single, unified self, players bring multiple identity facets---values, beliefs, competencies, emotional dispositions---any of which may be activated, magnified, or recombined through character play. \textcite{bowman2025transformative} develop this further through the concept of \textit{mosaic identity}, where the self is understood as a composite of overlapping fragments that are differentially activated across contexts. Role-playing does not create identity from nothing; it selects, amplifies, and reorganizes existing facets of the player's self.

% The debater-character relationship exemplifies this fragmented self model. As established in the background chapter, debate characters represent fragmented selves whose defining feature is epistemic commitment to a particular position. The debater locates within themselves beliefs, knowledge, and rhetorical dispositions that can support the assigned side, magnifying these into the character's core identity while suppressing contradictory facets. This is not pretense but genuine identity work: the cognitive demands of competitive performance require authentic inhabitation of the epistemic stance, producing what \textcite{bowman2010functions} describes as developing a ``theory of mind'' outside one's own. \textcite{lankoski2012embodied} explain this through grounded cognition theory: performing an epistemic stance neurologically activates the same processes as holding that stance, meaning that the debater-character's identity work is cognitively genuine rather than merely theatrical.

% This identity framework is essential for understanding the bleed mechanisms that follow. Memetic bleed operates precisely through identity facets: arguments rehearsed and positions inhabited do not simply vanish when the round ends, because they were constructed from genuine fragments of the debater's self. The facets activated during play---the capacity to see merit in an opposing view, the rhetorical skill of making a case one does not initially endorse---persist and integrate into the debater's broader identity mosaic. Emancipatory bleed, similarly, works through identity: for marginalized students, the experience of being taken seriously as intellectuals activates and strengthens identity facets---confidence, intellectual authority, public voice---that may have been suppressed in other contexts.

\subsubsection{Bleed as Transformative Mechanism}

Following \textcite{hugaas2024bleed}, bleed describes the bidirectional spillover of psychological contents between player and character (in debate, between the debater and their assigned position). In debate, I identify 3 types of bleed at play.

Memetic bleed, the transfer of ideas and perspectives, explains how debaters develop what I call epistemic humility: skepticism towards one's own opinions and openness to the beliefs of others. The opinions of the role played in a debate bleed directly toward the debater, becoming arguments that can be summoned again outside the round. More importantly, after years of arguing positions contrary to their beliefs, debaters should bleed into accepting that reasonable people hold opposing views and that argumentative power does not equate with being right. This mechanism provides a theoretical account of what forensics scholars have documented empirically: Colbert's ``disequilibrium'' \parencite{colbert1995enhancing} describes the trigger that opens the channel for memetic bleed, while the ``democratic ethos'' that \textcite{hicks1998public} identifies as debate's civic outcome is the intended result of that bleed integrating into the debater's empowered attitude. This should render it impossible for value convergence, the process of members of a community becoming intolerant to other opinions, to emerge. Later on I hypothesize why I observed it doing so.

This bleed of a humble attitude connects to \textcite{bogost2007persuasive}'s concept of procedural rhetoric. Bogost argues that every game, through its systems and rules, reflects a world-view that players may in time be persuaded of. The rules of debate inherently postulate that anyone can support anything and that the dialectic process reaches not truth but a showcase of who is more capable in argument (the same motion can be decided either way in different rounds). Thus experienced debaters ideally learn to acknowledge how there may be much truth in the points made also by the less skilled in argument.

Procedural bleed, the transfer of physical abilities and habits, explains skill development. The rhetorical habits, speaking patterns, and argumentative reflexes developed through competition bleed into debaters' ordinary communication, producing the documented improvements in public speaking and critical thinking. \textcite{cirlin1997sociological}'s account of how debaters develop persuasive skill describes precisely this form of procedural bleed and its transfer beyond the activity.

Emancipatory bleed \parencite{kemper2017battle,kemper2020wyrding} explains debate's particular power for students from marginalized backgrounds. For students where eloquent public speaking was not modeled, debate provides a transformational container \parencite{bowman2021magic} to discover and develop their voice. Being taken seriously as intellectuals, having ideas weighed on merits rather than dismissed based on identity, produces lasting changes in self-concept and capability. \textcite{snider2003gamemaster}'s documentation of debate empowering previously ``silenced student populations'' through Urban Debate Leagues provides empirical evidence for this form of bleed in competitive debate contexts.

%\subsubsection{Debrief and Integration Practices}

%\textcite{bowman2025transformative} present a taxonomy of debrief practices essential to transformative design, distinguishing several dimensions: structured versus unstructured debriefs, emotional versus intellectual versus educational processing, and practices of de-roling (explicitly stepping out of character), narrativization (constructing meaning from the experience), and post-game integration (connecting in-game experiences to out-of-game life). The debrief is not merely a postscript to play but the primary site where raw experience is converted into lasting transformation \parencite{crookall2014engaging,daniau2016transformative}.

%Traditional debate's built-in debrief---the oral critique---is revealing when analyzed through this taxonomy. It is structured but narrowly so: the judge explains their decision, identifies strengths and weaknesses, and offers advice for improvement. This processing is almost exclusively intellectual and evaluative. The judge assesses what worked and what did not as argumentation, with no attention to emotional processing, no de-roling (debaters are expected to simply move on to the next round), and minimal integration connecting the argumentative experience to personal development. The absence of emotional processing is particularly significant given that debate regularly requires inhabiting positions contrary to one's beliefs, engaging with emotionally charged topics, and experiencing public evaluation. \textcite{montola2010positive} argues that negative emotions in role-playing---discomfort, vulnerability, frustration---can be transformative rather than merely distressing, but only when the container provides adequate processing space. Traditional debate formats offer no such space: the oral critique processes the \textit{game performance} but not the \textit{human experience}.

%The RPG-specific debrief literature reinforces this concern: \textcite{atwater2016debrief}'s foundational review characterizes debriefs as ``non-therapeutic transitional tools for moving between perspectives'' and cautions that debriefing alone does not make activities safer---safety must be designed into the format itself. This analysis reveals debrief design as a major untapped lever for enhancing debate's transformative potential. Formats might incorporate emotional check-ins, explicit de-roling moments separating the debater from the position argued, or structured reflection connecting argumentative experiences to the debater's developing worldview.

%\subsubsection{Safety Frameworks}

%\textcite{bowman2025transformative} present a comprehensive safety model organized across three temporal phases and a community-level dimension. \textit{Before-game} safety includes briefing participants on content and expectations, establishing consent frameworks, and calibrating the intensity of play. \textit{During-game} safety encompasses real-time safety mechanics enabling participants to signal discomfort, escalate or de-escalate intensity, and exit if needed, along with facilitator monitoring of participant well-being. \textit{After-game} safety involves structured debrief, emotional check-ins, and aftercare practices supporting integration. The \textit{community-level} dimension addresses systemic safety: codes of conduct, reporting mechanisms, power dynamics awareness, and community culture. Complementing this temporal model, \textcite{bowman2025safety} propose a participant-experience framework of ``Zones of Safety, Challenge, and Risk'': the \textit{comfort zone}, where participants engage with familiar material; \textit{growth edges}, where productive discomfort enables learning; and \textit{hard limits}, where content risks genuine harm. Applied to debate, the comfort zone encompasses familiar topics and preferred argumentation styles; growth edges correspond to challenging motions and arguing unfamiliar sides; and hard limits involve topics touching personal trauma or identity threats.

%Analyzed through this model, competitive debate's safety provisions are unevenly developed. Before-game safety is minimal: debaters receive motions with little content warning, and there is no calibration of emotional intensity. During-game safety relies almost entirely on the competitive frame itself---the alibi that ``it's just a debate'' is expected to provide sufficient protection, but this alibi erodes precisely when topics become personally relevant or when competitive stakes are high. After-game safety, as discussed in the debrief section, processes game performance but not emotional experience. At the community level, while organizations like ARDOR maintain ethical guidelines and equity officers, the competitive structure itself produces hierarchies and pressures that can undermine psychological safety.

%This analysis is significant because the thesis identifies competitive pressure and hierarchical dynamics as central problems---and these are fundamentally safety concerns. The methods chapter already commits to safety practices including pre-session briefing, post-session debrief, equity officer access, and the ability to exit without explanation. The safety framework from \textcite{bowman2025transformative} provides the theoretical grounding for these design choices and suggests additional interventions: content calibration before rounds, in-game safety mechanics adapted from larp practice, and structured emotional processing after rounds.

%\subsubsection{Designable Surfaces}

%A core design principle articulated in \textcite{bowman2025transformative} holds that in transformative game design, \textit{everything is a designable surface}: every rule, role, ritual, spatial arrangement, temporal structure, and social norm shapes participants' experience and therefore their potential for transformation. This principle reframes the designer's task from creating a single ``game'' to intentionally shaping a constellation of interconnected surfaces, each of which can be independently tuned to support or undermine transformative goals.

%Applied to competitive debate, the designable surfaces include: the motion or topic (what participants argue about, and how it is framed); side assignment (whether debaters choose or are assigned positions, and how this structures their relationship to the content); speech structure and timing (the sequence, duration, and format of speeches, which shape cognitive load and interaction patterns); the judge's role and evaluative criteria (who the judge role-plays as, and what standards they apply); feedback and debrief structure (how post-round processing is organized, as discussed above); competitive structure (how winners are determined, how status is distributed, what incentives the tournament creates); and community norms (the informal expectations, traditions, and social dynamics surrounding the formal game).

%Current debate formats have evolved these surfaces primarily through tradition, convention, and incremental competitive optimization rather than through intentional transformative design. The World Schools format's speech timings, for instance, derive from practical compromise rather than from analysis of optimal cognitive load for perspective-taking. Judging criteria emphasize competitive fairness---a worthy goal---but do not explicitly consider how evaluative standards shape the bleed patterns discussed above. This principle provides the theoretical warrant for the thesis's design project: if every surface shapes transformation, then every surface is an opportunity for intentional intervention.

\subsubsection{Zone of Proximal Development and Scaffolded Learning}

Vygotsky argues that higher cognitive functions develop through the scaffolding, by a more capable other, of challenges. When learners operate in the zone of proximal development where challenge is achievable only with guidance, they internalize new competences. \textcite{vygotsky1978mind}'s framework explains how debate creates naturally occurring scaffolding: power-pairing keeps challenges within this zone, while judges provide the guided instruction that feeds internalization. The debate cycle naturally embodies \textcite{kolb1984experiential} experiential learning model, with the debrief, which \textcite{crookall2014engaging} and \textcite{daniau2016transformative} identify as essential to transformation, built into debate through the oral critique. This scaffolding explains why the ``disequilibrium'' that \textcite{colbert1995enhancing} identifies as debate's cognitive engine produces growth rather than frustration: coaching, graduated competition, and structured feedback keep the cognitive challenge within the zone where productive accommodation occurs.

\subsection{Diagnosing Forensics Pedagogy Critiques through Transformative Game Design Mechanisms}

The critiques identified within forensics pedagogy, echoing the issues identified during the problem statement, can be productively reframed through the theoretical mechanisms of transformative game design, illuminating not only why these problems occur but also suggesting design-based interventions.

\subsubsection{Value Convergence as Corrupted Memetic Bleed}

What I coined as value convergence can be seen as a failure in memetic bleed: the format-level outcome it suppresses is \textit{value pluralism}, and the participant-level outcome is \textit{epistemic humility} (the same transformation seen from two sides). Whereas forensics literature locates the cause in a ``closed system'' of insulated tournament practice, the role-playing-game framing allows me to isolate a clear mechanism within that system: standard formats ask every judge to role-play the same archetype, the average unbiased person. This role is vaguely specified, conflating ``unbiased'' with ``lacking values'' and thus procedurally teaching that the unspoken values that debaters try to converge to in an attempt to make the game more objective and \DIFdelbegin \DIFdel{competitively unfrastrating }\DIFdelend \DIFaddbegin \DIFadd{less competitively frustrating }\DIFaddend are an acceptable baseline of society and thus correct.

Following \textcite{bowman2025transformative}, the container's effectiveness requires psychological safety for all participants. When judges face social consequences for heterodox opinions the alibi protecting honest judgment erodes and what bleeds into debaters is strategic equilibrium rather than appreciation for intellectual diversity. Debaters, optimizing for competitive success, do ``argument borrowing'', adapting strategically: arguments that succeed in elimination rounds or from high-status debaters become templates, and approaches aligned with perceived values propagate. These two effects in time lead community members to see said values as the only acceptable ones, integrating them in their identity. 

Design interventions might therefore (a)~diversify audiences debaters must persuade, (b)~strengthen judge alibi through structural features, and (c)~explicitly distinguish in debriefing between ``what won this debate'' and ``what constitutes good reasoning.''

\subsubsection{Information Density and Misaligned Procedural Bleed}

The critique of information-dense speaking represents misalignment in procedural bleed. When formats reward rapid delivery optimized for expert judges, debaters develop procedural habits adapted to highly specific audiences. The skills that bleed out are real but narrowly optimized, representing failure of transfer. Rubric-perception asymmetry creates the condition for misaligned bleed; the judge-as-expert role converts it into a bleed pattern; community governance propagates that pattern across a circuit. I will refer to the positive transformation along this axis as \textit{transferable persuasion}: rhetorical habits that survive the move from expert audiences to ordinary ones.

Design interventions might include (a)~evaluation criteria rewarding accessibility alongside substance, or (b)~judging paradigms simulating general audiences. (c)~By constructing vivid characters for judges to role-play, the bleed-in from character to judge can itself become memetic, the judge taking on the character's understanding while maintaining expert-level rule comprehension.

\subsubsection{Competitive Hierarchy as Compromised Container}

Hierarchy is downstream of a deeper mechanism. As foreshadowed in §3.1.1, sustained competitive pressure erodes the alibi enabling transformation: when tournament success conflates with status and self-worth, debate can no longer be treated as ``just a game,'' and stable hierarchies form along the axis of who has won. I will refer to the positive outcome along this axis as \textit{container safety}: a play environment whose stakes do not collapse into personal worth. \textcite{bowman2013social} documents social conflict within role-playing communities arising from power differentials between facilitators and players, competitive play dynamics, creative agenda disputes, and unprocessed bleed. These patterns parallel competitive debate: power differentials map onto hierarchies between experienced and novice debaters and competitive play produces status systems based on tournament success.

Emancipatory bleed requires a container where being taken seriously is possible; rigid hierarchies block this for those positioned as subordinate. The zone of proximal development requires calibrated challenges; competitive anxiety may prevent productive struggle or reflective observation.

Design interventions might (a)~distribute status more broadly, (b)~emphasize growth over victory, (c)~strengthen alibi by making explicit the distinction between game performance and personal worth, and (d)~incorporate safety practices adapted from larp design.

%\subsubsection{Implications for Design}

%This reframing generates specific predictions. Formats designed to maximize transformation should: (1) structure memetic bleed toward epistemic humility through debriefing practices connecting argumentative experiences to personal development; (2) align procedural bleed with transferable communication skills by strengthening judge alibi to role-play as specific general audience members; and (3) maintain container safety by preventing hierarchical success from conflating with personal worth.

\section{METHODS}
\label{sec:methods}

\subsection{Design Philosophy}

This research employs Research through Design (RtD) as methodological framework. RtD positions the design process itself as inquiry, where knowledge emerges through making and iterating on designed artifacts \parencite{coulton2017games,zimmerman2007research}. Following \textcite{zimmerman2003play}'s iterative methodology, each playtest generates data informing the next iteration.

\DIFaddbegin \DIFadd{More specifically, this is an }\textit{\DIFadd{evocative}} \DIFadd{design experiment in the sense of \mbox{%DIFAUXCMD
\textcite{waern2015experimental}}\hskip0pt%DIFAUXCMD
. They distinguish }\textit{\DIFadd{controlled}} \DIFadd{design experiments, which vary an isolated factor to answer descriptive or evaluative questions, from evocative ones, which are ``open for the fact that they can yield unexpected results'' and aim to understand ``the behaviour and experiences that a design choice will evoke in players.'' The present study is of the latter kind: it does not test whether the format passes a fixed threshold but asks how the design behaves and what it evokes in those who play it. Design research of this kind ``is judged by its relevance to design rather than its repeatability'' \mbox{%DIFAUXCMD
\parencite{zimmerman2007research}}\hskip0pt%DIFAUXCMD
.

}

\DIFaddend \subsection{Development Process}

Development unfolded through four iterative phases: brainstorming potential mechanics against identified problems; evaluating whether candidate designs addressed them without introducing new difficulties; testing surviving elements in pilot playtests with post-session debriefs; and synthesizing findings into the next iteration.

The process ran three iterative cycles. The first playtest (one debate) was held in March 2026. A second design iteration was drafted based on its findings. The second playtest was conducted in April 2026 with seven veteran debaters in a familiar informal setting, yielding a third design iteration. Both sessions are documented in Annexes A--G.

\subsection{Research Methods}

Each session ran approximately four hours, comprising workshops introducing rules and roles (deliberately omitting theoretical underpinnings to avoid demand characteristics), the debate, and a group debrief.

\subsubsection{Participant Observation}

Following \textcite{baker2006observation}'s taxonomy, I adopted a participant-as-observer role during both playtest sessions. My researcher status was known to participants; I served as facilitator during workshops and debriefs and as debater during the debate itself.

Observation focused on debate flow and behaviors relevant to the design\DIFdelbegin \DIFdel{predictions}\DIFdelend \DIFaddbegin \DIFadd{'s intentions}\DIFaddend : entertainment value, rule legibility, and elements distinct from standard formats.

Informed by the advice of \textcite{spradley1980participant}, I recorded field notes in two phases: condensed and expanded. Condensed notes, taken during debates, captured key moments, quotations, and behavioral observations using a brief written observation template. Expanded notes, written within 24 hours, elaborated on condensed notes, added context, and recorded analytical reflections. Both sessions were audio-recorded with participant consent to enable field note verification and capture verbal content missed in real time.

\subsubsection{Participants}

Participants in the first session were recruited through ARDOR Muntenia, the NGO managing the debate community in Bucharest, as part of a Jean Monnet high school debate club weekend meeting; 8 senior debaters aged 18--20 participated. Roles within the debate were assigned randomly. The second session was organized through personal contacts: participants were veteran debaters known to the researcher. The debate was held with 7 senior debaters aged 28 to 38, with roles assigned randomly as in the first session.

Although most Romanian debaters are minors, ethical guidelines made adult participants the best alternative. Both sessions used convenience sampling, appropriate for proof-of-concept RtD but not representative; the significant age and life-experience gap between cohorts limits the evidential weight of cross-session consistency.

\subsubsection{Post-Session Debriefs}

Semi-structured group debriefs complemented observation. After each play session an approximately one-hour conversation was held with each playtest's participants together: the researcher introduced pre-planned open questions derived from the theoretical predictions (Annex C), but allowed the conversation to develop organically, following participant-introduced threads and encouraging elaboration on emergent themes. Social conformity effects and dominant voices cannot be fully eliminated in group format but organic conversation development and active follow-up on quieter participants were used to partially mitigate this.

\subsection{Data Analysis}

Analysis followed a design-oriented evaluative approach. The primary purpose was \DIFdelbegin \DIFdel{not generating generalizable theory but }\DIFdelend informing ongoing design decisions and articulating how the designed format produces \DIFdelbegin \DIFdel{experiences consistent with its design predictions}\DIFdelend \DIFaddbegin \DIFadd{the experiences it intends}\DIFaddend , assessed against the researcher's extensive observational experience of existing formats. This aligns with experimental game design methodology \parencite{waern2015experimental}, where design experimentation becomes research when ``the aim is to understand something generic, some more fundamental aspect of game design.''

Before data collection, five evaluation fields were operationalized into a taxonomy of \textit{potential clues}: observable behaviors or reported experiences \DIFdelbegin \DIFdel{counting as evidence for or against }\DIFdelend \DIFaddbegin \DIFadd{bearing on how the format performs along }\DIFaddend each field's \DIFdelbegin \DIFdel{prediction}\DIFdelend \DIFaddbegin \DIFadd{dimension}\DIFaddend . Both positive and negative clues were specified at this stage. This pre-specification \DIFdelbegin \DIFdel{ensures the analysis remains }\DIFdelend \DIFaddbegin \DIFadd{serves as a discipline against confirmation bias, keeping the analysis }\DIFaddend responsive to disconfirming evidence rather than converging only on \DIFdelbegin \DIFdel{confirming }\DIFdelend \DIFaddbegin \DIFadd{flattering }\DIFaddend observations. The complete taxonomy in Annex~I defines what counts as evidence of effectiveness or ineffectiveness in each field.

Between sessions, rapid analysis identified what worked and should be preserved, what caused confusion, and what suggested new opportunities, driving the next iteration.

After all playtest sessions were concluded, systematic analysis of field notes and session transcripts proceeded in four steps. First, both sources were read in full and individual clues were extracted: each behavior observed or participant statement that could \DIFdelbegin \DIFdel{constitute evidence for }\DIFdelend \DIFaddbegin \DIFadd{bear on }\DIFaddend any field's \DIFdelbegin \DIFdel{prediction}\DIFdelend \DIFaddbegin \DIFadd{dimension}\DIFaddend . Second, each clue was assigned to one of the pre-specified taxonomy categories in Annex~I. Third, each clue was assessed for evidential strength using explicit criteria: behavioral evidence above self-report; unsolicited observations above responses to direct prompts; evidence consistent across both sessions above single-session occurrences; and novel behaviors above continuities, which may reflect participants' existing training rather than effects of the format. Fourth, the clues within each field were synthesized into a field-level \DIFdelbegin \DIFdel{judgment: the prediction was considered confirmed when positive clues outweighed negative ones }\DIFdelend \DIFaddbegin \DIFadd{account: a characterisation of how the format behaved along that dimension. Where positive clues dominated }\DIFaddend in aggregate strength, \DIFdelbegin \DIFdel{held with caveats when }\DIFdelend \DIFaddbegin \DIFadd{the account is correspondingly confident; where the }\DIFaddend evidence was mixed, \DIFdelbegin \DIFdel{and not confirmed when negative clues dominated or predicted positive clues were absent}\DIFdelend \DIFaddbegin \DIFadd{it holds the competing readings together rather than collapsing them; where predicted clues failed to appear or negative ones accumulated, it says so plainly}\DIFaddend . Throughout, the meaningful absence of predicted negative clues was \DIFdelbegin \DIFdel{treated as supporting evidence}\DIFdelend \DIFaddbegin \DIFadd{itself read as informative}\DIFaddend .

The \DIFaddbegin \DIFadd{main purpose of the pre-specified clue taxonomy is to maximise the breadth of relevant insight derived from the playtesting and to guard against the cognitive tunneling and inattentional blindness that may surface due to the structural circularity of authoring, facilitating, and observing the game at once. It also favours noticing unexpected results and eases reasoning about their implications: it makes explicit, a priori, the mechanisms that ought to have been at play but were not, or were not solely responsible.

}

\DIFadd{The }\DIFaddend five fields each encompass a spectrum of player behaviors pointing toward a specific characteristic the game must possess for transformative potential: (1) long-term engagement, (2) \DIFdelbegin \DIFdel{growth-eliciting play}\DIFdelend \DIFaddbegin \DIFadd{productive pressure}\DIFaddend , (3) useful growth direction, (4) container safety, and (5) epistemic humility.

Fields 1 (engagement) and 2 (productive pressure) are foundational: without voluntary engagement no transformational container can form, and without productive intellectual pressure inside it the container does not produce growth. Field 3 (growth direction) \DIFdelbegin \DIFdel{is the empirical test of the }\DIFdelend \DIFaddbegin \DIFadd{examines the }\DIFaddend information-density problem and \DIFdelbegin \DIFdel{of }\DIFdelend the broader transferable-skills claim. Field 4 (container safety) \DIFdelbegin \DIFdel{is the empirical test of the }\DIFdelend \DIFaddbegin \DIFadd{examines the }\DIFaddend competitive-hierarchy problem. Field 5 (epistemic humility) \DIFdelbegin \DIFdel{is the empirical test of the }\DIFdelend \DIFaddbegin \DIFadd{examines the }\DIFaddend value-convergence problem. \DIFdelbegin \DIFdel{The analysis led for each field to a verdict of either }\textit{\DIFdel{strongly confirmed}}%DIFAUXCMD
\DIFdel{, }\textit{\DIFdel{confirmed}}%DIFAUXCMD
\DIFdel{, }\textit{\DIFdel{confirmed in outcome with informative deviation in mechanism}}%DIFAUXCMD
\DIFdel{, }\textit{\DIFdel{confirmed at proof-of-concept level}}%DIFAUXCMD
\DIFdel{, or }\textit{\DIFdel{not confirmed}}%DIFAUXCMD
\DIFdelend \DIFaddbegin \DIFadd{Each field thus yields an account of how the format performed along its dimension}\DIFaddend .

Together, these steps \DIFdelbegin \DIFdel{can }\DIFdelend \DIFaddbegin \DIFadd{do not intend to }\DIFaddend demonstrate: (a) specific predicted effects of the format did or did not materialize at the behavioral or participant-report level; (b) specific mechanics appear connected to specific observed effects; and (c) observed effects align with the format's intended design goals. They cannot demonstrate persistence, superiority, or generalizability\DIFdelbegin \DIFdel{, appropriate limitations for proof-of-concept Research through Design where }\DIFdelend \DIFaddbegin \DIFadd{. As established above, such research is judged by its relevance to design rather than its repeatability \mbox{%DIFAUXCMD
\parencite{zimmerman2007research}}\hskip0pt%DIFAUXCMD
: its }\DIFaddend contribution lies in \DIFdelbegin \DIFdel{demonstrating that }\DIFdelend \DIFaddbegin \DIFadd{showcasing how }\DIFaddend systematic design can address identified problems and in articulating \DIFdelbegin \DIFdel{design principles }\DIFdelend \DIFaddbegin \DIFadd{transferable design principles or questions }\DIFaddend meriting further investigation. They also cannot eliminate the risk of motivated interpretation: all coding was performed by a single researcher who is also the format's designer, with no inter-rater verification; the pre-specified taxonomy and reflexive notes are partial mitigations but not substitutes for independent coding.

\subsection{Ethical Considerations}

This research adheres to Uppsala University's ethical guidelines for research involving human participants. All participants were informed of the study's purpose and data practices during and after recruitment, and could withdraw consent to any aspect of their participation or data use at any time. Peers whose input was used in developing this format were informed of context and credited. All data was stored securely in Uppsala University's encrypted OneDrive, separated by content type, and will be deleted after thesis presentation. All data has been anonymized and is not publicly available.

The ARDOR Ethical Guide and Trainer Safety Guide standards were followed to ensure psychological safety. Motions avoided personal trauma, and debaters could refuse any motion. Debaters could raise safety concerns with a third-party equity officer and report abusive or distressing behavior. During debates, any debater could exit without explanation.

All participants received a pre-session briefing and post-session debrief to address any concerns.

\subsection{Positionality and Reflexivity}

As stated in the introduction, I approach this research as a complete community member \parencite{adler1987membership}; the fifteen years shaping my knowledge also shape my values, assumptions, and emotional investments, creating several positionality challenges. \textit{Insider knowledge} enables nuanced recognition of subtle format problems but risks normalizing issues I cannot perceive from within the community. \textit{Structural circularity} is the most consequential: I authored the theoretical framework, designed the format to satisfy it, and then evaluated \DIFdelbegin \DIFdel{whether itdid}\DIFdelend \DIFaddbegin \DIFadd{it}\DIFaddend , predisposing every stage of the process toward confirmation, not only interpretation. \textit{Power dynamics} from my status as former Team Romania coach and successful debater may have shaped participant candor and debrief dynamics. \textit{Acquaintance proximity} in the second session deepened rapport and observation access while amplifying social desirability and confirmation risks.

Several measures were taken to partially mitigate these biases. The pre-specified clue taxonomy locked evaluation criteria before data collection, limiting post-hoc rationalization of findings. Reflexive notes at the end of every playtest actively sought disconfirming evidence. Awareness of power dynamics informed deliberate facilitation choices to encourage dissenting voices. No measure fully eliminates circularity: the pre-specified taxonomy was designed by the researcher whose predictions it tests, making all findings design-informed evidence rather than independent verification. Method-specific limitations are noted in their respective subsections above.


\section{DESIGN OF THE INSIDE OUT FORMAT}
\label{sec:design}

\DIFaddbegin \begin{figure}[H]
    \centering
    \includegraphics[width=\textwidth]{{otherdocuments/Screenshot 2026-06-10 at 20.45.54.png}}
    \caption{\DIFaddFL{Overview of the }\emph{\DIFaddFL{Inside Out}} \DIFaddFL{format as published on the project website, showing team structure, speech order, timing guide, and motion structure.}}
    \label{fig:format-overview}
\end{figure}

\DIFaddend \subsection{From Theory to Design}

The following sections argue why each structural element of the \DIFdelbegin \DIFdel{Inside Out }\DIFdelend \DIFaddbegin \emph{\DIFadd{Inside Out}} \DIFaddend format exists and how it serves the theoretical objectives, both for initial design decisions and for the two significant structural additions made in light of playtest. The full first-iteration ruleset is in Annex A; its chronological evolution in Annexes B--D.

\subsection{Design Summary}

\DIFdelbegin \DIFdel{Inside Out }\DIFdelend \DIFaddbegin \emph{\DIFadd{Inside Out}} \DIFaddend is a role-playing debate format for nine participants: eight debaters organised into four teams of two, and one Evaluator. Three teams, the \textbf{Value Defenders}, are each assigned one of three values held by a fictional named Subject. The fourth team, the \textbf{Subject team}, plays that person: a specific individual facing a genuine personal dilemma, whose role is not to convince but to weigh the arguments they hear and reach a decision.

A debate round follows a structured sequence. The round starts with a set of Rolling In rituals. Then the three Value Defender teams deliver first speeches (\textbf{A1, B1, C1}) in turn, each arguing for their assigned value and proposing a concrete, self-constructed decision it recommends. The Subject team then delivers their own first speech (\textbf{S1}), committing to a preliminary decision and explaining which arguments have so far moved them and why. The Defenders then deliver second speeches (\textbf{A2, B2, C2}), responding to the Subject's preliminary position and refining their cases. Finally, the Subject delivers a closing synthesis (\textbf{S2}): they name the decision they have reached and explain, transparently, what arguments and values actually moved them.

A separate person, the \textbf{Evaluator}, does not judge or decide the winner. They observe the round, then award \textbf{commendations}, named recognitions for specific demonstrated skills, and deliver oral feedback oriented toward growth rather than competitive outcome. The round closes with \textbf{Rolling Out}: participants de-role, shake hands, and the Evaluator invites a brief, optional check on whether anything in the debate felt personally close to home.

\DIFdelbegin \DIFdel{Inside Out }\DIFdelend \DIFaddbegin \emph{\DIFadd{Inside Out}} \DIFaddend departs from the World Schools (WS) format in three structural ways: the \textbf{team structure and roles} (four teams with distinct functions, plus a separate Evaluator); the \textbf{shape of motions} (personal, circumstantial dilemmas specifying the Subject's context and values, rather than abstract propositions); and the \textbf{rituals and special speeches} (Rolling In/Out, the interstitial S1, and S2) that frame the debate as genuine deliberation rather than adversarial performance. The following sections justify each departure.

\subsection{Design Rationale}

Each structural choice in \DIFdelbegin \DIFdel{Inside Out }\DIFdelend \DIFaddbegin \emph{\DIFadd{Inside Out}} \DIFaddend traces to the theoretical framework's diagnosis of how traditional formats produce corrupted bleed, misaligned skill transfer, and container compromise.

\subsubsection{Why Multiple Teams Defending Values?}

Traditional formats position two teams in binary opposition, a structure whose procedural rhetoric \parencite{bogost2007persuasive} teaches that every issue has exactly two sides. \DIFdelbegin \DIFdel{Inside Out }\DIFdelend \DIFaddbegin \emph{\DIFadd{Inside Out}} \DIFaddend replaces this with three Value Defender teams, each advocating for a different value held by the Subject. This restructuring addresses several identified problems simultaneously.

The multi-team structure \DIFdelbegin \DIFdel{produces }\DIFdelend \DIFaddbegin \DIFadd{is intended to produce }\DIFaddend a procedural rhetoric of complexity and humility: the format asserts through its rules that decisions involve balancing multiple legitimate human values, not choosing between right and wrong. This \DIFdelbegin \DIFdel{directly counters }\DIFdelend \DIFaddbegin \DIFadd{is designed to counter }\DIFaddend the corrupted memetic bleed driving value convergence by positioning value pluralism at the core of the conflict. Because what is assigned is a value rather than a policy position, the debater-character relationship is explicitly about inhabiting an identity facet rather than advocating an instrumental stance. And because the Subject is a person with specific background, circumstances, and values, the motion reinforces that decisions are products of context.

Crucially, the format assigns each team a \textit{value} rather than a \textit{solution}. As identity facets, values are emotionally resonant and generative of self-narrative in ways that instrumental solutions are not. Prescribing the value while leaving the solution open forces each team into authentic advocacy. This invites creative, contextualized, contextualizing and immersion-conducive role-play. Furthermore, the structure cultivates collaborative alongside combative skills. Value Defenders must persuade the Subject, a sympathetic listener making a genuine decision, not defeat an opponent. This \DIFdelbegin \DIFdel{corrects }\DIFdelend \DIFaddbegin \DIFadd{is meant to correct }\DIFaddend the misaligned procedural bleed by developing communication habits adapted to contexts beyond competitive debate, where genuine persuasion requires calibrating arguments to the specific beliefs and values of one's audience.

\subsubsection{Why Split Subject and Evaluator?}

In traditional formats, a single judge both determines the winner and provides feedback, entangling evaluation with adjudication. \DIFdelbegin \DIFdel{Inside Out }\DIFdelend \DIFaddbegin \emph{\DIFadd{Inside Out}} \DIFaddend separates these functions: the Subject makes decisions within the debate, while the Evaluator observes and offers commendations. This enacts framework predictions §3.2.1(a) (diversify audiences debaters must persuade) and (b) (strengthen judge alibi); the design's stronger move is to split the judge function rather than fortify a single role.

This split means winning the debate is not a goal in itself but a method for growing, symbolised through commendations rather than victories. Decision-making becomes synthesis rather than arbitration, modeling how real-world deliberation works and \DIFdelbegin \DIFdel{steering }\DIFdelend \DIFaddbegin \DIFadd{intended to steer }\DIFaddend memetic bleed toward epistemic humility. The Subject learns to synthesize; the Value Defenders learn to persuade collaboratively; and the Evaluator, freed from adjudicating competition, can focus on recognizing communicative virtue, maintaining the transformational container by decoupling status from competitive victory.

The Subject thus fills the role the framework attributes to judges in standard formats: \S\ref{sec:theory}'s prediction that a vivid character should be present for the format's listening role to take on is enacted here. The Subject also creates a natural accessibility incentive: if a Value Defender speaks incomprehensibly, their arguments are simply not weighed in the decision, a consequence that incentivises clarity. The Subject does dual duty: as biased character it sources value-pluralism bleed (§3.2.1); as non-expert listener it anchors accessibility (§3.2.2).

The format's final speech (S2) extends this logic to the round's conclusion. The Subject's closing synthesis requires them to name a concrete decision and explain what arguments and values actually moved them. The round ends not with one side claiming victory over the other but with an acknowledgment of how competing values were weighed and reconciled.

\subsubsection{Why Commendations and Oral Feedback?}

Replacing scores and victory with commendations moves the format from a paradigm of competing against others toward one of improving oneself. The format recognises that debaters can demonstrate excellence in fundamentally different ways, a principle made concrete by commendation categories spanning qualitatively distinct skills (``Believable Passion'', ``Argument Ingenuity'', ``Resilience'') and by a badge progression system that builds unique skill profiles over time. The theoretical framework predicts that design interventions for the compromised container should (a)~distribute status more broadly, (b)~emphasise growth over victory, and (c)~strengthen alibi distinguishing performance from worth (§\ref{sec:theory}).

The commendation system also resolves a persistent tension in modern debate formats between content density and persuasive discourse, precisely the misaligned procedural bleed identified in the theoretical framework. Commendations dissolve this conflict: debaters can pursue recognition for persuasive discourse in some speeches and recognition for academic rigor and complexity in others.

This growth orientation extends to the post-round debrief. In traditional formats, the oral critique centres on the adjudicator's decision: who won, why, and what would have changed the result, orienting debaters entirely toward competitive performance. \DIFdelbegin \DIFdel{Inside Out }\DIFdelend \DIFaddbegin \emph{\DIFadd{Inside Out}} \DIFaddend reframes this as feedback centred on demonstrated skills and developmental trajectory: the Evaluator explains the commendations awarded, articulates why particular moments demonstrated particular skills, and offers advice oriented toward growth rather than winning a future round.

\subsection{Design Additions from Playtesting}

Two structural additions from the first playtest merit justification. Other changes from that iteration (Points of Information, silent thinking time, motion enrichment, ritual simplification) are refinements of mechanisms already justified above.

\subsubsection{Why the Opening Subject Speech?}

An opening Subject speech (S0, here referred to as S1 in the final speech order) was added: a monologue as the character expressing why this decision is difficult and why each value is relevant.

The opening speech does not operate in isolation but in concert with the Rolling In and Rolling Out rituals that bracket every round. Rolling In functions as the container's liminal entry threshold \parencite{bowman2025transformative}; the Opening Speech completes it: the Subject's voice performing uncertainty as a specific character makes that character vivid to all participants at the moment the container opens. Rolling Out closes the arc.

By requiring the Subject to speak first as their character, the format also ensures that all subsequent debate proceeds from a particular person's situation rather than a universal proposition. The theoretical framework predicts (§\ref{sec:theory}, prediction (c) on information density) that a design intervention against misaligned procedural bleed should involve constructing vivid characters for the listening role, so that bleed-in from the character's situation becomes memetic rather than filtered through the listener's own background and values. \DIFdelbegin \DIFdel{Inside Out }\DIFdelend \DIFaddbegin \emph{\DIFadd{Inside Out}} \DIFaddend vests this listening role in the Subject (the Evaluator only post-evaluates and does not role-play). The Opening Speech is the structural mechanism that enacts this prediction: hearing the Subject speak as the character before any argument is made anchors the debate's frame to a specific person's dilemma rather than an abstraction. This makes the motion irreducibly personal, resisting the drift \textcite{goodnight1982spheres} identified from personal and social spheres toward the technical. It also replaces the first design's interstitial S1 (in which the Subject committed to a preliminary decision after the first Defender round).

\subsubsection{Why the Huddle?}

The Huddle is a structured interval after the first round of Defender speeches in which the Subject team may question any Defender, grant and revoke the floor, and prompt direct clashes, replacing the interstitial S1 entirely.

Theoretically, the Huddle simulates a council of competing interests. The dynamic is simultaneously cooperative and conflictual: Defenders can form alliances, acknowledge each other's arguments, and concede ground to gain the Subject's trust, while still competing for influence over the final decision. This procedural structure embeds a dynamic absent from binary debate formats: the collaborative-conflictual texture of real deliberative bodies, where stakeholders negotiate among themselves while simultaneously advocating to a decision-maker. The procedural rhetoric of the Huddle \DIFdelbegin \DIFdel{teaches }\DIFdelend \DIFaddbegin \DIFadd{is intended to teach }\DIFaddend that resolution does not arise from one side defeating the other, but from interests being forced to confront, clarify, and (sometimes) reconcile.

The Huddle also \DIFdelbegin \DIFdel{generates }\DIFdelend \DIFaddbegin \DIFadd{aims to generate }\DIFaddend a distinct pattern of procedural bleed \parencite{hugaas2024bleed} not produced by traditional debate. Traditional debate bleeds adversarial monologue habits: constructing extended causal chains under pressure, optimized for expert judges. The Huddle \DIFdelbegin \DIFdel{bleeds }\DIFdelend \DIFaddbegin \DIFadd{is intended to bleed }\DIFaddend deliberative habits instead: directed questioning, identifying where two arguments have not yet engaged and creating the conditions for them to do so, adapting argument in real time to a specific interlocutor's voiced needs. These are precisely the skills \textcite{cirlin1997sociological} associated with genuine persuasion, but the Huddle makes audience adaptation the primary structural mechanic of the interval rather than a secondary consideration. Such habits transfer to facilitated meetings, mediation, and collaborative decision-making: real-world contexts that traditional debate's adversarial procedural bleed works against rather than toward.

The remaining changes across the three iterations are refinements rather than new procedural rhetoric. Points of Information were introduced across all speeches to increase interactivity and fun; a two-minute silent thinking period was added before S2 for less experienced players; the motion's character description was enriched with fuller beliefs and past decisions to enable value-specific argumentation; Value Defenders were required to address the Subject in the second person; the role was renamed Inner Voices after participants coined the term organically in the second playtest; S2 instructions were extended to require emotional closure alongside intellectual synthesis (the Subject saying not only what decision they reached but how they now feel about it); and the rituals were compressed: Rolling In to a two-line call-and-response, Rolling Out simplified, and a Calling the Floor ritual added giving the Evaluator scripted phrases for each key moment.

\subsection{Synthesis and Design Tensions}

Taken together, the design decisions map onto the three identified problems as follows. Corrupted memetic bleed (value convergence) is addressed by the multi-team structure, which positions value pluralism at the core of the conflict; by the personal Subject, which grounds debate in context-dependent reasoning; and by commendations, which sever the evaluative feedback loop. Misaligned procedural bleed (information density) is addressed by the Subject as accessibility anchor, by the collaborative persuasion dynamic, and by commendation categories that reward persuasive discourse alongside analytical rigor. Compromised container (competitive hierarchy) is addressed by commendations replacing victory, by the Evaluator's decoupled role, and by oral feedback oriented toward growth. These three interventions track the three positive outcomes named in §\ref{sec:theory}: \textit{value pluralism / epistemic humility}, \textit{transferable persuasion}, and \textit{container safety}.

Several tensions in this design constitute productive questions for playtesting. Will competitive motivation survive without win-loss outcomes, or will debaters socialized in traditional formats experience commendations as insufficient stakes? The Subject team bears a unique cognitive burden in synthesizing three value streams, and whether this is productive challenge or overwhelming difficulty depends on experience level and motion complexity. The format's novelty (four teams, new roles, unfamiliar mechanics) may create barriers to entry for newcomers. Without victory as evaluative anchor, commendation judgments may feel more subjective.
\section{RESULTS}
\label{sec:results}

The research question\DIFaddbegin \DIFadd{, as posed in the introduction, }\DIFaddend breaks into two\DIFdelbegin \DIFdel{sub-questions: can a method of applying transformative game design theory to this domain be demonstrated, and does that method yield expected results?
}\DIFdelend \DIFaddbegin \DIFadd{: how transformative game-design theory can be operationalised into a debate format, and how participants respond to the format that results. The first is answered by the design process this thesis documents; the second by the playtests, read below.

}\DIFaddend 

\subsection{\DIFdelbegin \DIFdel{Demonstrating }\DIFdelend \DIFaddbegin \DIFadd{How }\DIFaddend the \DIFdelbegin \DIFdel{Method}\DIFdelend \DIFaddbegin \DIFadd{Format Was Designed}\DIFaddend : The \DIFdelbegin \DIFdel{Design }\DIFdelend Process as Contribution}

The first \DIFdelbegin \DIFdel{sub-}\DIFdelend question is answered by the \DIFdelbegin \DIFdel{thesis }\DIFdelend \DIFaddbegin \DIFadd{design process }\DIFaddend itself. The progression from theoretical analysis through iterative design to empirical testing \DIFdelbegin \DIFdel{operationalizes transformative game design theory in debate ; each design }\DIFdelend \DIFaddbegin \DIFadd{is the knowledge produced: it shows }\emph{\DIFadd{how}} \DIFadd{transformative game-design theory can be brought to bear on a debate format, with each }\DIFaddend iteration's reasoning \DIFdelbegin \DIFdel{is documented }\DIFdelend \DIFaddbegin \DIFadd{recorded }\DIFaddend in Annexes A, B, and G\DIFdelbegin \DIFdel{, retroactively collaging the method as narratively explicit and academically repeatable. }\DIFdelend \DIFaddbegin \DIFadd{. This progression offers a worked, transferable example of the method in operation.

}\DIFaddend 

\DIFaddbegin \DIFadd{A relevant mention on process is that the design did not deduce itself from theory although an a-posteriori justification was discovered; it emerged through iteration and intuition, and the theoretical analysis came alongside and after, illuminating why the moves that felt right worked, and where they might be pushed further. What the scholarly frame did add was a way of making the design's reasoning legible and ideally transferable, although at the expense of imposing limits on the number of playtest cycles, rather than fewer, would have sharpened the format further; the binding constraint was the difficulty of recruiting adult playtesters, not a limit of the approach.

}

\DIFaddend \subsection{\DIFdelbegin \DIFdel{Evaluating the Results}\DIFdelend \DIFaddbegin \DIFadd{How Participants Responded}\DIFaddend : Evidence from Playtests}

The second \DIFdelbegin \DIFdel{sub-question concerns whether the designed format shows signs of achieving its intended effects}\DIFdelend \DIFaddbegin \DIFadd{question is taken up through the playtests}\DIFaddend . Results were \DIFdelbegin \DIFdel{analyzed }\DIFdelend \DIFaddbegin \DIFadd{read }\DIFaddend using the clue-taxonomy method (Section~\ref{sec:methods}; full taxonomy in Annex~I). Each field corresponds to a \DIFdelbegin \DIFdel{prediction made }\DIFdelend \DIFaddbegin \DIFadd{design intention set out }\DIFaddend in Section~\ref{sec:design}\DIFdelbegin \DIFdel{and is assessed }\DIFdelend \DIFaddbegin \DIFadd{; the analysis characterises, }\DIFaddend as a constellation of clues\DIFaddbegin \DIFadd{, how the format behaved along the field's dimension and what it evoked in participants}\DIFaddend .

\subsubsection{Field 1: Engagement, Appeal, and Long-Term Viability}

Field~1 \DIFdelbegin \DIFdel{asks whether }\DIFdelend \DIFaddbegin \DIFadd{concerns how }\DIFaddend the format generates \DIFdelbegin \DIFdel{genuine }\DIFdelend voluntary engagement, \DIFdelbegin \DIFdel{by replacing }\DIFdelend \DIFaddbegin \DIFadd{specifically, how it replaces }\DIFaddend competitive appeal with imaginative engagement. Transformative potential cannot manifest in a game that players do not return to. The research question asks how game design theory can elevate a format's transformative potential, but this presupposes that participants inhabit the format long enough for its mechanisms to produce their effects. The transformational container forms and deepens through repeated participation; bleed arcs, skill accumulation, and epistemic shifts are all longitudinal processes that require the format to generate genuine voluntary re-engagement.

\DIFdelbegin \DIFdel{The prediction was }\textit{\DIFdel{confirmed}}%DIFAUXCMD
\DIFdel{: participant }\DIFdelend \DIFaddbegin \DIFadd{Participants }\DIFaddend were observed to display\DIFaddbegin \DIFadd{, }\DIFaddend and also self-reported\DIFdelbegin \DIFdel{high-levels }\DIFdelend \DIFaddbegin \DIFadd{, high levels }\DIFaddend of engagement and of desire to play again; negative signs of frustration, anxiety or disengagement were \DIFdelbegin \DIFdel{not observed or reported}\DIFdelend \DIFaddbegin \DIFadd{neither observed nor reported, }\DIFaddend and the following behaviours \DIFdelbegin \DIFdel{coroborate all of these}\DIFdelend \DIFaddbegin \DIFadd{corroborate this}\DIFaddend .

One important signal is the spontaneous adoption of ``inner voices'' as in-play address, a term absent from the ruleset. Two converging supports tighten the picture: animated post-session conversations extending beyond the structured debrief in both sessions, and players experimenting with uninstructed and unrewarded creative choices (accents, performance styles, personal props). Self-reports add a third strand: the format seen as ``surprisingly liberating,'' the Huddle phase ``very very fun,'' and similar positive characterizations, all offered unsolicited across both sessions.

A differential engagement pattern \DIFdelbegin \DIFdel{was observed and is itself evidence that emancipatory bleed \mbox{%DIFAUXCMD
\parencite{kemper2017battle, kemper2020wyrding} }\hskip0pt%DIFAUXCMD
is operating as theorized}\DIFdelend \DIFaddbegin \DIFadd{emerged that the design had not set out to produce}\DIFaddend . In the sessions observed, women and debaters with less successful traditional-format track records expressed consistent preference \DIFaddbegin \DIFadd{for the format}\DIFaddend , while three of the fifteen participants across both sessions (all competitive men with strong traditional-format track records) expressed reservations, one describing the format as ``anticlimactic'' and missing the ``ego-boosting'' experience of decisive victory. The \DIFdelbegin \DIFdel{format was not designed to serve the highly competitive profile that traditional formats already serve; its appeal to voices suppressed by traditional formats is what the frameworkpredicts when status-competition is replaced by imaginative investment under a strengthened alibi \mbox{%DIFAUXCMD
\parencite{bowman2021magic, deterding2017alibis}}\hskip0pt%DIFAUXCMD
}\DIFdelend \DIFaddbegin \DIFadd{pattern traces to a specific design move: removing the win/loss outcome withdrew the status-competition on which traditional standing is built, and in its absence the voices that standing usually crowds out came forward. Read through the framework, this is emancipatory bleed \mbox{%DIFAUXCMD
\parencite{kemper2017battle, kemper2020wyrding} }\hskip0pt%DIFAUXCMD
displacing status-competition, surfaced as an observed asymmetry in who the format served, not as a property the design was built to exhibit}\DIFaddend .

Three other clues \DIFdelbegin \DIFdel{were factored in the verdict}\DIFdelend \DIFaddbegin \DIFadd{qualify the picture}\DIFaddend : the commendation system produced a subtler negative clue, with Evaluators in both sessions finding commendation assignment cognitively overwhelming at small informal scale. The system was designed for tournament use where Evaluators have prior familiarity with the badge list; the difficulty observed here is a calibration issue specific to off-tournament context, not a structural failure. The 40-minute onboarding brief confirms a real learning curve, while both sessions producing engaged play within a single session confirms the cost is onboarding, not fundamental inaccessibility. The Rolling Out ritual recurred as a disengagement signal across both sessions. The redesigned version in the third iteration (Annex~H) addresses this directly but remains untested in this thesis cycle.

\subsubsection{Field 2: Productive Pressure Mechanisms}

Field~2 \DIFdelbegin \DIFdel{asks whether }\DIFdelend \DIFaddbegin \DIFadd{concerns how }\DIFaddend productive intellectual pressure \DIFdelbegin \DIFdel{persists after removal of }\DIFdelend \DIFaddbegin \DIFadd{behaves once }\DIFaddend win/loss outcomes \DIFaddbegin \DIFadd{are removed: what sustains it, if anything, in the absence of victory}\DIFaddend . Growth through play requires genuine intellectual and social challenge. Transformation through structured activity depends on productive pressure: the kind of challenge that stretches participants' capacities and motivates them to develop. Without productive pressure, repeated play (\DIFdelbegin \DIFdel{confirmed by }\DIFdelend \DIFaddbegin \DIFadd{established in }\DIFaddend Field 1) would be pleasant but developmentally inert.

The \DIFdelbegin \DIFdel{prediction }\DIFdelend \DIFaddbegin \DIFadd{design intention }\DIFaddend was that pressure would be redirected rather than removed, with the commendation architecture replacing victory as the evaluative anchor. \DIFdelbegin \DIFdel{The prediction was }\textit{\DIFdel{confirmed in outcome with an informative deviation in mechanism}}%DIFAUXCMD
\DIFdel{: pressure persisted }\DIFdelend \DIFaddbegin \DIFadd{What the playtests show is both an outcome and a surprise. Pressure did persist }\DIFaddend at high levels across both sessions \DIFdelbegin \DIFdel{, but the }\DIFdelend \DIFaddbegin \DIFadd{but its }\DIFaddend operative driver was \DIFdelbegin \DIFdel{an emergent social-recognition dynamic alongside }\DIFdelend \DIFaddbegin \DIFadd{not only }\DIFaddend the formal commendation system the design \DIFdelbegin \DIFdel{predicted}\DIFdelend \DIFaddbegin \DIFadd{installed; it was an emergent social-recognition dynamic the design had not anticipated, surfacing a mechanism not intentionally built in, exactly the kind of result an evocative experiment is positioned to catch}\DIFaddend .

On the outcome side\DIFdelbegin \DIFdel{(was pressurepreserved?)}\DIFdelend \DIFaddbegin \DIFadd{, when observing markers of positive, growth-orienting pressure}\DIFaddend , the first session produced elaborate, nuanced solutions with careful value-priority arguments and passionate speeches, and the second produced contextually rich, roleplay-oriented speeches with uninstructed worldbuilding detail. All of these \DIFdelbegin \DIFdel{prove }\DIFdelend \DIFaddbegin \DIFadd{evidence }\DIFaddend intellectual effort, triggered by a pressure to perform. The risk that emotional expressiveness might displace analytical rigor did not materialize: practical reasoning and value hierarchies remained present alongside the emotional engagement. A first-session self-report (``architecting solutions'' as ``very complex and interesting strategically'') aligns with the behavioral picture. A novel positive clue from the second session adds force: participants spontaneously initiated a group conversation about how to diminish another Inner Voice's contribution without being dismissive, a growth discussion standard debate does not produce because the skill is not rewarded there. The absence of apathy in both sessions, with no participant reducing effort in the win/loss vacuum, weighs as supporting evidence.

On the mechanism side, \DIFdelbegin \DIFdel{a surprising operative driver , }\DIFdelend \DIFaddbegin \DIFadd{the operative driver }\DIFaddend alongside the commendation architecture \DIFdelbegin \DIFdel{, }\DIFdelend was peer recognition: participants wanted recognition from fellow Inner Voices and from the Subject. This \DIFdelbegin \DIFdel{is an informative deviation showing }\DIFdelend \DIFaddbegin \DIFadd{shows }\DIFaddend the operative rule can be a social rule emergent from the formal structure \DIFdelbegin \DIFdel{, and it }\DIFdelend \DIFaddbegin \DIFadd{rather than the formal scoring rule itself: an instance of the second-order character of game design, in which play arises not from the rules as written but from the session those rules give rise to \mbox{%DIFAUXCMD
\parencite{salen2004rules}}\hskip0pt%DIFAUXCMD
. It }\DIFaddend converges with the Field~4 finding that informal social reward did the work of formal commendation when the latter was unreliable.

One playtest observation brought concern over the Subject role's scaffolding. In the first playtest, the Subject described not knowing ``what to say,'' uncertain how to hold emotional investment and evaluative neutrality at once; this was addressed in the second iteration's instructions. That session also offered a partial resolution: the Subject who navigated the role most successfully was the one who leaned into roleplay rather than maintaining analytical detachment. Inhabiting the character, responding emotionally, and asking from genuine curiosity resolved the cognitive tension between roleplay and evaluation, consistent with the broader argument in Section~\ref{sec:background} that debate is fundamentally a form of role-playing.

\subsubsection{Field 3: Growth Direction -- Useful Skills and Perspectives}

Field~3 \DIFdelbegin \DIFdel{is the empirical test of the }\DIFdelend \DIFaddbegin \DIFadd{takes up the }\DIFaddend information-density problem (§1.3, §2.4.3). It \DIFdelbegin \DIFdel{asks whether the format produces growth in }\DIFdelend \DIFaddbegin \DIFadd{concerns how the format shapes growth across }\DIFaddend the spectrum of useful skills documented in Section~\ref{sec:background} as the basis of debate's transformative reputation: critical thinking, communication and public-speaking proficiency, and empowerment of voices that traditional formats suppress. \DIFdelbegin \DIFdel{Without this confirmation the format might }\DIFdelend \DIFaddbegin \DIFadd{A format could }\DIFaddend be enjoyable and demanding \DIFdelbegin \DIFdel{but }\DIFdelend \DIFaddbegin \DIFadd{yet, }\DIFaddend through its demands\DIFdelbegin \DIFdel{it would not build useful skills}\DIFdelend \DIFaddbegin \DIFadd{, build little of lasting use; the question is which way this format's demands point, and what they build}\DIFaddend .

The \DIFdelbegin \DIFdel{prediction was that the format would }\DIFdelend \DIFaddbegin \DIFadd{design intention was to }\DIFaddend preserve standard debate's analytical gains while extending them along axes traditional formats develop weakly: persuasion addressed to a sympathetic non-expert listener\DIFaddbegin \DIFadd{, }\DIFaddend and accessibility to participants without prior debate training. The \DIFdelbegin \DIFdel{prediction was }\textit{\DIFdel{confirmed}}%DIFAUXCMD
\DIFdelend \DIFaddbegin \DIFadd{playtests show how that redirection played out, across three strands}\DIFaddend .

Critical thinking and analytical rigor were preserved, not displaced by the format's emotional register. Both sessions produced elaborate, nuanced solutions; participants navigated value hierarchies and inter-value tensions without prompting; practical reasoning sat alongside roleplay. First-session self-reports of ``architecting solutions'' as ``very complex and interesting strategically'' converge with the behavioural picture. Further self-reports articulating learning ``to think about the context of decisions,'' to be ``more human and more honest about human nature,'' and to recognise that ``sometimes decisions are personal, not universal'' indicate the analytical frame did not narrow but expanded toward contextually grounded reasoning.

Communication and public-speaking proficiency, the dimension on which traditional formats are most often critiqued, developed along axes traditional debate actively suppresses through its rewards for information density. Players in the second session adopted second-person address, direct eye contact with the Subject, and contextual worldbuilding from the opening of the debate, all uninstructed: audience adaptation became the primary rhetorical mode. Expressive range broadened in unscripted directions: accents, performance styles, a soundtrack played behind a Subject speech, worldbuilt characters with emotional stakes. Participants explicitly identified this dimension, noting that the format develops ``conversational skills and empathy without necessarily having this educational background.''

That formulation also names the empowerment dimension. Voices the traditional formats suppress (in this sample, women and debaters with less successful traditional-format track records) consistently expressed preference for the format, and Field~1's differential engagement pattern reads here as confidence development. Across both sessions, play diverged from standard-debate conventions speech by speech, reading these adjustments as effects of the format. The two predicted negative clues, unhelpful play rewarded and hollow performative emotionality, did not manifest in either session.

\subsubsection{Field 4: Container Safety and the Erosion of Hierarchy}

Field~4 \DIFdelbegin \DIFdel{asks whether }\DIFdelend \DIFaddbegin \DIFadd{concerns how }\DIFaddend removing win/loss outcomes \DIFdelbegin \DIFdel{reduces adversarial dynamics.  Fields 1--3 confirm that individual participants engage, are challenged, and develop useful skills within a session; Field 4 tests whether }\DIFdelend \DIFaddbegin \DIFadd{affects adversarial dynamics, how }\DIFaddend the format's structural logic reshapes the social environment those sessions occur within: \DIFdelbegin \DIFdel{whether }\DIFdelend \DIFaddbegin \DIFadd{how far }\DIFaddend it can convert a community culture into one that sustains rather than undermines the container.

The \DIFdelbegin \DIFdel{prediction was }\textit{\DIFdel{confirmed within the playtest window}}%DIFAUXCMD
\DIFdel{. Persistence in a mature circuit, where established status would be staked on performance, remains untested however.
}%DIFDELCMD < 

%DIFDELCMD < %%%
\DIFdel{The }\DIFdelend clearest behavioral signal is an emergent norm against adversarial rebuttal: in the second playtest, participants spontaneously articulated that ``outright attacking others' points feels wrong, disrespectful, and unnecessary'', a norm absent from the ruleset, arising from the format's structural logic. This is procedural rhetoric in the precise sense \textcite{bogost2007persuasive} theorizes: rebuttal carries less competitive benefit while contextual contribution does, inverting the rhetorical incentive structure of standard debate. Reduced anxiety converges from another angle: across both sessions, participants independently reported the format as ``much much less anxiety-inducing,'' described a ``feeling of safety regardless of result,'' and noted that ``everyone distinguished himself by his own style'': unsolicited, multi-source reports across two independent groups, consistent with the transformational container forming as theorized \parencite{bowman2021magic, bowman2025transformative}.

Brave, unconventional play forms a third strand: vulnerable, emotionally intimate approaches in the first session and roleplay-heavy delivery, soundtrack-opened Subject speeches, and worldbuilt characters with emotional stakes in the second drew immediate gestural reward (laughing, nodding, leaning forward). This experimentation was possible because the Inner Voices framing operates as Deterding's \parencite{deterding2017alibis} alibi, distancing performer from position and removing the identity threat that constrains risk-taking in win/loss formats. The fourth indicative is what was \textit{not} observed: none of the three predicted negative clues (distress, arrogant or unempathetic behavior, praise of such behavior) manifested across either session, despite participants being experienced competitive debaters with established adversarial habits and despite the second session's moments of genuine emotional intensity (see Field~5).

One further qualification: the commendation system's calibration immaturity at small scale, noted in Field~1, recurs here in a more theoretically revealing form: where the formal commendation channel was unreliable, the room's gestural and verbal social reward did the same evaluative work, a second instance of the social-recognition mechanism named in Field~2.

\subsubsection{Field 5: Epistemic Humility}

Field~5 \DIFdelbegin \DIFdel{is the empirical test of the }\DIFdelend \DIFaddbegin \DIFadd{takes up the }\DIFaddend value-convergence problem. A format could \DIFdelbegin \DIFdel{confirm }\DIFdelend \DIFaddbegin \DIFadd{satisfy }\DIFaddend all four earlier fields while still producing participants who are enthusiastic, pressured, skilled, and socially safe but epistemically unchanged, still convinced that their values are objectively correct and others' are deficient. It \DIFdelbegin \DIFdel{tests the deepest predicted transformation: that the design produces }\DIFdelend \DIFaddbegin \DIFadd{concerns the deepest of the format's intended transformations: }\DIFaddend a shift in how participants hold their own beliefs, as an emergent consequence of inhabiting the format's structural logic.

The \DIFdelbegin \DIFdel{prediction was }\textit{\DIFdel{confirmed at proof-of-concept level}}%DIFAUXCMD
\DIFdel{: the analysis centers on the single highest-strength clue in the dataset and is supported by behavioral, articulated, and absent-as-evidence strands.
}%DIFDELCMD < 

%DIFDELCMD < %%%
\DIFdel{The }\DIFdelend focal observation is an unplanned bleed event in the second session: a participant disclosed privately to the researcher that the scenario in play resonated with a real-life situation they were facing that same day. This is the highest-strength clue in the dataset because three weighting principles align: it is behavioral, private, and it instantiates the transformational mechanism directly, with the format's content crossing the play/life boundary.

\DIFdelbegin \DIFdel{As proof-of-concept evidence, it shows the mechanism functioning exactly as theorized.
}\DIFdelend 

Three in-play events correlate to fill out the picture. Synthesis orientation emerged across both rounds of Inner Voices speeches, appearing in the first round and intensifying in the second: spontaneous drift toward converging on a solution despite no rule incentivizing it, the clearest behavioral indicator that the format's structural logic orients participants toward integration rather than dominance. Players in the second session organically renamed themselves ``Inner Voices'' during play, positioning themselves as parts of a shared internal deliberation rather than as opposing sides (subsequently formalized in the third iteration). The second-session debrief \DIFdelbegin \DIFdel{produced a meta-confirmation}\DIFdelend \DIFaddbegin \DIFadd{coroborates this}\DIFaddend : participants independently identified that the S0 institutionalizes surfacing personal biases openly before argumentation, reframing bias as perspective rather than contamination, an analysis they arrived at without prompting, which is itself evidence the mechanism functioned: the format made its own epistemic logic visible to its players. First-session self-reports converge from another angle: one participant articulated the attitudinal shift unprompted (``every team says valuable stuff and it is important more to see what is important to us''), and others observed that arguing someone is ``morally wrong'' for a personal decision feels ``a bit absurd'' in this format.

\DIFdelbegin \DIFdel{The strands together trace one phenomenon: value pluralism replacing value convergence as the format's procedural lesson \mbox{%DIFAUXCMD
\parencite{bogost2007persuasive}}\hskip0pt%DIFAUXCMD
.
}\DIFdelend 

\DIFdelbegin \DIFdel{Three qualifications limit what can be claimed. The bleed event is a single observation and cannot demonstrate persistence. Both sessions used small samples of experienced debaters; whether the same configuration produces the same effect at scale is what the next design cycle would test. No longitudinal follow-up was conducted; the durability of the shift beyond the playtest window is unestablished. By contrast, the absence of unequivocal condemnation of positions or belief-holders, or praise of such condemnation, holds across both sessions despite motions involving genuine ethical weight and participants practiced in adversarial argumentation, strengthening confidence beyond what the positive clues alone would justify.
}\DIFdelend %DIF > Three qualifications limit what can be claimed. The bleed event is a single observation and cannot demonstrate persistence. Both sessions used small samples of experienced debaters; whether the same configuration produces the same effect at scale is what the next design cycle would test. No longitudinal follow-up was conducted; the durability of the shift beyond the playtest window is unestablished. By contrast, the absence of unequivocal condemnation of positions or belief-holders, or praise of such condemnation, holds across both sessions despite motions involving genuine ethical weight and participants practiced in adversarial argumentation, strengthening confidence beyond what the positive clues alone would justify.

\DIFdelbegin \subsection{\DIFdel{Proof of Concept Assessment}}
%DIFAUXCMD
\addtocounter{subsection}{-1}%DIFAUXCMD
\DIFdelend %DIF > \subsection{Proof of Concept Assessment}%

\DIFdelbegin \DIFdel{Overall, all clues point towards Inside Out being and being perceived as an interesting and promising deviation from the standard of debate formats. Because of various elements of novelty, players were enthusiastic about playing it some more, tried hard when playing and employed useful-to-build skills (and also different ones from the standard of most debates). Clues allow me to me hopeful that removing a clear win-loss result and adding the commendation system might be a good first step in making the debate community less hierarchical. The playtests also offered more clues pointing towards it maybe managing to allow participants to achieve epistemic humility than honestly anticipated. All of these can be summarized as the game becoming an acceptable (at least by myself as an experienced game-design practicioner) proof of concept and a starting point for a longer design process.
}\DIFdelend %DIF > Overall, the clues converge on \emph{Inside Out} as an interesting and promising deviation from the standard of debate formats, both in itself and as perceived by its players. Drawn in by various elements of novelty, players were enthusiastic to play it again, tried hard while playing, and exercised useful-to-build skills---including ones the standard of most debates does not develop. The clues leave me hopeful that removing a clear win/loss result and adding the commendation system may be a good first step toward a less hierarchical debate community, and they pointed toward the format fostering epistemic humility more than I had honestly anticipated. What this amounts to is not a verdict but a starting point: an acceptable proof of concept---acceptable, at least, to my own judgment as an experienced game-design practitioner---and the opening of a longer design process, together with a body of design knowledge about how transformation might be engineered into an argumentative, interpersonal format.%

 % The most interesting design finding is a refinement of that account itself. Field~2's mechanism deviation (social recognition driving pressure rather than commendations) and Field~4's converging observation (informal social reward doing the work of formal commendation when calibration was unreliable) together suggest that rule-structure shapes behavior, but the operative rules need not be the formal scoring rules: they can be social-recognition rules that emerge from the formal structure. 

\subsection{Limitations}

Several limitations bound what these results can claim. First, the original design envisioned approximately four iterative playtest sessions; two were completed within the thesis timeline. \DIFdelbegin \DIFdel{Both yielded substantive observational and debrief data sufficient for proof-of-concept assessment, but the }\DIFdelend \DIFaddbegin \DIFadd{The }\DIFaddend reduced iteration count narrows the breadth of design refinement reflected here. Second, two recurring design gaps, Subject role scaffolding (Field~2) and the Rolling Out ritual's immersion-breaking effect (Field~1), were partially addressed across iterations but remain open questions; both should be prioritized in the next design cycle. Third, both playtest groups consisted of experienced debaters already familiar with competitive formats; the format's performance with novice participants remains untested.

\DIFdelbegin \DIFdel{The methodological constraints articulated in Section~\ref{sec:methods} continue to bound inference: the constellation across the five fields supports a coherent design account but cannot demonstrate persistence, superiority, or generalizability.
}\DIFdelend 

\section{Discussion}
\label{sec:discussion}

\subsection{What Including Debate Reveals About Role-Playing Games}
\label{sec:rpg-insights}

The meeting of playtest and theoretical expectations gave birth to a handful of insights:

Firstly, on alibi. The framework treats alibi as one concept with two glosses: \textcite{deterding2017alibis}'s \textit{justification-alibi} (the rationale an adult adopts as licence to play) and \textcite{bowman2025transformative}'s \textit{performance-alibi} (the in-play permission to act under cover of character). The playtests show these are coupled, not parallel. The Rolling Out closing ritual, a framing device of the kind most LARPs use without difficulty, was reported as ``robotic'' in the first session and met with mumbling in the second. Yet the same participants spontaneously adopted ``Inner Voices'' as in-play address. Both moves were formally licensed by the same in-play permission. What discriminated them was justification: chanting felt decorative under a rationale that read play as ``learning argumentation,'' but natural where the rationale was ``inhabiting a character.'' The justification an adult brings to the table places a prior on the performance-alibis they will exercise once in-play. The contribution to the framework is to treat the adult-justification surface not as a separate designable surface from the in-play one, but as a filter on it.

The second refinement concerns evaluator roles in relation to value convergence. The mechanism named in §\ref{sec:theory} (judges asked to role-play the same ``unbiased average person'' archetype) was confirmed by playtest: \DIFdelbegin \DIFdel{Inside Out}\DIFdelend \DIFaddbegin \emph{\DIFadd{Inside Out}}\DIFaddend 's structural choices, prescribing for the Subject a clearly biased character and asking players to voice their interpretative biases during the S1 speech, led to players' realization of the need for tolerance towards others' beliefs and circumstantial decisions. This is a clean demonstration of the thesis's methodological claim: only treating the judge as a roleplayed character names the parameter that needed redesign, where the standard reform vocabulary of bias and calibration used in the debate community does not manage that.

This is further substantiated by the apparent tension in the Subject role (emotional investment and evaluative neutrality) proving easier to resolve by leaning into \textit{steered} roleplay \parencite{montola2015steering, pohjola2015steering}, mirroring my and \textcite{seo2022good}'s thesis that the cognitive overload of debate is only manageable through genuine character inhabiting. Steered roleplay proves thus to be employable in the pursuit of performance and memetic bleed, itself a teachable capacity that can scaffold the long-term cognitive abilities forming debate's transformative potential.

The third refinement concerns pressure. Peer recognition alongside formal commendations, was the operative driver, mirroring the pressure found in LARP communities without a formal performance metric, and raising the question of whether container safety is better served by hushed judgement or by open, constructively shaped critique.

\DIFaddbegin \DIFadd{A fourth refinement concerns emancipatory bleed and the shape of competition itself. Traditional debate stages a binary, zero-sum contest. }\emph{\DIFadd{Inside Out}}\DIFadd{'s non-binary structure and as generated RtD knowledge it suggests that designing }\emph{\DIFadd{against}} \DIFadd{the binary, victory-centred grammar of an activity can itself be a feminist, or more broadly emancipatory, design move. In an act of nuancing subversion, the differential response fell neatly along a masculine/feminine axis, which means the design move revealed a binary whilst seeking to subvert binarity altogether. This invites to deliberation on how deliberative games can be designed with anticipated gender-differential responses in mind.

}

\DIFaddend Treating debate as a special (evaluative, conflictual, performance-demanding) category of role-playing game has thus produced insights that travel back to other role-playing games as well.

\subsection{Directions for Continued Study}
\label{sec:further-study}

The playtests open four research directions, each pointing to an established literature outside transformative-game-design ready to be imported.

First, although the commendation system was received enthusiastically, evaluators struggled to assign commendations consistently and named a risk of ``farming'': stylised play optimised for commendation coverage rather than development. The commendations were designed to be a small part of a larger design of a long term progression tracking system. Only once the design of this entire surface, based on theories from gamification, game-progression design, game-economy design and balancing is finalized and implemented at the level of community adoption, can the commendation system be assessed properly.

Second, participants' reconstruction of how the information-density meta took root in British Parliamentary and contaminated World Schools located the cause not in rules but in who held social power in the circuit once the founding teachers disengaged. This refines the framework's judge-as-expert mechanism (§\ref{sec:theory}): the judge role is where bleed misaligns; governance is where the misalignment is sustained over time. Governance, how rules are interpreted, evaluators trained, motions authored, and meta drift sustained, is a designable surface above the format that my efforts have yet to touch. 

Third, naming peer recognition as a pressure category opens questions the playtests cannot answer: how it scales beyond the bounded session, how it interacts with progression, and what failure modes (clique formation, in-group bias, the homogenising drift a unitary evaluator role produces in standard debate) it carries.

Fourth, \DIFdelbegin \DIFdel{Inside Out}\DIFdelend \DIFaddbegin \emph{\DIFadd{Inside Out}}\DIFaddend 's design moves are decomposable: the commendation system, the split between an evaluated Interlocutor and an unevaluated external Evaluator, three teams, a judge who roleplays a named character, and teams prescribed values but granted creative freedom over their expression are each candidates for single-mechanism transplantation, inviting a methodology that treats transformative mechanisms as transplantable units rather than constituents of holistic experiences.

Importing progression theory, community management, peer-recognition dynamics, and mechanic-level analysis, together with specific mechanical refinement, coalesce into the research programme that can not only sharpen this thesis but also benefit the entire field of transformative game design study.

\section{Conclusion}

Alfred Snider first recognized debate as a game and championed its transformative power \parencite{snider1982gaming,snider1984gaming}, yet none who followed applied game-design know-how to the question his argument makes inevitable: how can we make it a better game? This thesis takes one step into that question.

This thesis \DIFdelbegin \DIFdel{answers a ``how can'' question in three stages: whether }\DIFdelend \DIFaddbegin \DIFadd{explores the topic in three movements: establishing that }\DIFaddend the conceptual tools apply (Sections \ref{sec:background}--\ref{sec:theory}: debate is a transformative role-playing game\DIFaddbegin \DIFadd{, }\DIFaddend and game design can be brought to bear on it); \DIFdelbegin \DIFdel{whether }\DIFdelend \DIFaddbegin \DIFadd{carrying }\DIFaddend a concrete design \DIFdelbegin \DIFdel{can be carried }\DIFdelend through (Sections \ref{sec:methods}--\ref{sec:design}: the \DIFdelbegin \DIFdel{Inside Out }\DIFdelend \DIFaddbegin \emph{\DIFadd{Inside Out}} \DIFaddend format and its iterative development); and \DIFdelbegin \DIFdel{whether the attempt produced value }\DIFdelend \DIFaddbegin \DIFadd{reading what that design evoked in play }\DIFaddend (Results and Discussion: \DIFaddbegin \DIFadd{a }\DIFaddend clue-by-clue \DIFdelbegin \DIFdel{confirmation and }\DIFdelend \DIFaddbegin \DIFadd{account and its }\DIFaddend theory-facing extensions).

\DIFdelbegin \DIFdel{In its attempt to contribute to the broader body of knowledge of the field of transformative game design, it does not manage to answer the research question more than at the level of proof of concept. It presents a design and it draws out the little playtest data available an analysis that argues how the design shows promise. 
}%DIFDELCMD < 

%DIFDELCMD < %%%
\DIFdel{To my own body of knowledge the thesis contributes most in becoming an artifact of academic practice and resistance to frustration. Although Inside Out does constitute a first}\DIFdelend \DIFaddbegin \emph{\DIFadd{Inside Out}} \DIFadd{is a first, }\DIFaddend and not strictly necessary\DIFaddbegin \DIFadd{, }\DIFaddend step in a much longer process of reform, whose later stages involve refinement, community adoption, governance design, and the slow work of shifting collective habits\DIFdelbegin \DIFdel{, the utility of the scholarly wrapper as part of which it was developed failed yet to shine bright enough because of a few reasons. The deductive arc this thesis narrates is but a necessary pretence. The actual design emerged intuitively and was retrofitted to theory; that it held testifies to the scholarly aplication of the argumentative skills debate built in me, not to the merits of this theory-born deductive design approach. One cannot but wonder how far the design would have evolved if effort were to have been invested in testing and refining and not in overcoming the difficulty of finding only over-18 playtesters and in wordsmithing convoluted theoretical analysis.
}%DIFDELCMD < 

%DIFDELCMD < %%%
\DIFdel{What this thesis carries forward, then, is neither a finished format nor great insights but three boons for whose utility one can argue. Firstly, my emotional momentum required }\DIFdelend \DIFaddbegin \DIFadd{. What the thesis carries forward into that process is threefold: a working design and the documented reasoning behind it, ready for the next cycle; the momentum }\DIFaddend to pursue the \DIFdelbegin \DIFdel{subsequent }\DIFdelend stages of reform \DIFdelbegin \DIFdel{that the design itself }\DIFdelend \DIFaddbegin \DIFadd{the design alone }\DIFaddend cannot deliver; \DIFdelbegin \DIFdel{secondly, a misguided yet instrumentalizable academic legitimization of this format and my scholarly self, a tool for creating alibi of play and funding opportunities;thirdly, a wanton dream of such research being picked up by those rare colleagues that find themselves also at the identity }\DIFdelend \DIFaddbegin \DIFadd{and a scholarly articulation that makes the project legible to others who might join it---those at the uncommon }\DIFaddend intersection of debate heritage, game-design \DIFdelbegin \DIFdel{career and scholarly ambitions}\DIFdelend \DIFaddbegin \DIFadd{practice, and scholarly ambition}\DIFaddend .

\DIFdelbegin \DIFdel{And with that the drums of ambition beat to the march of dreams to persevere in research galore, to fund and discover how to construct a financially sustainable organism }\DIFdelend \DIFaddbegin \DIFadd{The design knowledge the thesis generates may also travel beyond debate. What }\emph{\DIFadd{Inside Out}} \DIFadd{explores bears on any context where people reason together under structure. In education, the format's separation of an evaluated decision-maker from an external, growth-oriented evaluator suggests a model for classroom deliberation that develops perspective-taking without collapsing into a contest to be won. In mediation and facilitated group decision-making, its multi-party, synthesis-seeking structure offers a procedural template for surfacing and reconciling competing interests. These are speculations and exciting seeds for potential further design and stufy.

}

\DIFadd{My longer aim, beyond this thesis, is to build something financially sustainable }\DIFaddend that orchestrates and supports debate-centric educational gameplay. \DIFdelbegin \DIFdel{Of this march, the first steps I spoke of }\DIFdelend \DIFaddbegin \DIFadd{The concrete path toward it runs through European funding mechanisms such as Erasmus+ and Horizon Europe grants, together with private investors operating in the education sphere, pursued through consortia that bring together Critique Gaming (the game-development company through which I have worked in this space for many years), debate-world educational institutions such as IDEA and ARDOR, academic institutions and potentially governmental educational partners. The goal is a long-term, properly resourced programme in which game designers, scholars, NGO executives, and debate specialists collaborate to build an extended circuit of debate events and recurrent debate-based educational processes, anchored in the existing infrastructure of debate clubs and competitions, with a pipeline for rigorous quantitative and qualitative data collection, active iterative development of debate formats, and systematic measurement of transformative impact. That programme would also constitute the platform for a doctoral thesis capable of doing what this one cannot: sustained longitudinal study, robust mixed-methods data, and a design process iterated at full scale. The first steps toward it were named }\DIFaddend at the end of the discussion. \DIFdelbegin \DIFdel{May there be many more to come. }\DIFdelend \DIFaddbegin \DIFadd{As a pilot, the study described in this thesis could help justify and jump-start all the above. Alfred }\DIFaddend ``Tuna'' \DIFdelbegin \DIFdel{proved that one can fund a lifestyle }\DIFdelend \DIFaddbegin \DIFadd{Snider built a life }\DIFaddend of meaningful teaching\DIFdelbegin \DIFdel{helping }\DIFdelend \DIFaddbegin \DIFadd{, reaching }\DIFaddend hundreds of young people\DIFdelbegin \DIFdel{simply by writing academic papers on how }\DIFdelend \DIFaddbegin \DIFadd{, on the conviction that }\DIFaddend debate is a game \DIFaddbegin \DIFadd{worth taking seriously as one. I pray to be but even half as graced}\DIFaddend . 
\DIFdelbegin \DIFdel{One can hope this lightning strikes twice.
}\DIFdelend %%%%%%%%%%%%%%%%%%%%%%%%%%%%%%%%%%%%%%%%%
% ANNEXES — hub file
% Each annex group is maintained in its own file:
%   annex_format.tex    — Annexes A, B (format design iterations 1 & 2)
%   annex_materials.tex — Annexes C, D, E (debrief form, playtest motion, consent)
%   annex_reports.tex   — Annexes F, G, H (playtest reports + design iteration 3)
%   annex_clues.tex     — Annex I (clues from playtests)
%%%%%%%%%%%%%%%%%%%%%%%%%%%%%%%%%%%%%%%%%

%%%%%%%%%%%%%%%%%%%%%%%%%%%%%%%%%%%%%%%%%
% ANNEXES
%%%%%%%%%%%%%%%%%%%%%%%%%%%%%%%%%%%%%%%%%

\newpage
\section*{Annex A -- Inside Out Debate Format: First Design Iteration}
\addcontentsline{toc}{section}{Annex A -- Inside Out Debate Format: First Design Iteration}

The \textbf{Inside Out} format is an experimental debate format, designed using best practices from tabletop RPG game design. It seeks to be a more \textit{fun}, \textit{performance-growth oriented}, and \textit{real-life skill-building} alternative to traditional competitive debate formats. Rather than winning arguments, debaters explore values, make meaningful decisions, and earn recognition for the skills they demonstrate.

\subsubsection{Team Structure}

Debaters are organised in four teams of two debaters each and are assigned positions within the debate at random:

\begin{itemize}
    \item \textbf{Team S (Subject):} Represents the Subject of the Motion, the person who must face a decision. Must weigh all values and choose.
    \item \textbf{Teams A, B, and C (Value Defenders):} Each argues that a particular Value held by the Subject should be prioritised in the decision, and consequently, what the Subject should decide.
\end{itemize}

\subsubsection{The Motion}

The Motion outlines:

\begin{itemize}
    \item What is the decision that has to be made.
    \item What is the background of the Subject and other relevant practical factors that can be taken into account within the decision.
    \item What are the three Values of the Subject, each assigned to one Value Defender team.
\end{itemize}

It is important to note that the Motion does not prescribe to the Value Defenders a particular position on what decision the Subject has to make. They have to construct it themselves.

\subsubsection{Speech Order}

The four teams hold alternating speeches in the following order:

\begin{center}
A1, B1, C1, S1, A2, B2, C2, S2
\end{center}

Within their speeches, the Value Defenders argue for the relevance of their value, construct a specific decision that best honours it, and engage with the arguments forwarded by the other Value Defenders.

Within their speeches, the Subject team must outline a strong reasoning that engages with and explains the interplay between all arguments of the Value Defenders, picking a particular decision. Within S1 they pick a preliminary decision and within S2 a final decision.

\subsubsection{Evaluation}

The debate is overseen by an Evaluator, who performs two tasks:

\paragraph{Commendations.} The Evaluator awards Commendations to debaters. Commendations represent acknowledgements of useful skills, knowledges, or attitudes. A list of awardable commendations is constructed by the Chief Evaluators of the Debate Tournament, along with the total number of Commendations that Evaluators can offer within that Tournament. Example commendations include Believable Passion, Witty Solution, Argument Ingenuity, and Resilience (see Section~9 of this Annex for the full list).

Over time, debaters who gather a particular commendation at least three times earn a Badge. Gathering it additional times upgrades the badge, an unlimited number of times, but each upgrade requires gathering three more Commendations than last time.

\paragraph{Oral Feedback.} The Evaluator delivers an Oral Feedback session, in which they explain to each team the commendations they received, the reasoning behind them, and give advice on how each debater can improve in the future based on their performance.

\subsubsection{Example Motion}

\begin{quote}
You are the parent of a 16-year-old teenager. Your family is middle class and lives in the suburbia of a metropolis. You have higher education. You live in a country in which Higher Education is free, but there are only a few universities that are deemed as good. The year is 2026. You evaluate that your kid is talented in music, is a good leader and has had very good STEM skills, but is also quite distracted, does not have many friends and has failed to learn a foreign language. Your kid expressed a moderate interest in studying Musicology at university after they finish high school, but is also considering not going to higher education. You have a good relationship with your child. They asked you to suggest what they should do. What do you recommend your kid they should do after they finish high school?

\medskip

Three Beliefs that you hold, relevant for this decision, are:
\begin{itemize}
    \item In life, financial stability and independence is essential for a happy life.
    \item It is important to be a well-rounded, cultured individual.
    \item The people that you surround yourself with end up defining who you are.
\end{itemize}
\end{quote}

\subsubsection{Comparison with Traditional Formats}

The table below contrasts Inside Out with World Schools debate along six dimensions.

\begin{center}
\begin{tabular}{|l|p{5cm}|p{5cm}|}
\hline
\textbf{} & \textbf{World Schools} & \textbf{Inside Out} \\
\hline
Teams & 2 teams $\times$ 3 members & \textbf{4 teams $\times$ 2 members} \\
\hline
Sides & Proposition vs.\ Opposition & \textbf{3 Value Defenders + 1 Subject} \\
\hline
Motion style & ``This House\ldots'' policy/value motions & \textbf{Personal dilemma with background \& values} \\
\hline
Goal & Win the argument: defeat the other side & \textbf{Explore values: help the Subject decide} \\
\hline
Evaluation & Scored on Content, Style \& Strategy (points) & \textbf{Commendations for skills \& attitudes (badges)} \\
\hline
Outcome & Win / Loss + speaker rankings & \textbf{No win/loss: growth-oriented feedback} \\
\hline
\end{tabular}
\end{center}

\medskip

\noindent\textbf{Key shift:} Inside Out moves from adversarial persuasion to collaborative deliberation. There is no winner; instead, debaters earn commendations that recognise individual skills, and the Subject synthesises all perspectives into a reasoned decision.

\subsubsection{Player Instructions}

\paragraph{Value Defenders.}
\begin{itemize}
    \item Formulate arguments for why your value is \textbf{important, useful, and relevant} to this dilemma.
    \item Propose a \textbf{specific decision} that best honours your value; argue why it is right and leads to good outcomes, especially through the lens of your value.
    \item Engage with other Defenders' \textbf{values}: show yours is more important, more relevant to the dilemma, or safer to follow.
    \item Engage with other Defenders' \textbf{decisions}: attack, support (claiming it maximises your value), or endorse with addendums. Choose your positioning wisely.
    \item Your core mission: \textbf{define what honouring your value looks like} in practice, and persuade the Subject it is the right path.
\end{itemize}

\paragraph{The Subject.}
\begin{itemize}
    \item \textbf{Enter without an opinion.} You are completely torn: you are here to listen and be persuaded, not to argue.
    \item Listen as if you \textit{are} the Subject: \textbf{let the speeches move you} as they would move this person.
    \item Also take a mental step back: \textbf{evaluate how persuasive each argument would be} for the Subject, and adjust your lean accordingly.
    \item Take notes on arguments and their interplay: record what the Defenders say, \textbf{not your own personal thoughts}.
    \item \textbf{Collaborate quietly} with your partner: you are both embodying the same person, not competing.
    \item In your speeches, share which decision you are leaning toward and \textbf{why}, weighing all arguments to the extent they were persuasive. You must pick a specific decision, not merely name a value.
\end{itemize}

\paragraph{The Evaluator.}
\begin{itemize}
    \item Analyse the debate and each player's performance. Decide which \textbf{commendations} to award: aim for one per debater (the same commendation can go to more than one person).
    \item Approach the round as an \textbf{educator} helping each player grow, not as a judge declaring a winner.
    \item Open with a \textbf{narrative of the debate} from your outside perspective: what each team did and how it all unfolded.
    \item Engage directly with each team: ask what they \textbf{felt they did well} and what they found difficult. Share your view and add your own observations.
    \item Ask each team what they want to \textbf{do better next time}, and give concrete advice on how to get there.
    \item Close by \textbf{awarding commendations} to each player.
\end{itemize}

\subsubsection{Rituals}

\paragraph{Rolling In.}
\begin{enumerate}
    \item \textbf{Team commitment chant.} Each team member states their commitment aloud before the debate begins:

    Subject: \textit{``I am [character name]. I am open to being moved.''}

    Each Defender: \textit{``I believe in [value]. I will argue from that belief.''}
\end{enumerate}

\paragraph{Rolling Out.}
\begin{enumerate}
    \item \textbf{Subject's thanks.} The Subject team thanks each Defender individually: \textit{``Thank you for speaking your truth.''}
    \item \textbf{Defenders' thanks.} Each Defender team thanks the Subject: \textit{``Thank you for listening and carrying the decision.''}
    \item \textbf{De-roling.} Everyone stands. Each person states their real name and shakes hands with every participant: \textit{``I am [name]. Thank you for playing with me.''}
    \item \textbf{Bleed check.} The Evaluator asks quietly: \textit{``Did anything in that debate feel close to home for anyone?''} (a brief, optional moment to name any emotional residue before moving into feedback). No one is required to answer.
\end{enumerate}

\subsubsection{Full Commendation List}

Commendations are organised across five categories. Each commendation has three award tiers: \textit{Earned} (threshold quality), \textit{Strong} (clear excellence), and \textit{Exceptional} (reserved for outstanding, rare performances). Some commendations carry role restrictions (Defenders only, or Subject only) as noted.

\medskip
\noindent\textbf{Style} \textit{(6 commendations: delivery, presence, and rhetorical quality)}

\paragraph{Believable Passion.} You were truly invested in your point, showing genuine, believable conviction and care for the value you defended. Distinct from Honest Connection (self-disclosure) and Broad Range (technical delivery control).
\begin{description}
    \item[Earned] Tone conveyed genuine care at key moments; voice quality shifted naturally, eye contact held at key claims, word choices felt spontaneous.
    \item[Strong] Conviction was felt in the substance of the speech; the room or the Subject visibly reacted.
    \item[Exceptional] The delivery made the dilemma feel real and worth deeply caring about. Reserve for speeches where the affective dimension was genuinely transformative.
\end{description}

\paragraph{Entertaining Discourse.} The speech was genuinely enjoyable through wit, lightness, or a well-timed observation that made the room smile or lean in. Distinct from Poetic Appeal (literary craft).
\begin{description}
    \item[Earned] At least one moment was distinctly enjoyable: a clever comparison, a light touch, or a well-timed observation.
    \item[Strong] Liveliness was woven into the argumentation consistently; humor or vividness served to make points land more effectively.
    \item[Exceptional] Entertainment and argument were seamlessly fused throughout, reaching the level of a sharp public talk where wit and substance were inseparable.
\end{description}

\paragraph{Broad Range.} Your delivery showed impressive stylistic range, varying tone, rhythm, tempo, and volume to serve the speech. Distinct from Believable Passion (authenticity of feeling); rewards technical delivery control.
\begin{description}
    \item[Earned] At least two noticeably different delivery registers were used; the speech did not feel flat.
    \item[Strong] Variation was deliberate and purposeful: shifts in pace or tone coincided with structural moments in the argument.
    \item[Exceptional] The full palette was deployed with mastery; the debater controlled the room's attention through delivery alone.
\end{description}

\paragraph{Honest Connection.} Something genuine about yourself entered the speech, a real disclosure rather than a rhetorical move, that created contact between the debater and the room. Distinct from Believable Passion (conviction); rewards self-disclosure and personal presence.
\begin{description}
    \item[Earned] At least one moment of genuine self-disclosure (an admission, first-person reflection, or personal statement of feeling) that felt real rather than strategic.
    \item[Strong] The personal register appeared at key structural moments; eye contact moved around the room, direct second-person address, or deliberate pause.
    \item[Exceptional] The speech was uniquely theirs; you felt the debater was genuinely present in the room with you.
\end{description}

\paragraph{Poetic Appeal.} Language was used with aesthetic care, where rhythm, imagery, or metaphor elevated the speech as a piece of writing, independent of its entertainment value. Distinct from Entertaining Discourse (enjoyment response); rewards literary and rhetorical craft.
\begin{description}
    \item[Earned] At least one passage was noticeably well-crafted: a sentence that would stand alone as good writing, or a rhythmic construction that made a point feel inevitable.
    \item[Strong] Poetic quality sustained across multiple moments; rhetoric served the argument.
    \item[Exceptional] The speech had genuine literary quality; specific lines are still in your head afterward. Language itself became part of the argument.
\end{description}

\paragraph{True Embodiment.} You convincingly inhabited your role, arguing and deciding from inside it rather than about it from outside. For the Subject: thinking, feeling, and deciding as the character. For Defenders: arguing from genuine conviction in the value.
\begin{description}
    \item[Earned] The role's reality guided the speech throughout; stated leanings tracked the actual arguments made.
    \item[Strong] The embodiment elevated the entire round; you could see the speaker being moved, reconsidering, weighing.
    \item[Exceptional] Not applicable; omit Exceptional tier for this commendation.
\end{description}

\medskip
\noindent\textbf{Strategy} \textit{(7 commendations: positioning, framing, and tactical intelligence)}

\paragraph{Specific Topicality.} Arguments were sharply specific to this dilemma, this Subject, and these values, with no generic claims. Distinct from Advantageous Framing (evaluative lens); rewards the content of the arguments themselves.
\begin{description}
    \item[Earned] Most substantive claims grounded in the specifics of the motion: the Subject's background, the values as stated, the actual decision at stake.
    \item[Strong] Topicality was strategic; the debater drew on details others missed.
    \item[Exceptional] Every substantive move calibrated to this specific situation; the arguments would fall apart if the Subject or dilemma were changed.
\end{description}

\paragraph{Advantageous Framing.} You defined the terms of discussion in a way that benefited your position before others could. Distinct from Specific Topicality (argument content) and Helpful Guiding (neutral Subject navigation).
\begin{description}
    \item[Earned] Explicitly established a coherent evaluation framework relevant to their value's claim.
    \item[Strong] The framing did real structural work; other debaters had to address it, and the Subject engaged with it directly.
    \item[Exceptional] The framing was original and visibly influenced the course of the debate. Award rarely.
\end{description}

\paragraph{Witty Solution.} You proposed a concrete, value-informed solution to the dilemma, specific enough that it could only have come from someone defending your value. Distinct from Advantageous Framing (criteria); rewards the solution itself.
\begin{description}
    \item[Earned] The proposed solution was specific and thought through, demonstrably flowing from the defended value.
    \item[Strong] Notably smart: navigated a genuine tension between values, or revealed an angle others missed.
    \item[Exceptional] Elegant and surprising, more persuasive than alternatives precisely because of how the value was applied.
\end{description}

\paragraph{Brave Pivot.} You changed position or conceded a point honestly, and it made your case stronger.
\begin{description}
    \item[Earned] Explicitly acknowledged an opponent's point had merit, or adjusted position in response to what was said.
    \item[Strong] The pivot was strategic: conceding less important ground to better defend more important ground.
    \item[Exceptional] Changed their mind in real time under pressure, arguing why the pivot extends rather than contradicts the previous position. Award with corresponding rarity.
\end{description}

\paragraph{Efficient Alliance.} \textit{(Defenders only)} You recognized when another Defender's arguments complemented yours and built on that alignment.
\begin{description}
    \item[Earned] Acknowledged at least one point of genuine agreement with another Defender that strengthens the case for a decision.
    \item[Strong] Explicitly constructed an alliance: ``both Team A and I agree that X, which shows the Subject that the weight of argument points toward Y.''
\end{description}

\paragraph{Careful Positioning.} \textit{(Defenders only)} Your stance toward other Defenders' decisions was deliberate and chosen rather than reactive, with positive effect.
\begin{description}
    \item[Earned] Response to other Defenders' decisions was evidently chosen rather than accidental; the stance was coherently maintained.
    \item[Strong] Positioning was in service of a coherent strategy; you could reconstruct the reasoning behind every stance taken, forming a consistent picture.
\end{description}

\paragraph{Helpful Guiding.} \textit{(Subject only)} You helped the room navigate the structure of the debate clearly, accurately, and fairly. Distinct from Advantageous Framing; the Subject's role is neutral navigation, not partisan framing.
\begin{description}
    \item[Earned] At least one moment of genuine meta-level navigation, e.g.\ ``the key question before us is now X, because Y has been established and Z has been challenged.''
    \item[Strong] The guiding was structurally important, identifying a genuine turning point or consolidating arguments fairly to all sides.
\end{description}

\medskip
\noindent\textbf{Content} \textit{(9 commendations: argumentation, analysis, and engagement with the debate)}

\paragraph{Word Economy.} You communicated clearly and efficiently; every minute of speaking time counted. Rewards form of expression, not argument substance.
\begin{description}
    \item[Earned] The speech did not waste significant time; arguments were stated without excessive preamble, redundant restatement, or filler.
    \item[Strong] Concision actively served clarity: complex ideas were expressed with precision; understood fully on first hearing.
    \item[Exceptional] The debater said something hard to say in few words; specific lines had a density of meaning that rewarded attention.
\end{description}

\paragraph{Rebuttal Ingenuity.} Your rebuttal identified a non-obvious flaw in an opponent's argument: a hidden assumption, internal tension, or unconsidered consequence. Distinct from Argument Ingenuity (novel constructive arguments) and Comprehensive Breadth (coverage).
\begin{description}
    \item[Earned] Identified a non-obvious weakness rather than the first available attack.
    \item[Strong] Engaged with the strongest version of the argument and demonstrated why it fails even there.
    \item[Exceptional] Changed the argumentative landscape; the rebuttal was difficult or impossible to satisfactorily respond to.
\end{description}

\paragraph{Rebuttal Prioritization.} You made deliberate, strategic choices about where to concentrate your rebuttal effort, and those choices were correct. Distinct from Comprehensive Breadth (coverage); rewards deliberate focus over breadth.
\begin{description}
    \item[Earned] Concentrated force on the most important clashes, explicitly or implicitly flagging the allocation.
    \item[Strong] Prioritization revealed genuine strategic intelligence; the choice of what to ignore was as deliberate as what was attacked.
\end{description}

\paragraph{Rebuttal Accuracy.} You represented opponents' arguments faithfully before responding (no mischaracterization, no straw-manning). No Exceptional tier; accuracy is a threshold quality. Award only when faithful representation was genuinely non-trivial.
\begin{description}
    \item[Earned] Paraphrased opponents' arguments correctly; engaged with the most charitable understanding.
    \item[Strong] Engaged with specific moments in specific speeches; accuracy made the rebuttal more effective, not just fairer.
\end{description}

\paragraph{Robust Logic.} Your arguments were well-constructed: valid, plausible, and without significant logical gaps. Distinct from Impressive Argumentative Depth; logic asks whether the reasoning chain holds structurally.
\begin{description}
    \item[Earned] Main arguments had clear claim-warrant-impact structure; a reasonable person could follow without supplying missing steps.
    \item[Strong] Logic held up under challenge: no internal contradictions, no steps requiring the audience to take things on faith.
    \item[Exceptional] Genuinely tight; would require a sophisticated objection to dislodge; obvious objections anticipated within the argument itself.
\end{description}

\paragraph{Impressive Argumentative Depth.} You went beyond the surface, interrogating implications, mechanisms, or philosophical roots. Distinct from Robust Logic (structural soundness) and Argument Ingenuity (novelty).
\begin{description}
    \item[Earned] At least one argument went meaningfully beyond its initial claim, explaining why the value's demand is well-founded or what it would look like for the Subject to live by it.
    \item[Strong] Depth sustained throughout; consistently pushed past surface-level argumentation.
    \item[Exceptional] After the speech, you found yourself reconsidering something you had assumed you understood, or the argument illuminated a dimension of the value no other speech had touched.
\end{description}

\paragraph{Comprehensive Breadth.} You engaged with the full scope of the debate, including inconvenient arguments. Distinct from Rebuttal Ingenuity (sharpness) and Rebuttal Prioritization (deliberate focus); rewards attentiveness across the full debate.
\begin{description}
    \item[Earned] Addressed the most important arguments from all other Defenders, not only the weakest or most convenient.
    \item[Strong] Demonstrated careful listening throughout; the speech reflected the actual state of the debate.
    \item[Exceptional] Managed to synthetically yet explicitly cover every significant aspect of the debate.
\end{description}

\paragraph{Argument Ingenuity.} You produced an argument that was unexpected, non-obvious, and substantively illuminating. Distinct from Rebuttal Ingenuity (reactive) and Impressive Argumentative Depth (thorough development); rewards novelty and surprise.
\begin{description}
    \item[Earned] At least one argument was notably fresh: not typically heard in this kind of debate, logically coherent, and able to withstand an obvious first objection.
    \item[Strong] Ingenious in the full sense: first surprised you, then convinced you, revealing something about the value or dilemma that prior thinking had missed.
    \item[Exceptional] The argument changed the debate; others had to respond to it and it became a fixed reference point. Very rare, award accordingly.
\end{description}

\paragraph{Accurate Synthesis.} \textit{(Subject only)} You gave a fair, accurate, and well-organized account of the full state of the debate. Distinct from Comprehensive Breadth (Defenders' engagement); rewards the Subject for faithfully describing and weighing what happened.
\begin{description}
    \item[Earned] Correctly identified the main arguments from each Defender and did not misrepresent or omit them, even where unfavorable.
    \item[Strong] Showed genuine comprehension of the debate's structure: an account of how arguments interacted, where they clashed, and what the real decision points were.
    \item[Exceptional] Identified the strongest and weakest versions of each Defender's case and grounded the final decision in a genuinely comprehensive account of the round.
\end{description}

\medskip
\noindent\textbf{Fairplay} \textit{(3 commendations: observe team dynamics \& debrief, not only speeches)}

\paragraph{Teammate Support.} You actively helped your partner formulate something difficult in a way that visibly contributed to the team's performance. Look for: visible collaboration during prep time; a second speaker building clearly on something the first struggled with. Award when the help was substantive and its impact visible. Ask in the debrief if not fully observable.

\paragraph{Teammate Solidarity.} You showed genuine support for your partner during a hard moment, without distancing yourself. Look for: body language and verbal behavior toward a partner after a difficult moment; a second speaker who picked up where a struggling first left off without implicitly throwing them under the bus. The difficult moment must have been real and the solidarity genuine. If ambiguous, ask in the debrief.

\paragraph{Strategic Problem-Solving.} You demonstrated intelligent, adaptive problem-solving in the face of strategic challenges in the round. Look for: a visible strategy adjustment that worked; a debater who recognized a weakness in their own case and addressed it explicitly; coordinated strategic adaptation across both speeches. Both identifying and solving the problem must be visible and shared between the two team members.

\medskip
\noindent\textbf{Growth Mentality} \textit{(4 commendations: use the debrief; always name the specific moment in feedback)}

\paragraph{Applied Feedback.} You visibly tried to apply feedback from a previous round; imperfect application still earns this. Evidence may come from direct comparison with a prior round or from the debater's own account in the debrief. What matters is the genuine attempt rather than the outcome.

\paragraph{Out of Comfort Zone.} You attempted something genuinely difficult for you and made a real go of it. Rewards the attempt, not the result. Comfort zone is individual; use the debrief to identify what was genuinely challenging. Ask directly if needed.

\paragraph{Tracking.} You showed clear awareness of your own growth across debates: where you started, what you have worked on, and where you are heading. Distinct from Applied Feedback (applying specific feedback in one round); rewards the broader self-awareness of one's development arc over time. Evidence may come from the debrief: the debater can articulate what they have improved, name skills still being worked on, and describe their trajectory with accuracy.

\paragraph{Resilience.} You encountered a genuinely difficult moment and came out intact, or stronger. The difficult moment must have been real. Award for: a strong attack responded to with a composed second speech; a delivery stumble overcome mid-speech; an unexpected move handled with adaptability. Name the specific moment in feedback; this commendation must be concrete to be useful.

%%%%%%%%%%%%%%%%%%%%%%%%%%%%%%%%%%%%%%%%%

\medskip

\noindent\textbf{Note:} The most current design iteration and additional resources can be found at \url{https://olaruandreidan.github.io/InsideOutDebateFormat/}.

%%%%%%%%%%%%%%%%%%%%%%%%%%%%%%%%%%%%%%%%%

\DIFaddbegin 

\subsubsection{\DIFadd{Published Digital Artifact}}

\DIFadd{The following screenshots document the format as published on the project website (}\url{https://olaruandreidan.github.io/InsideOutDebateFormat/}\DIFadd{), representing the final design state after both playtest cycles.

}

\begin{figure}[H]
    \centering
    \includegraphics[width=\textwidth]{{otherdocuments/Screenshot 2026-06-10 at 20.46.16.png}}
    \caption{\DIFaddFL{Player instructions page: Inner Voices, the Subject, and the Evaluator.}}
    \label{fig:website-player-instructions}
\end{figure}

\begin{figure}[H]
    \centering
    \includegraphics[width=\textwidth]{{otherdocuments/Screenshot 2026-06-10 at 20.46.34.png}}
    \caption{\DIFaddFL{Rituals section and full commendation list as published.}}
    \label{fig:website-rituals-commendations}
\end{figure}

\DIFaddend 

\newpage
\section*{Annex B -- Inside Out Debate Format: Second Design Iteration}
\addcontentsline{toc}{section}{Annex B -- Inside Out Debate Format: Second Design Iteration}

The second design iteration is based on observations and participant feedback gathered during the first playtest (Jean Monet debate club, Bucharest, March 2026). See Annex F for a detailed first playtest report. The changes below supersede the corresponding elements in the first iteration (Annex A).

\subsubsection{Speech Order Changes}

The interstitial S1 speech between the first and second rounds is removed. In its place, two new elements are introduced:

\paragraph{S0 -- Subject Opening Speech.} A new opening speech before any Defender speaks. The Subject team introduces the dilemma and argues why it is genuinely hard: all three values are relevant, the decision is real, and there is no obvious right answer. This speech gives the debate proper gravitas before Defenders begin.

\paragraph{The Huddle.} A 10-minute interval inserted after C1 (and before A2). The Subject team may ask direct questions to any Defender team, grant them the floor to answer, and revoke their right to speak at any point. The Huddle allows the Subject to control the conversation toward resolving argumentative tensions, forcing arguments that need to interact to interact.

\paragraph{Thinking Time.} A 2-minute silent thinking period is added before S2, giving the Subject team time to integrate all arguments before delivering their final decision.

The new speech order is:

\begin{center}
S0, A1, B1, C1, [The Huddle], A2, B2, C2, [2-min thinking time], S2
\end{center}

\textbf{Points of Information (POIs)} are introduced during all speeches: any debater may signal to the current speaker to ask a brief question.

\subsubsection{Motion Changes}

\begin{itemize}
    \item The Subject character description in the motion should include a longer list of the Subject's personal beliefs and past decisions, to flesh out the character more fully and enable value-specific argumentation.
    \item Value Defenders should address the Subject in the second person (``you, as the Subject, should\ldots'') rather than in the third person, reinforcing the personal and direct character of the deliberation.
\end{itemize}

\subsubsection{Evaluation Changes}

Commendations are awarded in private: after the debate, the Evaluator approaches each team separately rather than delivering commendations in a public session. The oral feedback narrative and growth-oriented debrief remain public.

\subsubsection{Speech Instruction Clarifications}

\begin{itemize}
    \item Subject instructions explicitly state that the Subject must pick a specific \textit{decision} (a concrete course of action), not merely name a value they find most important.
    \item More explicit guidance is provided on what the Subject team should prepare during prep time: for S0, they should prepare to explain the dilemma's genuine difficulty and why all values are relevant; for S2, they should prepare to synthesise all arguments and justify a final specific decision.
\end{itemize}

%%%%%%%%%%%%%%%%%%%%%%%%%%%%%%%%%%%%%%%%%

\newpage

\section*{Annex C -- Debrief Form}
\addcontentsline{toc}{section}{Annex C -- Debrief Form}

The following questionnaire was administered to participants after the first playtest. Questions aim to gather data on entertainment value, format distinctiveness, participant intentions, and transformative effects.

\begin{enumerate}
    \item How fun was the debate? (scale 1--10)
    \item How fun was it compared to a debate in a traditional format? (5-choice scale)
    \item Was there anything surprisingly entertaining or engaging about it?
    \item Was there anything about the format that was quite difficult, to the point of frustration?
    \item Was there anything about the rules of the format that was confusing?
    \item Would you play this format some more? (Yes / No)
    \item If you would play it some more, in what circumstance would you?
    \item If there was to exist an alternative debate circuit centred around this format, would you attend those competitions? Even at the expense of competitions in the WS or BP format?
    \item Do you believe you played differently than you would have in a debate in WS/BP with a rather similar motion? If so, how?
    \item Do you feel different than how you usually feel after a debate? If so, how?
    \item Did you perceive yourself and others speaking similarly or differently than how you do in a regular debate? If so, how?
    \item Were you to redo this debate, what would you do differently?
    \item Were you to play more in this format, but on different motions, what would you do differently?
    \item What do you believe you learned from this debate?
    \item What do you believe this debate format does better, as a learning/empowerment tool, than WS or BP? Do you believe it does something worse?
\end{enumerate}

%%%%%%%%%%%%%%%%%%%%%%%%%%%%%%%%%%%%%%%%%

\newpage
\section*{Annex D -- Playtest Motion}
\addcontentsline{toc}{section}{Annex D -- Playtest Motion}

The following motion was used in the first playtest at the Jean Monet debate club, Bucharest, March 2026.

\paragraph{First version.} Used in the first playtest at the Jean Monet debate club, Bucharest, March 2026.

\begin{quote}
You are Alex. You are a lawyer. You are close to your older brother Matei. He is a very caring person to everyone, always eager to help even strangers. He always helped you when you were in need. However he was always quite dishonest with you, and you caught him lying even about important things. Matei is also an alcoholic and you believe he is quite bigoted.

\medskip

Matei was involved in a hit and run recently. He hit a cyclist that ended up hospitalised. The cyclist remembered his license plate only partially and he is now being investigated by the police. From your knowledge they only have circumstantial evidence, but he still risks going to jail. He asked you to lie that he was with you at home at the time of the accident, in order to provide him an alibi. Your assessment as a lawyer is that this will almost certainly exonerate him. However, a very small risk remains that additional material evidence could be found, at which point you might become an accessory to the crime.

\medskip

\textbf{What should you do?}

\medskip

Three Values that you hold, relevant for this decision, are:
\begin{itemize}
    \item You believe that helping people in need is very important, and a moral duty.
    \item You believe in justice: that people should receive what they deserve.
    \item You believe that family should stick together and have each other's back.
\end{itemize}
\end{quote}

\paragraph{Second version.} Used in the second playtest, April 2026.

\begin{quote}
Alex is a criminal defence lawyer with seven years of experience. The eldest of three siblings from a small town, Alex was the first in the family to attend law school. After graduating, Alex spent two years at a legal aid clinic representing clients who could not afford private counsel, before joining a private firm. Alex's first major case there was defending a driver charged with causing serious injury in a road accident; Alex secured his acquittal.

\medskip

Matei is Alex's older brother and the two have remained close throughout their lives. Matei is known among everyone who knows him for being eager to help others, often going out of his way even for strangers. He helped Alex through several difficult periods. He has also been dishonest with Alex repeatedly over the years, including about important matters. Matei struggles with alcohol and holds a number of views that Alex finds bigoted.

\medskip

Matei was recently involved in a hit-and-run: he struck a cyclist who was hospitalised. The cyclist recalled his license plate only partially, and Matei is now under police investigation. As far as Alex knows, the evidence against him is circumstantial. Matei has asked Alex to lie and say he was at home with Alex at the time of the accident. Alex's legal assessment is that this would almost certainly exonerate him, though a small risk remains that additional material evidence could surface, at which point Alex would become an accessory to the crime. \textbf{What should you do?}

\medskip

Values:
\begin{enumerate}
    \item You believe that helping people in need is very important, and a moral duty.
    \item You believe in justice, that people should receive what they deserve.
    \item You believe that family should stick together and have each other's back.
\end{enumerate}
\end{quote}

%%%%%%%%%%%%%%%%%%%%%%%%%%%%%%%%%%%%%%%%%

\newpage
\section*{Annex E -- Informed Consent Form}
\addcontentsline{toc}{section}{Annex E -- Informed Consent Form}

\bigskip

I, \underline{\hspace{7cm}} of \underline{\hspace{1.5cm}} years of age, voluntarily accept to participate in a playtest followed by an interview session, both moderated by Andrei Dan Olaru, as part of his \textit{``Applying Transformative Game Design Theory to Competitive Debate''} study. This study seeks to explore the creation of an alternative debate format, informed by modern transformative game design theory. I understand that:

\begin{enumerate}
    \item This session implies video recording of the entire debate playtest and the subsequent collective interview, plus some additional questions using an online questionnaire.
    \item Pseudonyms will be used when analysing the data for confidentiality, and the researchers commit to not revealing the identities of the participants. Individual raw data may be discussed in class or with the supervisor without the participant's name to trace back to. The data obtained will be analysed by grouping and some textual examples will be used anonymously.
    \item The recording will be safely stored in the university's OneDrive to maintain the saved raw interviews protected. The video will be erased one month after the dissertation's results. The transcripts will be erased up to one year after the dissertation's results.
    \item None of the questions will go against my integrity, no harm will come to me during this session, no answer is right or wrong and I may choose not to answer any question(s) without consequences of any kind.
    \item I am free to walk out of the playtest and interview at any time if I deem it necessary without prejudice or any consequences.
\end{enumerate}

For any additional information or difficulty, the participant may contact the researcher Andrei Dan Olaru at \texttt{andrei-dan.olaru.0744@student.uu.se}.

\bigskip\bigskip

\noindent
\begin{tabular}{p{6cm} p{1cm} p{6cm}}
\underline{\hspace{6cm}} & & \underline{\hspace{6cm}} \\
Participant's signature & & Researcher's signature \\
\end{tabular}

\bigskip\bigskip

I, \underline{\hspace{7cm}} have been informed of the conditions and freely accept to participate in the research by Andrei Dan Olaru.

\bigskip

\noindent
\begin{tabular}{p{6cm} p{5cm}}
\underline{\hspace{6cm}} & \underline{\hspace{3cm}} of March 2026 \\
Participant's signature & Date \\
\end{tabular}

%%%%%%%%%%%%%%%%%%%%%%%%%%%%%%%%%%%%%%%%%

\newpage

\section*{Annex F -- First Playtest Report}
\addcontentsline{toc}{section}{Annex F -- First Playtest Report}

The following is a first-person account of the first playtest session, held in March 2026 at the Jean Monet debate club in Bucharest. The session gathered six senior debaters (aged 18--20). Design changes arising from this playtest are documented separately in Annex B

\subsubsection{Setting and Participants}

The playtest involved six senior debaters from the Jean Monet debate club in Bucharest, all aged between 18 and 20. These were experienced debaters, though none had encountered the Inside Out format before. I facilitated the session as the Evaluator and observed throughout.

\subsubsection{The Brief}

The brief took approximately 40 minutes. The format was significantly different enough from traditional debate that this length of explanation was necessary even for experienced players. Participants had many questions, particularly about two elements: (1) how the Subject role works and the roleplay dimension, which was described as ``the newest and most different part of the format'', and (2) whether counter-argumentation between Value Defenders is permitted, a question I found somewhat surprising in hindsight.

\subsubsection{Preparation Time}

During prep time, the players struggled to figure out what to prepare. This was the first indication that the format demands a fundamentally different kind of preparation than traditional debate. Unlike preparing a policy or value case for traditional formats, players needed to figure out how to approach the Subject's dilemma through the lens of a specific value.

\subsubsection{Rituals}

The roll-in rituals felt awkward and cringe to participants. The rituals themselves were lengthy and formalized, which created discomfort rather than appropriate atmosphere-setting. More effective would be a brief introduction by the Evaluator followed by a small gesture or shared chant that everyone performs together.

Notably, the players did recognize and appreciate that the rituals acknowledged the effort and role of the judge, in this case shared between the Evaluator and Subject.

The calling rituals (used to introduce each speaker) were received positively and did not require revision.

The roll-out rituals also felt overly lengthy and became cringe. However, one participant noted they appreciated the moment of collective gratitude that the rituals attempted to create. A simpler version involving handshakes and brief role-specific phrases would likely work better.

\subsubsection{Speech Flow and Embodiment}

An interesting pattern emerged across the speeches: players began using World Schools parliamentary plural and third-person references early on, but speech by speech, they naturally shifted toward second-person address and self-identification with the Subject's values. One player would say ``we believe'' at the start of their first speech, but by the second round was saying ``I believe in this value'' and referring to the Subject directly as ``you.''

This raised an important design question: should the Value Defenders fully embody the voice of the value (speaking as if they are an inner voice of the Subject themselves), or should they position themselves as advisors speaking about the Subject in the second person? The natural drift toward the former suggests players intuitively moved toward stronger roleplay embodiment, but the rules should be clarified on this point.

\subsubsection{First Round of Value Defenders}

Players immediately grasped that they needed to argue for the importance and relevance of their assigned value to the dilemma. They picked up on the fact that the solutions offered could be complex, realistic, and nuanced, distinct from traditional debate solutions.

However, the distinction between arguing for a value and proposing a decision remained fuzzy for many. The fact that they must ``pick a solution and not necessarily a value'' was confusing initially; the rules require clarification and possibly example mini-speeches illustrating this distinction.

\subsubsection{Argumentation Style}

The format demands a form of argumentation that is fundamentally different from traditional debate. Rather than constructing large-scale causal chains and impact scenarios, Value Defenders must argue for their values on their own terms: their importance, their weight in the decision, their interaction with other values. This was hard for players: it is uncommon in World Schools or BP debate, and finding ``more things to say'' about values themselves (as opposed to practical outcomes) required a different cognitive mode.

One significant observation: players felt the Subject character needed to be more fully fleshed out. The richer the characterization (the beliefs they hold, decisions they have made in the past, their personality) the easier it would be for Value Defenders to construct arguments tailored to that specific person's identity and narrative consistency. Players noted that this reflects how people actually construct their own identities: by looking at the coherence of their past actions and trying to remain narratively consistent.

\subsubsection{Participant Experience}

Despite the challenges, or perhaps because of them, participants reported the format to be surprisingly liberating. They reported feeling ``a lot less anxiety'' than in traditional debate, while also finding the work ``actually very complex and interesting strategically.'' They noted that architecting solutions was intellectually demanding in a different way: every detail of the context became potential argumentative material rather than merely setting for claims about large-scale consequences.

One highly competitive player, the type that ``usually wins a lot at competitions'', reported finding the format less exhilarating. They felt that the cooperative nature took away from the ``thrill of collegial competitive jabbing'' and the ``adrenaline of winning,'' making the overall experience feel ``anticlimactic.'' This was a useful signal: the format appears to function as an alternative to traditional debate rather than a replacement, appealing to different motivations than ego-driven competitive success.

\subsubsection{Recognized Format Distinctiveness}

Players articulated several ways in which the format plays differently:

\begin{itemize}
    \item Explaining long, mechanistic causal chains is less necessary because there is no need to think through large-scale consequences (though such thinking could still occur if the dilemma involved someone with power over others).
    \item The power to formulate solutions makes it easier to stay topical and makes the whole debate less vague: the debate becomes about real dilemmas rather than abstract policy positions.
    \item It is less about proving large-scale third-order impacts and more about navigating genuine tensions between values.
\end{itemize}

Players reported speaking more slowly but with similar theatricality to traditional debate, contrary to expectations. Though some initially felt they spoke more monotonously, other participants who reported understanding that the Subject could be emotionally moved as well noted they were quite theatrical and expressive. Some said that were they to debate again, they would try to be more emotionally immersed, consciously shifting away from trying to be very ``sweaty'' toward winning.

\subsubsection{Reported Learning and Growth}

When asked what they believed they had learned, participants identified several insights:

\begin{itemize}
    \item The format teaches you to be more relaxed and to enjoy playing.
    \item It teaches you to think about the context of decisions rather than abstractly applying rules.
    \item It makes you more human and more honest about human nature.
    \item It engages with the identity of decision-making: the recognition that ``every team says valuable stuff and it is important more rather to see what is important to us.''
    \item This last point led to a striking reflection: that sometimes decisions are not universal, they are personal, and arguing that someone is morally wrong when deciding or believing something is a bit absurd.
\end{itemize}

Some participants felt the format is more suitable for beginners than for highly competitive debaters. I agree with the latter assessment but also believe that with a more complex motion, experienced debaters could recognize how the format could be played at higher skill levels.

\subsubsection{Participant Feedback on Rules}

Many participants requested the addition of Points of Information (POIs): brief interjections during speeches. This became a high-priority feature request.

The participants also noted that the Subject role needs explicit scaffolding during preparation time. Rather than leaving this open, it would be helpful to have them prepare a first speech in which they explain the dilemma, why it is so hard, and how all values are relevant. Giving the dilemma ``real gravitas'', proving to everyone that the decision is truly important and truly hard, would set the stage more effectively.

\subsubsection{Observational Notes and Design Recommendations}

Several design improvements emerged from observation:

\begin{itemize}
    \item \textbf{Ritual for calling people:} A more natural, shorter ritual is needed between speeches to call in the next speaker.
    \item \textbf{Subject preparation:} Give the Subject team something concrete to do during prep time: prepare a first speech that explains the dilemma and its difficulty, establishing that all values are genuinely relevant.
    \item \textbf{Motion design:} Include a longer list of things the Subject believes about themselves and decisions they have made in the past. This supports more specific, character-tailored argumentation by the Value Defenders.
    \item \textbf{Speech instructions:} Make explicit instructions clearer, particularly on the distinction between defending a value and proposing a solution, and on whether Value Defenders should embody an ``inner voice'' of the Subject or position themselves as advisors.
    \item \textbf{Interstitial speech:} Consider whether an interstitial speech between the first and second rounds of Value Defenders might help reset the debate or provide needed structure. (This was later addressed in subsequent iterations by introducing the Huddle.)
\end{itemize}

%%%%%%%%%%%%%%%%%%%%%%%%%%%%%%%%%%%%%%%%%

\newpage
\section*{Annex G -- Second Playtest Report}
\addcontentsline{toc}{section}{Annex G -- Second Playtest Report}

The following is a first-person account of the second playtest session, held in April 2026. The session gathered seven veteran debaters. Design changes arising from this playtest are documented separately in Annex H.

\subsubsection{Setting}

I held the second playtest together with seven veteran debaters, close friends. We held it at the apartment that part of us commonly rent and share with other nerds for the single purpose of playing TTRPGs there. Because of that, the main room is dominated by a hilariously large, roughly round table. We held the debate seated at it, which one player noted created a council-feel to the whole session, as opposed to a regular debate in which the seating mimics parliamentary discourse and thus feels profoundly more impersonal. I was not the Evaluator this time; I played as the second speaker of the A Inner Voice team.

\subsubsection{The Brief}

The brief took about 40 minutes, similar to how much it took with the high-schoolers. The fact that it required that amount of time even with people accustomed to parsing lengthy rulebooks speaks to how much the design differs from traditional formats. The division of the ``judge'' role into two separate individuals (the Subject and the Evaluator) was noted as interesting by several players. The role of the Subject was by far the hardest for players to understand, especially the purpose of the opening Subject speech. I had, at the point of writing the rules, found it genuinely difficult to pin down what the Subject should do in that first speech. I wrote vague guidelines, and this showed, but in the end I believe it was for the better: I let the player decide what would be most interesting and useful. In that I was fortunate: the role was drawn by a player with substantial game-design experience, a professional L\&D background, many years as a trainer, and deep experience as a roleplayer.

\subsubsection{Rituals}

The roll-in rituals were experienced as cringe, and were surely too lengthy. Players felt they should be replaced by a short introduction from the Evaluator followed by a brief shared gesture or chant. This aligns with how regular debates open: the judge makes an introduction, then everyone bangs the table vigorously and the debate begins. One player noted that the rituals do recognise the effort put in by the judge, in this case shared between the Subject and the Evaluator.

The calling rituals were received positively.

The roll-out rituals were also too lengthy and became cringe, but one player noted they liked the way in which people express gratitude to everyone. The consensus was that this could be replaced by a moment of handshakes in which each person says a short phrase based on their role. Players also noted that calling a defender by name during the calling ritual felt unnecessary.

\subsubsection{S0 -- The Opening Subject Speech}

The player who drew the Subject role delivered an interesting, roleplay-heavy speech, even playing a soundtrack. They fleshed out the motion with many evocative details that further emotionally contextualised the scenario, then tried to explain how they care about all three values and why each is relevant to the decision. This in part exposed some nuanced personal biases that bled from the player into the character. Several players, including that player, noted that this institutionalises the already existing practice of debaters ``playing for the judge'': attempting to exploit the judge's known biases. I found this observation striking. The speech, although I was sceptical of it initially, does hold a great deal of value: fleshing out the scenario, giving a floor for surfacing biases. What it needs is clearer language for what players are supposed to do.

\subsubsection{First Round of Inner Voices}

All players immediately adopted a roleplay stance. They spoke directly to the Subject, addressing them by name and in the second person, looking at them, and choosing speech patterns that mirrored those the Subject had displayed. They also took the liberty to add small details to the motion, to worldbuild.

I find this very promising but also somewhat concerning. Giving players unofficial worldbuilding powers may lead inexperienced players to twist the debate in an unhelpful direction. However, adding colour to the broad facts of the scenario does lend to a more realistic simulation of how people actually argue, raises emotional engagement, and heightens immersion.

It was interesting that even very veteran debaters struggled, especially in this first round, to both appeal to emotion and to argue specifically for their assigned value. Most focused on constructing interesting and nuanced Solutions and arguing why these were practical and useful, utilitarian and cerebral overall, except for the last team, who appealed explicitly to emotion, which set the tone interestingly for the second half.

Players also began referring to each other as ``the other inner voices'', an organic renaming that feels more natural than ``Value Defenders''. I plan to adopt it.

\subsubsection{The Huddle}

Everyone found the Huddle very, very fun. The Subject player later said it was a genuine eureka moment when they realised they could ask particular Inner Voices to have very specific conversations with each other. It is worth noting that this player has meeting facilitation experience. During the Huddle, the Subject first asked each team to clarify things they were still confused about, then asked teams to expand on arguments they found particularly interesting, then asked pairs of Inner Voices to directly clash on issues they felt needed resolving.

At a certain point the Subject ran out of ideas for what to ask. I believe that with two players on the Subject team this would happen less frequently. Still, ten minutes felt too long; I would cut it to eight.

\subsubsection{Second Round of Inner Voices}

The Huddle exacerbated something that had already started in the first round: a slow drift toward consensus on what decision should be taken. The second round was much less about the practicalities of which decision is best, and much more about arguing which is the dominant lens through which the decision should be viewed.

My intuition is that something similar might happen with other groups, not only with veteran debaters.

During my own speech I also pivoted toward consensus on the right Solution. I did this because I found it interesting and novel, something regular debate never allows, and I wanted to see how it felt, to demonstrate that the rules permit it, and to explore whether I could deliberately fish for the ``Brave Pivot'' commendation. I found that I felt little need to rebut the points made by others; arguing for my lens and shifting focus felt like a more natural and persuasive strategy than attacking others, and I allowed myself to be openly persuaded by fellow Inner Voices on some points.

A fellow player later noted that in this format to outright attack what others say ``feels wrong'', disrespectful and unnecessary, and that the format requires tact in how one goes about diminishing the relevance of another's points.

\subsubsection{Second Subject Speech}

The Subject player structured their speech quite like a regular debate adjudication: they explained what they understood each Inner Voice team to have argued and how convincing each case was, then expressed the solution they had reached and what had factored most in their decision. They later said that during the first speech they were in more of a roleplay mindframe, but that as the debate progressed they transitioned more and more to feeling and thinking like a debate judge.

My intention had been for the second speech to remain somewhat more roleplay-ish, with the Subject speaking about how something \textit{felt} persuasive, being a bit less analytical. I am not sure whether adjusting the instructions is needed, or whether this was simply the style of this particular player.

\subsubsection{The Evaluator Speech}

The Evaluator gave insightful feedback on how each of us could have been more persuasive to a regular person. She then noted it was very difficult to find feedback to give, because the format does not specify clearly toward what educational objectives the feedback should be directed.

This raised a broader question I should think about more: traditional debate exists within a complex, organic constellation of clubs and competitions. Most people tell themselves they do debate to learn (critical thinking, public speaking), and trainers and judges tell themselves they are there to teach useful skills. In practice, my sense is that the circuit has a life of its own, and that people invest effort on either side mostly for the social reputation it yields.

How should this format be played? Who should facilitate it? Who is it for? What should it help participants achieve?

Another interesting suggestion was that a structure allowing the Evaluator to be more actively involved (to facilitate and give real-time feedback, to encourage players) would be useful. I tend to agree. Renaming the Evaluator ``the Facilitator'' might help signal this.

The Evaluator also said that awarding commendations was difficult, and that in a larger circuit it would become easier: she would not feel pressure to give everyone a commendation and would reserve them for truly outstanding moments. She added that once the commendations were memorised the system would feel much more organic than the numeric speech-point quantification that most traditional formats require of judges.

\subsubsection{On Bleed}

I decided to skip the final bleed-releasing ritual, expecting it would feel awkward. It was interesting, then, that even without it, one of the players spontaneously shared during a brief intermission that the debate had hit uncannily close to home: on the way to the session, a family member had asked them to testify for them in court, which they had found morally wrong.

\subsubsection{On Commendations}

There was substantial discussion about the commendation system. The more ``gamer'' of the players expressed that it scratches a ``make a build'' itch, but felt that commendations should carry some bonuses as well. They also raised the concern that players might start farming them, choosing to play in unusual ways simply to collect a specific commendation and ``complete'' a build.

\subsubsection{Participant Feedback}

Most players expressed that the game was quite fun and that they believed it would be enjoyable even outside a casual friends setting.

One noteworthy pattern in the feedback: women and players who had been less successful as traditional debaters said they strongly preferred this format. Men who had been highly successful debaters expressed finding it less exhilarating. One participant put it directly: ``There are few things more ego-boosting than the feeling of total mind-domination that you get when you win a debate, and this cannot give you that.''

One point everyone agreed on was that the format is much less anxiety-inducing than regular debate, while still providing a structure for meaningful conversation, unlike open real-life discussions, where people tend to fight over talking space.

One player suggested that the format might benefit from transitioning more wholeheartedly toward improv theatre, becoming a roleplaying game without any performance-evaluation component. Another countered that there is value in the sense that you are always progressing, at least a little, with every debate.

%%%%%%%%%%%%%%%%%%%%%%%%%%%%%%%%%%%%%%%%%

\newpage
\section*{Annex H -- Inside Out Debate Format: Third Design Iteration}
\addcontentsline{toc}{section}{Annex H -- Inside Out Debate Format: Third Design Iteration}

The third design iteration is based on observations and participant feedback gathered during the second playtest (April 2026). The changes below supersede the corresponding elements in the second iteration (Annex B).

\subsubsection{Role Renaming}

\textbf{Value Defenders are renamed Inner Voices.} This name emerged organically during the second playtest, when players began referring to each other as ``the other inner voices''. It better reflects the roleplay-oriented nature of the format and the relationship between the Value Defender teams and the Subject.

\subsubsection{S0 -- Opening Subject Speech Instructions}

The instructions for the opening Subject speech are revised to foreground the roleplay dimension:

\begin{quote}
Hold a monologue as if you are the character, alone in a place of comfort, speaking to yourself about the dilemma. Express why this decision is truly difficult and how much getting it right matters to you. Explain why each value is relevant to this decision, and how you cannot simply disregard any of them. Do not lean toward any decision yet.
\end{quote}

\subsubsection{The Huddle}

The Huddle is shortened from 10 minutes to \textbf{8 minutes}.

\subsubsection{S2 -- Final Subject Speech Instructions}

The instructions for the final Subject speech are revised to include emotional closure:

\begin{quote}
Synthesise all arguments and deliver a decision over a specific solution: a concrete course of action from those presented by the Inner Voices. Explain how everything in the debate contributed to this decision; leave nothing important out. Say why you feel this is the right decision, and how you now feel about making it.
\end{quote}

\subsubsection{Seating}

A seating instruction is added: try to arrange chairs so that the Inner Voices and the Subject face each other. Inner Voices speak from where they sit, standing or seated. The Subject does the same. The Evaluator sits somewhere outside the main circle, present but out of sight, out of mind.

\subsubsection{Specific Times}

The following timing guidelines are provided:

\begin{itemize}
    \item \textbf{Speaking times:} 5 minutes per speech.
    \item \textbf{Prep time:} 30 minutes for beginners or high-school debaters; 15 minutes for experienced adults.
    \item \textbf{Evaluator oral critique:} 10 minutes.
    \item \textbf{Evaluator-led debrief:} 15 minutes.
\end{itemize}

\subsubsection{Rituals}

The following rituals replace those described in the first design iteration (Annex A).

\paragraph{Rolling In.}
\begin{enumerate}
    \item \textbf{Evaluator opens.} The Evaluator says: \textit{``You are here to decide together.''}
    \item \textbf{All respond.} Everyone answers: \textit{``Yes, we are.''} followed by shared clapping.
\end{enumerate}

\paragraph{Calling an Inner Voice.}
\begin{enumerate}
    \item The Evaluator calls: \textit{``Voice, speak for [Value].''}
    \item The speaker \textbf{must begin their speech} with the words \textit{``I will\ldots''}
\end{enumerate}

\paragraph{Rolling Out.}
\begin{enumerate}
    \item \textbf{Voices' thanks.} Each Inner Voice says: \textit{``Thank you for listening to me.''}
    \item \textbf{Subject's thanks.} The Subject says to the Inner Voices: \textit{``Thank you for speaking your truth.''} The Subject then says to the Evaluator: \textit{``Thank you for listening to me.''}
    \item \textbf{Bleed check.} Everyone states their real name and how they feel right now. The Evaluator then leads a brief debrief.
\end{enumerate}

\subsubsection{Commendation System}

The commendation system is updated with a tier structure governing how commendations are distributed:

\begin{itemize}
    \item Every debater should receive at least one \textbf{T1 (Earned)} commendation per speech. At higher levels of play, T1 commendations carry less competitive weight but remain a baseline acknowledgement of participation.
    \item The Evaluator has a \textbf{limited pool of T2 (Strong) and T3 (Exceptional)} commendations to offer per debate. These are awarded in private.
    \item Some commendations may be \textbf{exclusive to specific Evaluators}, reflecting their expertise or perspective.
    \item Some rounds or competitions may designate a \textbf{featured commendation}: a particular commendation given special significance in that context.
\end{itemize}

\subsubsection{Evaluator Instructions}

The Evaluator role is reframed toward active facilitation:

\begin{itemize}
    \item You are the \textbf{facilitator} of the entire event. Your primary job is for everyone to feel safe and to play to grow.
    \item Feel free to \textbf{intervene} when it feels natural. Encourage players when you think it is necessary. Gently redirect when someone may be making others uncomfortable. Invite players out of their comfort zone when you believe it is in their best interest. Be human; establish a relationship of trust and growth with the others.
    \item After the debate ends, \textbf{first facilitate a brief debrief} stemming from the bleed check moment, then transition to the oral critique.
    \item Prepare \textbf{potential debrief questions} whilst others are preparing for the debate. Debrief questions should \textit{not} focus (as the oral critique does) on skill growth, but rather on \textbf{meaningful emotional growth and acceptance of others' realities}. Explore how debaters would act in a similar scenario and what their own values are regarding the values explored in the debate. When in doubt, use these two anchor questions: \textit{``How important are these values for you?''} and \textit{``Were you in a similar scenario, what other values would influence you?''}
    \item Open the oral critique with a \textbf{narrative of the debate} from your outside perspective: what each team did and how it all unfolded.
    \item Engage directly with each team: ask what they \textbf{felt they did well} and what they found difficult. Share your view and add your own observations.
    \item Ask each team what they want to \textbf{do better next time}, and give concrete advice on how to get there.
    \item Award at least one \textbf{T1 (Earned)} commendation per speech to every debater. Use your limited pool of \textbf{T2 (Strong) and T3 (Exceptional)} commendations sparingly; these are awarded in private.
\end{itemize}

\newpage

\section*{Annex I: Clues from Playtests}
\addcontentsline{toc}{section}{Annex I: Clues from Playtests}

The following catalogue is the output of the three-phase systematic analysis described in the Methods chapter: initial descriptive coding of field notes and debrief responses, comparative organization around continuity, difference, and novelty, and articulation of design principles connecting mechanics to observed effects.

Unlike a traditional results catalogue, which reports only what happened, this annex is organized around a pre-defined taxonomy of \textit{potential clues}: the observable signs that would count as evidence \textit{for} or \textit{against} the format's intended effects across each evaluation field. This taxonomy was specified \textit{before} data collection, as described in the Methods chapter, to ensure that analysis remained responsive to disconfirming evidence and did not converge only on confirming observations. For each potential clue, the catalogue records whether it was observed during the playtests, with what evidential support, and at what strength. Clues that did not manifest are recorded explicitly: the absence of a negative clue is evidence of the format functioning as intended; the absence of a positive clue is a genuine gap.

Evidence is classified by direction and strength. \textbf{Positive clues} are behaviors or reports consistent with the format achieving its intended effect. \textbf{Negative clues} are behaviors or reports indicating a design gap or failure. Strength is assessed as \textit{Strong} (behavioral, multi-source, or unsolicited), \textit{Moderate} (self-reported or single-source), or \textit{Weak} (speculative or highly context-dependent). Consistent with the comparative analysis described in the Methods chapter, behavioral and emergent evidence carries more weight than self-reported responses, and novel or different behaviors (departures from standard debate practice) constitute stronger design evidence than continuities, which may reflect participants' existing training rather than effects of the new format. Self-reported evidence gathered in response to direct questions in the semi-structured group debrief is noted where relevant, as it remains evidentially valid even if weaker than unsolicited or behavioral clues.

\textit{Sources: P1 = First playtest field notes (Jean Monet debate club, March 2026); P2 = Second playtest field notes (April 2026); D1/D2 = Post-session debrief responses after first/second playtest; T2 = Transcript of second playtest.}

\bigskip

\noindent\textbf{A note on social desirability in Field 1.} Engagement and popularity clues are the most susceptible to social desirability effects: participants knew they were playtesting a format designed by the researcher in the room, and the exclusive nature of the invitation may have elevated their investment in displaying enthusiasm. Positive engagement evidence should therefore be read with additional skepticism; behavioral indicators are weighted higher than self-reports, and negative clues, which participants would have had more reason to suppress, carry elevated weight when they do appear.

\bigskip

%%% FIELD 1 %%%

\noindent\textit{Why Field 1 is a required characteristic.} Transformative potential cannot manifest in a game that players do not return to. The research question asks how game design theory can elevate a format's transformative potential, but "elevated transformative potential" presupposes that participants inhabit the format long enough for its mechanisms to produce their effects. The transformational container forms and deepens through repeated participation; bleed arcs, skill accumulation, and epistemic shifts are all longitudinal processes that require the format to generate genuine voluntary re-engagement. A format that produces one interesting session but fails to motivate return has demonstrated novelty, not transformative design. Field 1 therefore functions as the necessary entry condition for all subsequent fields: every transformative effect assessed in Fields 2--5 is predicated on the format being one people choose to play again.

\bigskip

\subsection*{Field 1: Engagement, Appeal, and Long-Term Viability}

\subsubsection*{Positive potential clues}

\begin{enumerate}

\item \textbf{Unsolicited expressions of enthusiasm during or after the debate}\\
\textit{Points toward:} Positive affective engagement; potential for voluntary repeated participation.\\
\textit{Relevance to field:} Unsolicited expressions isolate intrinsic motivation as the causal driver, removing the social-obligation confound that makes prompted enthusiasm ambiguous. For Field 1, intrinsic motivation is the operative mechanism: a format people are enthusiastically drawn toward without external encouragement is one they have reason to choose on their own, which is the basis of the voluntary repeated participation that long-term viability requires.\\
\textit{Observed:} Yes (P1, P2, D1, D2).\\
\textit{Evidence:} Participants described the experience as ``surprisingly liberating'' (P1, D1) and ``very very fun'' during the Huddle phase (P2, D2). Multiple unprompted positive comments recorded in both debrief forms without being solicited.\\
\textit{Mechanic:} Removal of binary win/loss outcomes, combined with out-of-sight Evaluator seating, reduces performance-oriented anxiety; personal dilemma motions create genuine stakes without threatening social standing, making organic enthusiasm more likely.\\
\textit{Strength:} Moderate. Self-reported; subject to social desirability effects; see caveat above.

\item \textbf{Passionate or animated discussion among participants without researcher prompting}\\
\textit{Points toward:} Genuine peer engagement; format generates internal conversational momentum.\\
\textit{Relevance to field:} Discussion that sustains itself beyond the formal debrief, without researcher steering, demonstrates that the format's content generates its own conversational momentum independent of structure. For Field 1, the key indicator is that engagement outlasts the session: a format that continues to be discussed after the researcher has stopped directing the interaction has generated the pull that makes planning future sessions feel like an obvious next step.\\
\textit{Observed:} Yes (P1, P2).\\
\textit{Evidence:} Post-debrief conversation in both sessions extended organically beyond the structured questions. In the second playtest, participants continued discussing the motion and their own values informally after the session ended. Both sessions also produced animated group conversations specifically about the experience of using second-person direct address and personal appeals during speeches: participants in both groups noted that this mode of speaking (addressing the Subject directly as ``you'' and appealing to their personal circumstances) felt novel, engaging, and surprisingly enjoyable to both perform and receive.\\
\textit{Mechanic:} Personal dilemma motions make participants' own values directly relevant, creating natural grounds for continued conversation; the debrief structure designates and legitimizes extended reflection rather than truncating it. The second-person address mode inverts standard debate's impersonal register, making the emotional texture of the debate itself a topic of enthusiastic reflection.\\
\textit{Strength:} Moderate. Observed; behavioral; but also consistent with social novelty of the event.

\item \textbf{Visible engagement signals during debate itself (laughter, focused silence, forward leaning)}\\
\textit{Points toward:} Format sustaining in-the-moment attention and flow.\\
\textit{Relevance to field:} In-the-moment engagement signals confirm that the experience of playing, rather than only the memory of having played, is positively absorbing. For Field 1, this is the real-time counterpart to post-session enthusiasm: a format that participants remember fondly but are visibly disconnected from during play has a mismatch between anticipated and actual experience that would degrade over repetition.\\
\textit{Observed:} Yes (P1 and P2).\\
\textit{Evidence:} Second playtest field notes record observable engagement during speeches and very novel and creative decisions: to play background music during S1 speech, to adopt speaking accents etc. In the first playtest, as players became more aware of what they could do in the debate, they started to have more fun and more passionate speeches.\\
\textit{Mechanic:} The Huddle phase builds emotional investment before speeches begin through interactive preparation; the Subject as a live, named decision-maker provides a concrete human target rather than an abstract adjudicator, sustaining speaker investment.\\
\textit{Strength:} Moderate. Observational; first-playtest engagement was inconsistent.

\item \textbf{Unprompted recognition of the format's benefits}\\
\textit{Points toward:} Participants forming their own positive evaluation independent of researcher framing.\\
\textit{Relevance to field:} When participants independently articulate why the format is valuable, they are constructing the self-sustaining narrative that motivates voluntary future participation and recruits new players. For Field 1, long-term viability at community scale depends on participants having reasons they can articulate to others, not only on individual desire to return.\\
\textit{Observed:} Yes (P1, P2, D1, D2).\\
\textit{Evidence:} Multiple participants articulated benefits spontaneously: accessibility, empathy development, skill transfer, reduced anxiety. In the second playtest, there was conversation on how the format would be a useful personal development tool for ``corpo trainings.''\\
\textit{Mechanic:} Not tied to one particular mechanic but rather the format as a whole.\\
\textit{Strength:} Moderate. Self-report; theoretically coherent; partially behavioral (the personal-development-tool framing implies intention to act).

\item \textbf{Expressed desire to replay or requests to organize future rounds}\\
\textit{Points toward:} Replayability; format generates sustained rather than one-off interest.\\
\textit{Relevance to field:} Expressed desire to replay is the most proximate forward-looking indicator of long-term viability available within a single-session playtest. For Field 1, this clue also captures whether participants perceive an ongoing growth arc within the format: naming what they would do differently next time signals that the format has made its own depth visible, which is the motivational basis for sustained multi-session engagement.\\
\textit{Observed:} Yes (P1, D1).\\
\textit{Evidence:} Participants said they would play again and would aim to be ``more emotionally immersed'' next time (P1, D1). In the second playtest players reported they would find it a fun game especially to play with friends, as a passtime, and to get their "build" done.\\
\textit{Mechanic:} The roleplay dimension of the format (emotional inhabitation of a character's values) is recognized by participants as an unexplored skill, creating an intrinsic learning arc; the commendation system provides named targets (e.g., ``Believable Passion'') for improvement across sessions.\\
\textit{Strength:} Moderate. Self-reported intention; actual re-engagement behavior untested.

\item \textbf{High enjoyment scores when directly asked}\\
\textit{Points toward:} Positive affective engagement; format generating a genuinely enjoyable experience rather than merely tolerable or dutiful participation.\\
\textit{Relevance to field:} Enjoyment is the affective substrate of voluntary repeated participation: people return to activities they found genuinely pleasant, and a format that is intellectually interesting but unpleasant to play will be studied rather than played. For Field 1, consistently high scores across participants indicate the format reliably produces positive experience, not only occasional highlights.\\
\textit{Observed:} Yes (D1, D2).\\
\textit{Evidence:} When asked directly in the debrief form how much they enjoyed the game, participants in both sessions reported very high scores. While solicited self-report is weaker evidence than unsolicited enthusiasm, the consistently high numbers across both groups, including participants who were not particularly enthused about other aspects, indicate that the enjoyment response was not merely social desirability but genuine. The convergence with unsolicited expressions in clue 1 above strengthens the overall picture.\\
\textit{Mechanic:} Same as clue 1.\\
\textit{Strength:} Moderate. Direct debrief form response; subject to social desirability effects; convergent with unsolicited expressions.

\end{enumerate}

\subsubsection*{Negative potential clues}

\begin{enumerate}

\item \textbf{Taciturn, flat, or deflated behavior after the round}\\
\textit{Points away from:} Engagement and emotional investment; would suggest the format failed to generate genuine stakes.\\
\textit{Relevance to field:} Post-round affect is the most immediate behavioral indicator of whether the format generated genuine investment rather than dutiful participation. For Field 1, deflated post-round behavior would signal the experience left no residue worth processing, no sense of unfinished business or growth arc to continue, making voluntary return unlikely.\\
\textit{Observed:} No.\\
\textit{Evidence:} Not manifested across either playtest. Post-round behavior in both sessions was animated and conversational.\\
\textit{Strength:} N/A.

\item \textbf{Direct criticism of the format itself}\\
\textit{Points away from:} Overall appeal; structural design coherence.\\
\textit{Relevance to field:} Structural criticism identifies design barriers to adoption and re-engagement. For Field 1, usability problems are long-term viability problems: a format with persistent confusing or frustrating elements will lose players between sessions even if the core experience motivates return, because the overhead of relearning and re-onboarding discourages casual re-engagement.\\
\textit{Observed:} Partially (P1, D1).\\
\textit{Evidence:} The commendation system was criticised as ``overwhelming'' and difficult to use, particularly the size of the badge list (P1, D1). This was a targeted structural criticism rather than a wholesale rejection. The issue of commendation size was not directly addressed by a design change; it is to some extent intrinsic to the system. However, participants recognized that cognitive load would decrease with familiarity: as an Evaluator learns the commendation list over repeated use, the overhead diminishes naturally. Participants also noted that a practical workaround is available from the start: an Evaluator can simply observe what debaters do well in their speeches and work backwards to identify which commendation matches, rather than scanning the list prospectively. This makes the system more accessible at small scale than it first appears.\\
\textit{Mechanic:} The commendation list in the first design iteration (Annex A) was designed for tournament scale, where Evaluators memorize a limited set of commendations over repeated use; at small informal scale with an unfamiliar system, the cognitive overhead of tracking all options exceeded what the context could support. The issue is thus primarily one of onboarding rather than structural incoherence.\\
\textit{Strength:} Moderate. Specific and actionable; confined to one subsystem; not generalized to the format overall; partially mitigated by intrinsic learning curve and retrospective matching strategy.

\item \textbf{Visible loss of focus or disengagement during speeches}\\
\textit{Points away from:} Sustained attention; flow state during play.\\
\textit{Relevance to field:} In-session disengagement indicates the format is failing to sustain attentional pull in real time. For Field 1, a format that participants remember fondly but were visibly disconnected from while playing has an experience-memory gap that narrows over repetition: disengagement during play would eventually erode even enthusiastic retrospective impressions.\\
\textit{Observed:} Partially (P1).\\
\textit{Evidence:} One participant appeared uncertain about what to do during the Subject role in the first playtest, producing visibly hesitant behavior. This was role-confusion-driven rather than disengagement with the format broadly, and was not repeated in the second playtest.\\
\textit{Mechanic:} The S0 and S1 instructions in the first design iteration did not specify the Subject's active task during other teams' speeches, leaving the Subject without a clear role during the middle of the debate.\\
\textit{Strength:} Weak. Single instance; role-specific; addressed in next design iteration.

\item \textbf{Expressed preference for returning to standard competitive debate}\\
\textit{Points away from:} Standalone appeal; would suggest format is only tolerable as an alternative, not as a preferred experience.\\
\textit{Relevance to field:} Expressed preference for existing formats is the most direct comparative appeal test: it measures whether Inside Out is perceived as genuinely distinct and preferable, or as a tolerable but inferior alternative. For Field 1, long-term viability in a community that already has established format options depends on the format offering something those options do not.\\
\textit{Observed:} Partially (P1, D1, P2, D2).\\
\textit{Evidence:} Across both playtests, a total of three participants (one in the first session and two in the second) expressed reservations tied to the format's non-competitive structure. All three were men who had been highly competitive and quite successful in traditional debate. In the first session, one participant described the format as ``anticlimactic'' and expressed missing the ``ego-boosting'' feeling of decisive victory (P1, D1); several others in that session suggested the format might be ``easier'' or ``for beginners,'' raising concern it could be perceived as less serious by the broader debate community. In the second session, in addition to similar sentiments, one participant raised a more structural concern: that the format's apparent benefits might be partly an artifact of playing without an established meta: as a meta develops around Inside Out (as it has in World Schools and BP), the same competitive pressures and strategic distortions that the format tries to eliminate could re-emerge. This is a theoretically significant observation worth tracking in future iterations.\\
\textit{Mechanic:} The removal of binary win/loss outcomes and zero-sum scoring is a deliberate design choice; these negative responses confirm that mechanic is operating as designed, not malfunctioning. The format is not intended to serve this participant profile as a primary audience, though the meta-concern raised in D2 is a genuine design question about long-term robustness.\\
\textit{Strength:} Moderate for scoping the target audience; weak as evidence of a format-wide problem. The meta-concern is Weak at this stage (speculative, without empirical basis) but theoretically interesting.

\item \textbf{``Cringe'' reactions to ritual elements or format conventions}\\
\textit{Points away from:} Immersive design; willingness to engage with the format's fictional premise.\\
\textit{Relevance to field:} Ritual failure makes the format feel awkward or embarrassing to inhabit, and willingness to perform awkward conventions is higher among enthusiasts than among casual players. For Field 1, persistent ritual cringe would progressively narrow the format's accessible audience, capping the long-term viability it needs to develop its transformative potential at scale.\\
\textit{Observed:} Yes (P1, D1, P2, D2).\\
\textit{Evidence:} The ending ritual (Rolling Out) was described by participants as ``very robotic'' and as something that ``takes you out of the immersion'' (P1, D1). This reaction recurred in the second playtest: participants were observed mumbling the scripted phrases rather than delivering them with investment, and the researcher (as designer) was asked by participants to justify the purpose of the ritual phrases, an indication of a breakdown in the fictional frame (P2, D2). This is the strongest negative engagement clue from either playtest. It indicates that immersion \textit{was} present before that point (participants could be ``in the story'') but was disrupted by the ritual structure. The ritual was redesigned in the third design iteration (Annex H); the revised version has not yet been tested.\\
\textit{Mechanic:} The first-iteration Rolling Out ritual required all participants to recite identical scripted phrases (``Thank you for listening to me,'' ``Thank you for speaking your truth''), producing a collective recitation effect (uniform, simultaneous, formulaic) incompatible with the personal and emergent character of the rest of the format.\\
\textit{Strength:} Strong. Multi-session; multi-report; behavioral disruption (mumbling, questioning designer); design response drafted but untested.

\end{enumerate}

\bigskip

%%% FIELD 2 %%%

\noindent\textit{Why Field 2 is a required characteristic.} Growth through play requires genuine intellectual and social challenge. The research question asks for a format that elevates transformative potential, but transformation through structured activity depends on productive pressure: the kind of challenge that stretches participants' capacities and motivates them to develop. The Inside Out format deliberately removes the win/loss evaluative structure that standard debate uses as its primary pressure mechanism, addressing the community hierarchy problem the research question identifies. But this design choice creates a central tension: does removing the competitive stakes also remove the cognitive and social demands that make debate a growth activity? Without productive pressure, repeated play (confirmed by Field 1) would be pleasantly novel but developmentally inert. Field 2 tests whether the pressure was redirected rather than removed, asking whether a non-competitive evaluative structure can sustain the challenge that transformative engagement requires.

\bigskip

\subsection*{Field 2: Productive Pressure Mechanisms}

\subsubsection*{Positive potential clues}

\begin{enumerate}

\item \textbf{Visibly high engagement and effort during speeches}\\
\textit{Points toward:} Cognitive challenge retained; participants investing effort proportional to competitive formats.\\
\textit{Relevance to field:} Sustained high effort in the absence of win/loss stakes is the most direct test of Field 2's central question: whether productive pressure has been successfully redirected. If participants invest comparable effort to standard competitive debate without a competitive outcome driving that investment, the format has demonstrated that productive pressure does not require a binary result.\\
\textit{Observed:} Yes (P1, P2, T2).\\
\textit{Evidence:} In the second playtest, all Value Defenders immediately adopted elaborate roleplay stances (second-person address, emotional tone, worldbuilding detail) suggesting high rhetorical investment (P2, T2). The first playtest also showed high investment: participants came up with elaborate and nuanced solutions, constructed careful arguments, and several delivered quite passionate speeches, fully ``try-hard'' play despite the unfamiliar format (P1).\\
\textit{Mechanic:} The format as a whole. Interaction mechanics in particular (the Huddle and Points of Information) were among the most liked and requested features across both playtests, suggesting they play an especially important role in sustaining engagement and effort.\\
\textit{Strength:} Strong. Behavioral.

\item \textbf{Strong performance quality: complex arguments, empathetic reasoning, strategic engagement}\\
\textit{Points toward:} Format generating the right kind of cognitive demand.\\
\textit{Relevance to field:} Performance quality differentiates productive pressure from mere social pressure: it tests whether the format's challenge demands complex thinking and empathetic reasoning, not just effort for its own sake. For Field 2, quality matters because growth requires challenge of the right kind: participants can try very hard at a low-complexity task without developing.\\
\textit{Observed:} Yes (P1, P2, D1, D2).\\
\textit{Evidence:} Participants described ``architecting solutions'' as ``very complex and interesting strategically'' (P1, D1). In the second playtest, participants navigated value hierarchies and inter-value tensions without prompting. The Subject's S0 and S2 speeches in the second playtest demonstrated sophisticated contextual reasoning (T2). Both sessions also produced strikingly high levels of acting and roleplay: players adopted character mannerisms, delivered emotionally engaged speeches, and one played a background soundtrack during their S1. A key question raised in debrief was whether this emotional register came at the expense of critical thinking, but the evidence suggests it did not: very elaborate solutions, practicalities-based argumentation, and careful value hierarchies were consistently present, indicating that critical thinking and emotional engagement coexisted rather than traded off (P1, P2, T2).\\
\textit{Mechanic:} The motion structure provides a specific character context that makes empathetic reasoning contextually available and strategically useful, rewarding perspective-taking as a competitive skill. Notably, the source of performance pressure shifts in this format: rather than the win/loss dynamic, the pressure becomes social: being recognized and appreciated by fellow Inner Voices and the Subject for the quality and resonance of one's contributions. This social pressure is intrinsically tied to the quality of the engagement rather than to a binary outcome, making it a more nuanced and arguably healthier motivator.\\
\textit{Strength:} Moderate. Partially self-report; partially observed; consistent across sessions.

\item \textbf{Extensive post-debate reflection without researcher prompting}\\
\textit{Points toward:} Format generating cognitive and emotional residue worth processing.\\
\textit{Relevance to field:} Unsolicited post-debate reflection indicates the format generated cognitive and emotional residue that participants feel compelled to process, the phenomenological signature of productive pressure that goes beyond task completion. For Field 2, this is evidence that the challenge architecture produces the kind of generative engagement participants continue working through after the formal activity ends.\\
\textit{Observed:} Yes (P1, P2, D1, D2).\\
\textit{Evidence:} Both debrief sessions extended well beyond structured questions. Participants in the first playtest returned to discussing the motion's dilemma long after the formal debrief. Participants in both sessions also discussed, with evident pride, how fun and interesting it had been to try out techniques that were new to them: the second-person address, the emotional appeals, the worldbuilding. This is a distinct form of reflection from dilemma-processing: it is reflection on craft and creative agency, suggesting the format activates not just cognitive and moral residue but performative pride in artistic experimentation.\\
\textit{Mechanic:} The whole format; in particular, the open-ended character of the roleplay task provides room for genuine creative experimentation that standard debate's tight genre conventions do not.\\
\textit{Strength:} Moderate. Observed; informally behavioral; consistent across both playtests.

\item \textbf{Expressed drive to improve; participants identifying their own growth edges}\\
\textit{Points toward:} Format creating the kind of productive pressure that motivates development.\\
\textit{Relevance to field:} Self-identified growth edges demonstrate that the format's pressure is experienced as intrinsically motivating rather than externally coercive: participants want to develop because the format has made them aware of capacities they do not yet have. For Field 2, this is the self-sustaining pressure indicator: a challenge architecture that makes participants their own most demanding evaluators has achieved the goal better than any external scoring system could.\\
\textit{Observed:} Yes (P1, D1, P2, D2).\\
\textit{Evidence:} Participants in the first session said they would aim to be ``more emotionally immersed'' on replay (P1, D1), identifying immersion as a dimension they had not yet mastered. In the second playtest, participants initiated a group conversation about a specific skill the format demands but they found difficult: being polite when doing counter-argumentation, navigating how to diminish the relevance of another Inner Voice's points without it feeling dismissive or adversarial (P2, D2). This reflects a more specific and format-native growth edge than emotional immersion, and was entirely self-identified. However, the very informal friends-hangout setting of the second session made the learning conversation less sustained and dynamic than it might have been in a more structured context; the topic was raised but not fully explored.\\
\textit{Mechanic:} The commendation system names specific skills (e.g., ``Believable Passion'') that participants can identify as aspirational targets; the roleplay dimension of the format provides a specific, named axis for development on replay that is recognized as distinct from standard debate skills. The polite counter-argumentation skill is format-specific: standard debate does not reward it, so recognizing it as a growth edge requires playing Inside Out.\\
\textit{Strength:} Moderate. Self-report; forward-looking; actual improvement behavior untested. The P2 polite counter-argumentation conversation is particularly interesting as a format-native, self-generated growth goal.

\end{enumerate}

\subsubsection*{Negative potential clues}

\begin{enumerate}

\item \textbf{Apathy or minimal effort during speeches}\\
\textit{Points away from:} Format generating meaningful productive pressure.\\
\textit{Relevance to field:} Apathy in a non-competitive context is the primary risk that the removal of win/loss stakes creates, and would constitute direct evidence that the design failed to redirect pressure successfully. For Field 2, the meaningful absence of apathy across both sessions is therefore strong supporting evidence: the motivational vacuum that removing competition could have created did not materialize.\\
\textit{Observed:} No.\\
\textit{Evidence:} Not manifested across either playtest. All participants engaged substantively with their speeches.\\
\textit{Strength:} N/A.

\item \textbf{Improvement discussion only when directly prompted by evaluator or researcher}\\
\textit{Points away from:} Intrinsic motivation to grow; would suggest format is experienced as externally evaluated rather than self-developing.\\
\textit{Relevance to field:} Growth discussion requiring external prompting would indicate that the format's pressure mechanism is experienced as externally imposed rather than internally generated. For Field 2, this would reveal a critical gap between the format producing challenge and the format producing self-directed growth orientation, the deeper claim transformative design requires.\\
\textit{Observed:} Partially (D1, D2).\\
\textit{Evidence:} Some debrief forms contained brief or minimal self-improvement reflections that appeared to be prompted rather than unprompted. However, other responses in both sessions were substantive and self-initiated. The pattern is inconclusive.\\
\textit{Strength:} Weak. Inconclusive; mixed within both debrief sets.

\item \textbf{Confusion about what to do within role; confusion about how to improve}\\
\textit{Points away from:} Immediate accessibility; clear pressure architecture.\\
\textit{Relevance to field:} Role confusion indicates the pressure architecture is opaque or inaccessible: participants may be motivated to improve but cannot identify how or toward what. For Field 2, this is a design-gap indicator that tests whether the pressure mechanism is scaffolded clearly enough to translate effort into directed growth, or whether it generates undirected effort that dissipates without producing development.\\
\textit{Observed:} Yes (P1, D1, P2).\\
\textit{Evidence:} In the first playtest, the Subject player described not knowing ``what to say'' or how to ``play very clean,'' uncertain how to hold both emotional investment and evaluative neutrality simultaneously (P1, D1). In the second playtest, the Subject ran out of prepared Huddle questions before the 10 minutes elapsed (P2), revealing that the role requires more scaffolding than the ruleset provided. Both gaps were addressed in the second design iteration.\\
\textit{Mechanic:} The S0 speech instructions in the first iteration did not specify the Subject's emotional register or clarify the relationship between roleplay inhabitation (appropriate in S0/S1) and evaluative synthesis (appropriate in S2); the Huddle preparation time did not scaffold the kinds of questions the Subject should prepare. Both gaps were directly addressed in the second design iteration's revised Subject role instructions.\\
\textit{Strength:} Strong as a design-gap indicator. Behavioral; multi-instance; directly theoretically informative.

\end{enumerate}

\bigskip

%%% FIELD 3 %%%

\noindent\textit{Why Field 3 is a required characteristic.} The research question targets specific, identifiable problems in existing debate practice (value convergence, information density, and competitive hierarchy) and demands a format whose growth outputs address those problems directly. It is insufficient for a format to be engaging (Field 1) and challenging (Field 2) if the challenge it produces develops the same narrow skill set as traditional formats: faster argument construction, adversarial rhetorical technique, and strategic point-scoring. The transformative potential the research question seeks to elevate is potential for genuine personal development: improved empathy, contextual reasoning, collaborative discourse, and perspective-taking. Field 3 tests whether the format's procedural rhetoric successfully encodes and transmits these skills through its rule structure, without explicit instruction.

\bigskip

\subsection*{Field 3: Growth Direction, Useful Skills and Perspectives}

\subsubsection*{Positive potential clues}

\begin{enumerate}

\item \textbf{Novel play patterns transferring real-world skills (empathy, contextual reasoning, negotiation)}\\
\textit{Points toward:} Format's procedural rhetoric successfully encoding and transmitting its intended skill set.\\
\textit{Relevance to field:} Uninstructed adoption of perspective-taking behaviors is the primary test of procedural rhetoric's effectiveness: if participants behave in the ways the format's procedural logic rewards without being told to, the rules are successfully encoding the desired skill orientation at the behavioral level. For Field 3, behavioral evidence of skill transfer is stronger than any self-report because it shows the format producing the relevant behaviors during play rather than generating retrospective claims about development.\\
\textit{Observed:} Yes (P1, P2, D1, D2, T2).\\
\textit{Evidence:} In the second playtest, all Value Defenders spontaneously adopted second-person address, direct eye contact with the Subject, and contextualizing worldbuilding details, none of which were instructed in the ruleset (P2, T2). In the first playtest, participants articulated learning to ``think about the context of decisions'' and to be ``more human and more honest about human nature'' (P1, D1). One participant noted the format develops ``conversational skills and empathy without necessarily having this educational background'' (P2, D2). Three additional novel play patterns are worth noting: (1) a consistently \textit{polite and collaborative form of counter-argumentation} emerged organically: players shifted focus rather than attacked, built on others rather than dismantling them, and recognized that outright rebuttal felt ``wrong'' in this context; (2) participants in the second playtest spontaneously moved toward \textit{alignment on a common solution} across the second round of speeches, a behavior directly inverse to standard debate and a paradigmatic instance of the format's procedural rhetoric encoding collaborative negotiation; (3) a \textit{push toward synthesis} was evident in both rounds, with players building on rather than contradicting prior contributions, mirroring the kind of constructive dialogue the format is designed to develop.\\
\textit{Mechanic:} The character assignment system (``you represent the Subject's value of X'') and the personal dilemma motion structure (providing a named character with specific circumstances) invite perspective-taking without instructing it; the Subject as a named individual rather than an abstract adjudicator provides a concrete human recipient for empathetic reasoning. The absence of a rebuttal-rewarding scoring category inverts the rhetorical incentive structure, making polite, collaborative, synthesis-oriented engagement strategically advantageous rather than weakly punished.\\
\textit{Strength:} Strong. Behavioral for second playtest; self-report for first playtest; consistent direction.

\item \textbf{Continuation of standard debate's positive skill patterns (research, structured argumentation)}\\
\textit{Points toward:} Format building on rather than displacing existing debate skills.\\
\textit{Relevance to field:} The preservation of existing valuable skills demonstrates the format builds on rather than displaces the intellectual infrastructure debate already develops well. For Field 3, this matters because the research question asks for a format that addresses value convergence without sacrificing the intellectual challenge of the activity: a format that develops empathy at the cost of rigorous argument construction would trade one problem for another.\\
\textit{Observed:} Yes (P1, P2).\\
\textit{Evidence:} Participants continued to employ structured argument construction and research-based reasoning within speeches. The format did not produce a drift toward purely emotional or anecdotal argument; the value-argumentation task retained intellectual rigor.\\
\textit{Mechanic:} The speech order (A1, B1, C1, S1, A2, B2, C2, S2) preserves the argument-sequencing logic of standard debate; the value-argumentation task still rewards research and logical structure, redirected toward value priority rather than policy consequence chains.\\
\textit{Strength:} Low. Observational; consistent across both playtests, but it may be due more to the profile of the participants as experienced debaters. 

\item \textbf{Good play (empathetic, collaborative, value-aware) visibly rewarded in Subject Speech, Oral Critique, and debrief}\\
\textit{Points toward:} Evaluative culture aligned with transformative design intent.\\
\textit{Relevance to field:} The evaluative culture's recognition of empathetic and value-aware play demonstrates that the format's feedback loop is aligned with its growth direction: the skills being developed are also the skills being rewarded. For Field 3, this is the systemic confirmation: not just that participants develop useful skills, but that the format's community names and incentivizes those skills, which is what makes growth direction durable rather than session-specific.\\
\textit{Observed:} Partially (P1, P2, D1, D2).\\
\textit{Evidence:} Evaluators acknowledged empathetic and value-aware play in both oral critiques. Both playtest oral critiques were very insightful and offered meaningful feedback that rewarded good play. However, in both playtests Subject decisions did more rather reward overly-analytical play, but that points more towards the profile of the participants as highly experienced debaters.\\
\textit{Mechanic:} The oral critique provides a structured feedback channel explicitly tied to the evaluator's observations; the commendation system names specific skills to reward. Both mechanisms are directionally correct but depend on evaluator familiarity with the system, which is immature at small scale.\\
\textit{Strength:} Moderate. Partially implemented; evaluative culture immature at small scale.

\end{enumerate}

\subsubsection*{Negative potential clues}

\begin{enumerate}

\item \textbf{Non-standard behaviors being rewarded that do not transfer real-world skills}\\
\textit{Points away from:} Skill development value; would suggest format incentivizes play that is idiosyncratic to Inside Out rather than transferable.\\
\textit{Relevance to field:} Rewarding behaviors idiosyncratic to Inside Out would indicate the format has created its own insular skill set: developmentally valuable only within the format's fictional frame and not transferable to the real-world contexts the research question is concerned with. For Field 3, this would be the central failure mode: a format confirming Fields 1, 2, and 4 while developing skills that only matter within itself.\\
\textit{Observed:} No.\\
\textit{Evidence:} Not manifested across either playtest. No evaluator praise for behaviors inconsistent with empathy, contextual reasoning, or collaborative discourse was recorded.\\
\textit{Strength:} N/A.

\item \textbf{Emergence of unhelpful new play patterns (e.g., hollow performative emotionality without substance)}\\
\textit{Points away from:} Skill transfer; would suggest format's emotional register invites performance rather than genuine reflection.\\
\textit{Relevance to field:} Emotional performance without substance would indicate the format's roleplay invitation is being taken as license for affect theater rather than genuine perspective inhabitation. For Field 3, this would reveal a gap between the format's intended procedural rhetoric (emotional engagement as a route to perspective-taking) and its actual output: emotional performance as a substitute for it.\\
\textit{Observed:} No.\\
\textit{Evidence:} Not manifested across either playtest. Emotional engagement observed in both sessions was substantively connected to argument construction and scenario reasoning.\\
\textit{Strength:} N/A.

\item \textbf{Immersion-breaking format elements preventing the bleed arc from completing}\\
\textit{Points away from:} Memetic and emancipatory bleed; a truncated immersion arc limits transformative potential.\\
\textit{Relevance to field:} The bleed arc is the mechanism through which in-game learning becomes available as a resource for real-world situations; if the transformational container breaks before bleed completes, the skills developed inside the format remain quarantined within the fictional frame rather than crossing the play/life boundary. For Field 3, this is the most structurally significant gap: in-round growth direction can be confirmed while the bleed arc failure prevents that growth from becoming genuinely transferable.\\
\textit{Observed:} Yes (P1, D1).\\
\textit{Evidence:} The ending ritual (Rolling Out) was experienced as ``very robotic'' and described as breaking immersion at the moment it was intended to complete the emotional arc (P1, D1). Because emancipatory bleed requires a sustained transformational container through closing, a ritual that breaks character at the end undermines the mechanism. This is the most significant growth-direction design gap identified.\\
\textit{Mechanic:} The first-iteration Rolling Out ritual required participants to recite identical scripted phrases in sequence, producing a uniform collective-recitation effect incompatible with the personal and emergent character of preceding play. The ritual was redesigned in the third design iteration (Annex H); that revised version has not yet been tested.\\
\textit{Strength:} Strong. Multi-report; behavioral disruption; theoretical mechanism clearly implicated; design response drafted but untested.

\end{enumerate}

\bigskip

%%% FIELD 4 %%%

\noindent\textit{Why Field 4 is a required characteristic.} The transformational container, the social space in which growth occurs, must be sustained not only within a single session but across the community that repeatedly inhabits the format. The research question identifies "community hierarchy" as a specific structural problem in existing formats, and transformative game design theory holds that the container's integrity is a prerequisite for bleed and sustained growth: participants will not take the vulnerable risks that genuine transformation requires if the surrounding community rewards status competition, punishes unconventional play, or makes failure socially costly. Fields 1--3 confirm that individual participants engage, are challenged, and develop useful skills within a session; Field 4 tests whether the format's structural logic reshapes the social environment those sessions occur within: whether it can convert an adversarial community culture into one that sustains rather than undermines the container.

\bigskip

\subsection*{Field 4: Community Dynamics and Reduced Toxicity}

\subsubsection*{Positive potential clues}

\begin{enumerate}

\item \textbf{Unsolicited expressions of safety and reduced anxiety}\\
\textit{Points toward:} Transformational container functioning; psychological safety achieved without sacrificing engagement.\\
\textit{Relevance to field:} Anxiety reduction reported without direct prompting, across independent groups, demonstrates that the format's structural modifications are producing the psychological safety effect, not researcher reassurance or social expectation. For Field 4, the source of safety matters: safety created by researcher presence would not persist in the format's independent use; safety created by structural logic is durable and community-portable.\\
\textit{Observed:} Yes (P1, P2, D1, D2).\\
\textit{Evidence:} ``Much much less anxiety-inducing than standard debate'' was consistently reported across both playtests and multiple independent participants (P1, P2, D1, D2). Participants described the experience as ``quieter, more rested'' (P1) and noted a ``feeling of salvation regardless of the result'' (P2, D2). The consistency across two independent groups strengthens confidence in this finding.\\
\textit{Mechanic:} Four mechanics operate jointly to reduce performance anxiety: (1) the Evaluator seating arrangement places the observer out-of-sight and peripheral, removing the primary anxiety-generating figure of standard debate; (2) no public scoring occurs during the round; (3) binary win/loss outcomes are removed, eliminating zero-sum social capital stakes; (4) the personal dilemma motion has no objectively correct answer, removing the possibility of being demonstrably wrong.\\
\textit{Strength:} Very strong. Multiple independent reports; consistent across two sessions; theoretically predicted prerequisite for bleed.

\item \textbf{Brave, vulnerable play acknowledged and rewarded}\\
\textit{Points toward:} Evaluative culture recognizing and incentivizing the emotional risk the format asks of participants.\\
\textit{Relevance to field:} Community acknowledgment of brave play demonstrates that the format's evaluative culture is aligning with its design intent in real time: unconventional, emotionally risky contributions are received as strengths rather than exploitable weaknesses. For Field 4, this is the confirmation that the container is actively maintained by the community rather than merely assumed by the format's rules.\\
\textit{Observed:} Yes (P1, D1, P2, D2).\\
\textit{Evidence:} Multiple instances of brave and unconventional play were observed across both playtests. In the first session, the C team chose a deliberately vulnerable and conversational approach (emotionally intimate and unhurried) that is highly uncommon in the standard formats those participants were trained in; this was noted by others and acknowledged in the oral critique (P1, D1). In the second session, brave play took several forms: the S1 player chose to play a background soundtrack during their opening speech, a choice with no precedent in debate and a real social risk; the C team built their entire case around a worldbuilt character (the Subject's mother) and her emotional stakes; players across the room adopted highly expressive, roleplay-heavy delivery styles; and participants responded to these choices with visible gesture-based approval and enthusiastic nonverbal acknowledgment throughout: laughing, nodding, leaning forward. The unplanned bleed event (a participant sharing privately that the scenario hit close to a real family situation) was also noted and acknowledged in the oral critique (P2, D2). Formal commendation assignment for brave play was not consistently enacted due to system immaturity.\\
\textit{Mechanic:} The bleed-check moment in the third-iteration Rolling Out ritual explicitly creates a designated space for emotional acknowledgment; the Evaluator role instructions frame psychological safety as a primary responsibility, making responsive acknowledgment of vulnerable play structurally expected. The removal of win/loss outcomes reduces the cost of unconventional choices, making brave play a lower-risk option.\\
\textit{Strength:} Strong; multi-instance; behavioral; evaluative response was informal rather than systematic.

\item \textbf{Participants articulating liberation from standard debate's performance pressure (emancipatory bleed)}\\
\textit{Points toward:} Emancipatory bleed occurring: participants experiencing the format as freeing rather than constraining, and building confidence in their ability to speak up.\\
\textit{Relevance to field:} Emancipatory bleed reports, participants experiencing the format as freeing from the evaluative regime of existing formats, demonstrate that the format is dismantling rather than merely replacing standard debate's performance anxiety architecture. For Field 4, emancipatory bleed is the mechanism that specifically addresses the "community hierarchy" problem: it is the experiential evidence that the format is reaching and liberating the voices that status competition suppresses.\\
\textit{Observed:} Yes (P1, P2, D1, D2).\\
\textit{Evidence:} ``More relaxed, not performing execrably, felt like playing really okay'' (P1, D1); ``gives a little feeling of salvation, regardless of the result'' (P2, D2); ``everyone distinguished himself by his own style, which in ordinary debate competitions does not happen'' (P2, D2). These reflections articulate a liberation from the evaluative regime of standard debate, congruent with previously reported clues of an overall much lower level of anxiety experienced by all participants.\\
\textit{Mechanic:} Character assignment externalizes personal value expression into a role (``I am arguing for the value of X''), reducing the social exposure of holding a minority position; the removal of binary outcomes eliminates the competitive cost of expressing an unconventional view; the personal style encouraged by the roleplay premise means individual expressiveness is rewarded rather than penalized.\\
\textit{Strength:} Strong. Consistent across both playtests; multiple independently-worded formulations of the same attitudinal shift.

\item \textbf{Emergent norms of collaborative rather than adversarial discourse}\\
\textit{Points toward:} Format's procedural rhetoric successfully reshaping community conduct without explicit instruction.\\
\textit{Relevance to field:} Spontaneous norm articulation (participants identifying without instruction that adversarial rebuttal "feels wrong") is the highest-strength behavioral indicator that the format's procedural rhetoric has reshaped social orientation at the community level, not just the individual level. For Field 4, emergent norms are the gold standard of structural success: they show the format's logic has been collectively internalized, producing community behavior the ruleset did not explicitly mandate.\\
\textit{Observed:} Yes (P2, D2).\\
\textit{Evidence:} Participants in the second playtest spontaneously articulated that ``outright attacking others' points feels wrong, disrespectful, and unnecessary.'' This was not a rule in the format but an emergent norm arising from the format's structure, a paradigmatic instance of procedural rhetoric in action. Players also noted that shifting focus and building on others' contributions felt ``more natural and more persuasive'' than traditional rebuttal.\\
\textit{Mechanic:} The Subject as shared persuasion target means all Value Defenders are competing for the same person's attention rather than defeating each other; no scoring category rewards rebuttal specifically, meaning attacking a competitor's argument carries no direct competitive benefit while building on it does, thus inverting the rhetorical incentive structure of standard competitive debate.\\
\textit{Strength:} Strong. Behavioral; emergent; unsolicited; theoretically predicted.

\end{enumerate}

\subsubsection*{Negative potential clues}

\begin{enumerate}

\item \textbf{Expressions of distress during or after the round}\\
\textit{Points away from:} Psychological safety; would suggest the format activates personal stakes without providing sufficient containment.\\
\textit{Relevance to field:} Distress in a format that deliberately elicits personal stakes would indicate the format lacks sufficient containment for the vulnerability it invites: the transformational container is failing its primary function. For Field 4, this is the most consequential failure mode: a format that elicits vulnerability and fails to hold it safely is not merely ineffective but actively harmful, which is why the absence of distress despite moments of genuine emotional intensity is particularly significant.\\
\textit{Observed:} No.\\
\textit{Evidence:} Not manifested across either playtest. The unplanned bleed event (participant sharing a real family situation) was handled calmly within the debrief structure; no distress was recorded. The bleed-check and debrief protocols appear to have functioned as intended.\\
\textit{Strength:} N/A. Notable absence given the emotional intensity of some moments.

\item \textbf{Arrogant, dismissive, or unempathetic behavior during debate}\\
\textit{Points away from:} Community health; toxicity reduction.\\
\textit{Relevance to field:} Adversarial social behavior from participants trained in competitive debate would indicate that the format's structural modifications have not successfully overridden the behavioral dispositions those participants bring from standard formats. For Field 4, this would reveal a gap between structural design and cultural practice: the format's rules may reward collaborative behavior while the community's habits continue to reward adversarial conduct.\\
\textit{Observed:} No.\\
\textit{Evidence:} Not manifested across either playtest. No dismissive or contemptuous behavior toward other participants was recorded in either field notes or transcripts.\\
\textit{Strength:} N/A.

\item \textbf{Such behavior being praised by the evaluator or community}\\
\textit{Points away from:} Evaluative culture aligned with design intent; would signal the format's recognition system is rewarding the wrong behaviors.\\
\textit{Relevance to field:} Praise of dismissive behavior would indicate the format's evaluative culture has been captured by existing norms rather than reshaping them: the recognition system, the primary lever for community change, is pointing in the wrong direction. For Field 4, the evaluative culture test is the most systemic indicator: if the community rewards the behaviors the format is designed to discourage, the social logic will prevail over the structural logic over time.\\
\textit{Observed:} No.\\
\textit{Evidence:} Not manifested across either playtest.\\
\textit{Strength:} N/A.

\end{enumerate}

\bigskip

%%% FIELD 5 %%%

\noindent\textit{Why Field 5 is a required characteristic.} The research question centers on "value convergence" as the primary epistemic problem the format must address, and epistemic humility (the capacity to hold beliefs as provisional, to recognize the legitimacy of competing value frameworks, and to treat disagreement as inquiry rather than contest) is its direct structural opposite. Fields 1--4 together demonstrate that the format is engaging, challenging, skill-building, and community-safe; but none of those fields tests whether the format addresses value convergence at the level of how participants hold their own beliefs. A format could confirm all four earlier fields while still producing participants who are enthusiastic, pressured, skilled, and socially safe but epistemically unchanged, still convinced that their values are objectively correct and others' are deficient. Field 5 tests the deepest targeted transformation: whether the format's procedural structure (Inner Voices framing, S0 bias disclosure, the absence of a winner, the multi-perspective structure) produces a measurable shift toward value pluralism as an emergent consequence of inhabiting the format's logic.

\bigskip

\subsection*{Field 5: Epistemic Humility}

\subsubsection*{Positive potential clues}

\begin{enumerate}

\item \textbf{Unprompted awareness that multiple competing perspectives hold genuine value}\\
\textit{Points toward:} Core attitudinal shift targeted by the format: epistemic humility arising from procedural engagement with value conflict.\\
\textit{Relevance to field:} Attitudinal shifts toward value plurality arising without researcher prompting demonstrate that the format's procedural structure produced the epistemic orientation the research question targets, rather than a researcher-advocated conclusion. For Field 5, the unprompted character of this evidence is essential: the research question asks whether the format can address value convergence structurally; if the shift requires prompting, it confirms only that the research team valued epistemic humility, not that the format produces it.\\
\textit{Observed:} Yes (P1, D1, P2, D2).\\
\textit{Evidence:} ``Every team says valuable stuff and it is important more to see what is important to us'' (P1, D1). This formulation, articulated without prompting after a single encounter, captures precisely the attitudinal shift the design targets: moving from ``which side is right'' to ``what do I value and why.'' Additional supporting evidence from the first playtest includes participants articulating that arguing someone is ``morally wrong'' for their personal decisions feels ``a bit absurd'' in this context, and recognition that the format makes personal value priorities legible rather than evaluating them against an external standard (P1, D1). In the second playtest, a notable extended conversation emerged in the debrief about how the Subject's opening speech (S0) institutionalises the practice of surfacing personal biases, making explicit what a player personally cares about before the debate even begins. Participants contrasted this with other formats, where the pretense of objectivity is maintained throughout: in those formats, they argued, biases are falsely supposed not to exist, propagating the idea that only one particular type of ``objective'' person can judge well and legitimizing social pressure to perform that identity. The S0 mechanism, by contrast, invites players to name their values openly from the start, reframing bias not as contamination but as perspective (P2, D2).\\
\textit{Mechanic:} The three-value motion structure assigns each team a value that is genuinely defensible for the character; the Subject's final synthesis requirement frames the question as ``what matters most to this person in this context'' rather than ``which position is objectively correct,'' making all three perspectives legitimate from the outset. The S0 speech institutionalises bias transparency rather than suppressing it, which makes epistemic humility structurally available rather than personally effortful.\\
\textit{Strength:} Moderate. Multi-session; self-report; directly on-target with design intent; P1 formulation notable for being unprompted; P2 debrief conversation particularly theoretically rich.

\item \textbf{Demonstrated understanding that personal values are context-dependent, not universal}\\
\textit{Points toward:} Internalization of the format's procedural rhetoric: that decisions are contextually embedded and legitimately vary by person and circumstance.\\
\textit{Relevance to field:} Recognizing that decisions are personal rather than universal is the specific intellectual content that distinguishes epistemic humility from mere tolerance: it is not just accepting that others hold different views but understanding that those views are legitimately contextual. For Field 5, this is the cognitive-structural indicator: the format has made contextual embeddedness of values experientially legible through play, not through assertion.\\
\textit{Observed:} Yes (P1, P2, D1, T2).\\
\textit{Evidence:} ``Sometimes decisions are personal, not universal'' (P1, D1). The Subject's S0 and S2 arc in the second playtest demonstrated a similar trajectory: moving from roleplay inhabitation (personal, contextual) toward evaluative synthesis (what values do I hold about this?), with explicit self-awareness of that transition (P2, T2).\\
\textit{Mechanic:} The character-specific motion context (personal history, relationships, circumstances) makes universal moral claims strategically inappropriate to the scenario; the S0 instructions ask the Subject to describe the scenario from a personal perspective before argumentation begins, embedding contextual reasoning before any abstract principle is introduced.\\
\textit{Strength:} Moderate. Multi-session; self-report and behavioral; theoretical coherence strong.

\item \textbf{Spontaneous movement toward synthesis rather than competitive winning}\\
\textit{Points toward:} Format orienting participants toward value integration rather than dominance: the deepest available behavioral indicator of epistemic humility in action.\\
\textit{Relevance to field:} Synthesis behavior emerging from the format's structural logic, without instruction and in inversion of standard debate's escalation pattern, is the deepest behavioral indicator that the epistemic shift has reached the level of action, not only self-perception. For Field 5, the behavioral character of this evidence is decisive: self-reports of attitudinal shift are susceptible to social desirability effects; transcript-confirmed spontaneous synthesis demonstrates the format has made integration rather than dominance the natural cognitive direction.\\
\textit{Observed:} Yes (P2, T2).\\
\textit{Evidence:} In the second round and also less in the first one of Value Defender speeches, participants naturally moved toward synthesis rather than intensifying conflict, the reverse of what standard debate produces. This was not instructed; it emerged from the format's procedural logic (P2, T2). Additionally, the Subject's opening statement ``I'm not sure if that's the solution; that's why this thing is eating at me'' (P2, T2, transcript-confirmed) demonstrates genuine epistemic uncertainty from the start of the round, produced by the format's demand that the Subject hold the dilemma rather than resolve it prematurely. Two further details reinforce this orientation toward collegiality and synthesis: (1) participants in the second playtest organically renamed themselves ``Inner Voices'' (which was subsequently formalized in the third design iteration), a label that positions them as parts of a shared internal deliberation rather than as opposing sides, colleagues rather than opponents; (2) the rules themselves permit multiple Inner Voices to argue for the same solution, structurally allowing consensus rather than requiring disagreement, which removes the adversarial constraint that standard debate formats impose.\\
\textit{Mechanic:} Speeches A2/B2/C2 are expected to build on A1/B1/C1 rather than dismantle competitors' positions. The Subject's S1 is an explicitly preliminary decision requiring ongoing engagement, structurally rewarding continued dialogue. The Huddle allows the Subject to ask clarifying questions, modeling collaborative inquiry as a legitimate competitive move. The Inner Voices naming and the same-solution permission together reframe competition from adversarial to collegial at the level of both identity and rules.\\
\textit{Strength:} Strong. Behavioral; transcript-confirmed; emergent; theoretically predicted.

\end{enumerate}

\subsubsection*{Negative potential clues}

\begin{enumerate}

\item \textbf{Unequivocal condemnation of beliefs or belief-holders}\\
\textit{Points away from:} Epistemic humility; would indicate the format is activating moral certainty rather than productive uncertainty.\\
\textit{Relevance to field:} Condemnation of positions or belief-holders would demonstrate that the format activates moral certainty rather than productive uncertainty: instead of making epistemic humility the natural cognitive orientation, the format's emotional intensity is channeling into absolutism. For Field 5, this is the mirror-image failure mode: a format that elicits strong value engagement could either produce humility (the target) or produce entrenchment (its opposite), and this clue tests which outcome the format's structural logic generates.\\
\textit{Observed:} No.\\
\textit{Evidence:} Not manifested across either playtest or in the transcripts. No unequivocal condemnation of a position, values system, or person was recorded.\\
\textit{Strength:} N/A. Absence is notable given that the motions used both involved choices with genuine ethical weight.

\item \textbf{Such condemnation being normalized or praised}\\
\textit{Points away from:} Community epistemic culture; would signal the format's evaluative frame rewards dogmatic certainty.\\
\textit{Relevance to field:} Normalization of condemnation would indicate the epistemic failure has reached the community level: not just that individuals experience certainty but that the evaluative culture rewards it. For Field 5, the evaluative culture test is the systemic indicator: the research question's goal of a format that "addresses the issues of value convergence" requires the format's community to sustain epistemic humility collectively, not only momentarily within individuals.\\
\textit{Observed:} No.\\
\textit{Evidence:} Not manifested across either playtest.\\
\textit{Strength:} N/A.

\end{enumerate}



%%%%%%%%%%%%%%%%%%%%%%%%%%%%%%%%%%%%%%%%%
% REFERENCES
%%%%%%%%%%%%%%%%%%%%%%%%%%%%%%%%%%%%%%%%%

\newpage
\printbibliography[title={REFERENCES}]

\end{document}
